\documentclass[10pt,a4paper]{article}
% ---------------------------------------------------------------------------
% Teaching edition: the agentm2m Claude Code plugin, its MCP server and the
% hosted service, explained for undergraduate software engineering students.
% Compiles with pdflatex (+ bibtex) or with tectonic (XeTeX).
% Every "Observed" output can be reproduced with docs/plugin-walkthroughs/.
% ---------------------------------------------------------------------------

\usepackage{iftex}
\usepackage[T1]{fontenc}
\ifPDFTeX\usepackage[utf8]{inputenc}\fi
% --- typography: Palatino text+math, Helvetica headings, Bera Mono code ---
\usepackage[sc]{mathpazo}
\linespread{1.06}
\usepackage[scaled=0.9]{helvet}
\usepackage[scaled=0.84]{beramono} % has a real bold face (txtt has none; faked bold crashed xdvipdfmx)
\ifPDFTeX
  \usepackage[protrusion=true,expansion=true]{microtype}
\else
  \usepackage[protrusion=true,expansion=false]{microtype}
\fi
\usepackage[margin=2.3cm,top=2.6cm,bottom=2.6cm,headsep=0.8cm]{geometry}
\usepackage{amsmath,amssymb,amsthm}
\usepackage{booktabs,multirow,array,tabularx}
\usepackage[table]{xcolor}
\usepackage{enumitem}
\usepackage{listings}
\usepackage{tikz}
\usetikzlibrary{arrows.meta,positioning,shapes.geometric,fit,backgrounds,calc}
\usepackage{pifont}
\usepackage[strings,nohyphen]{underscore} % plain _ in \texttt identifiers, no hyphen at breaks
\usepackage{titlesec}
\usepackage{needspace}
\usepackage{etoolbox}
\usepackage{fancyhdr}
\usepackage[section]{placeins} % floats never drift into the next section
\usepackage[font=small,labelfont={bf,sf},labelsep=period,skip=5pt]{caption}
\usepackage[skins,breakable]{tcolorbox}
\usepackage{hyperref}
\usepackage[capitalise,noabbrev]{cleveref}

% --- palette (same as the AutoM2M teaching edition) ---
\definecolor{accent}{RGB}{28,62,112}
\definecolor{accentlight}{RGB}{232,238,247}
\definecolor{rawbg}{RGB}{255,248,230}
\definecolor{rawbar}{RGB}{214,150,40}
\definecolor{modelbg}{RGB}{235,244,255}
\definecolor{modelbar}{RGB}{60,110,180}
\definecolor{enginebg}{RGB}{236,250,238}
\definecolor{enginebar}{RGB}{60,145,80}
\definecolor{llmbg}{RGB}{248,238,252}
\definecolor{llmbar}{RGB}{140,80,160}
\definecolor{wirebg}{RGB}{242,242,242}
\definecolor{wirebar}{RGB}{110,110,110}
\definecolor{extbg}{RGB}{255,255,255}
\definecolor{extbar}{RGB}{120,120,120}
\definecolor{badbg}{RGB}{252,236,236}
\definecolor{badbar}{RGB}{170,50,50}

\hypersetup{colorlinks=true, linkcolor=accent, citecolor=enginebar,
  urlcolor=accent, pdftitle={Plugging the Team In: the AgentM2M Plugin, its MCP Server and the Hosted Service, Explained},
  pdfauthor={}}

% --- headings ---
\titleformat{\section}{\sffamily\Large\bfseries\color{accent}}
  {\thesection}{0.7em}{}[{\color{accent!35}\titlerule[0.6pt]}]
\titleformat{\subsection}{\sffamily\large\bfseries\color{accent}}{\thesubsection}{0.6em}{}
\titleformat{\paragraph}[runin]{\sffamily\bfseries\color{accent}}{}{0pt}{}[.]
\pretocmd{\section}{\needspace{10\baselineskip}}{}{}
\pretocmd{\subsection}{\needspace{7\baselineskip}}{}{}
\titlespacing*{\section}{0pt}{3.2ex plus 1ex minus .2ex}{1.6ex plus .2ex}
\titlespacing*{\subsection}{0pt}{2.6ex plus 1ex minus .2ex}{1.0ex plus .2ex}
\titlespacing*{\paragraph}{0pt}{1.6ex plus .5ex minus .2ex}{0.8em}

% --- header / footer ---
\pagestyle{fancy}
\fancyhf{}
\renewcommand{\sectionmark}[1]{\markboth{\thesection\quad #1}{}}
\fancyhead[L]{\sffamily\footnotesize\color{black!55}Plugging the Team In --- Teaching Edition}
\fancyhead[R]{\sffamily\footnotesize\color{black!55}\nouppercase{\leftmark}}
\fancyfoot[C]{\sffamily\footnotesize\color{black!60}\thepage}
\renewcommand{\headrulewidth}{0.4pt}
\renewcommand{\headrule}{{\color{black!20}\hrule width\headwidth height\headrulewidth}}
\fancypagestyle{plain}{\fancyhf{}\renewcommand{\headrulewidth}{0pt}%
  \fancyfoot[C]{\sffamily\footnotesize\color{black!60}\thepage}}

\setlist{itemsep=2pt, topsep=4pt}
\setlength{\parindent}{0pt}
\setlength{\parskip}{5pt plus 1pt}
\setlength{\emergencystretch}{3em}
\widowpenalty=10000 \clubpenalty=10000
\renewcommand{\arraystretch}{1.12}
\setlength{\tabcolsep}{5pt}
\newcolumntype{Y}{>{\raggedright\arraybackslash}X}
\newcolumntype{L}[1]{>{\raggedright\arraybackslash}p{#1}}

\newcommand{\cmark}{{\color{enginebar}\ding{51}}}
\newcommand{\xmark}{{\color{red!65!black}\ding{55}}}
\newcommand{\pmark}{\ding{109}}

% --- shorthands ---
\newcommand{\tool}{\textsc{AgentM2M}}
\newcommand{\auto}{\textsc{AutoM2M}}
\newcommand{\cc}{Claude Code}
\newcommand{\W}[1]{\textsf{W#1}}
\newcommand{\D}[1]{\textsf{D#1}}
\newcommand{\devteam}{\textsc{DevTeam}}
\newcommand{\cls}[1]{\textsf{#1}}
\newcommand{\code}[1]{\texttt{#1}}
\newcommand{\file}[1]{\texttt{#1}}
\newcommand{\mcptool}[1]{\texttt{#1}}
\newcommand{\skill}[1]{\texttt{/agentm2m:#1}}

% --- listing styles ---
\lstdefinelanguage{ATL}{
  morekeywords={module,create,from,rule,to,lazy,unique,helper,def,
                context,uses,thisModule,and,or,not,in},
  morekeywords=[2]{@llm,@check}, alsoletter={@},
  sensitive=true, morecomment=[l]{--}, morestring=[b]'}

% No ligature-like replacement of -> or <- anywhere: students copy listings.
\lstset{frame=l, framerule=2.2pt, framesep=7pt, xleftmargin=10pt,
        framexleftmargin=0pt,
        basicstyle=\ttfamily\footnotesize, breaklines=true,
        breakatwhitespace=false, columns=fullflexible, keepspaces=true,
        showstringspaces=false, numbers=none, upquote=true,
        keywordstyle=\bfseries, keywordstyle={[2]\bfseries\color{purple}},
        commentstyle=\itshape\color{black!55},
        captionpos=t, aboveskip=8pt, belowskip=8pt}
% "raw" = a file or text a person authors: a task, a command, a configuration file, a skill
\lstdefinestyle{raw}{backgroundcolor=\color{rawbg}, rulecolor=\color{rawbar}}
% "model" = source code: Python, rule modules, typed teams
\lstdefinestyle{model}{backgroundcolor=\color{modelbg}, rulecolor=\color{modelbar}}
% "engine" = deterministic output of the engine, checker or server
\lstdefinestyle{engine}{backgroundcolor=\color{enginebg}, rulecolor=\color{enginebar}}
% "llm" = a prompt sent to / an answer received from the LLM (here: Claude Code)
\lstdefinestyle{llm}{backgroundcolor=\color{llmbg}, rulecolor=\color{llmbar}}
% "wire" = a protocol message between two programs (line breaks added)
\lstdefinestyle{wire}{backgroundcolor=\color{wirebg}, rulecolor=\color{wirebar}}
\captionsetup[lstlisting]{singlelinecheck=false, margin=10pt}

% keep every listing on one page (all are short)
\BeforeBeginEnvironment{lstlisting}{\par\vspace{2pt}\noindent\begin{minipage}{\linewidth}}
\AfterEndEnvironment{lstlisting}{\end{minipage}\par\vspace{2pt}}
\crefname{lstlisting}{Listing}{Listings}
\Crefname{lstlisting}{Listing}{Listings}

% --- remarks as soft boxes ---
\newcounter{remark}
\newenvironment{remark}[1][]{%
  \refstepcounter{remark}%
  \begin{tcolorbox}[enhanced, breakable, colback=accentlight, colframe=accentlight,
      borderline west={2.2pt}{0pt}{accent}, arc=0pt, boxrule=0pt,
      left=8pt, right=8pt, top=5pt, bottom=5pt, before skip=8pt, after skip=8pt]
  \textsf{\bfseries\color{accent}Remark~\theremark\if\relax\detokenize{#1}\relax\else\ (#1)\fi.}\enspace}%
  {\end{tcolorbox}}

\newcommand{\what}[1]{\par\smallskip\noindent{\sffamily\bfseries\color{accent}#1}\quad}
\newcommand{\legend}[3]{\tikz[baseline=(c.base)]{\node(c)[fill=#1,inner xsep=6pt,inner ysep=2.5pt,
  font=\sffamily\footnotesize\bfseries]{\strut #3};\fill[#2](c.north west)rectangle([xshift=2.2pt]c.south west);}}

% --- teaching-edition boxes ---
\definecolor{thinkbg}{RGB}{255,246,236}
\definecolor{thinkbar}{RGB}{200,110,30}
\definecolor{keybg}{RGB}{236,247,240}
\definecolor{keybar}{RGB}{40,130,80}
\newcounter{think}
\newenvironment{think}[1][Check your understanding]{%
  \refstepcounter{think}%
  \begin{tcolorbox}[enhanced, breakable, colback=thinkbg, colframe=thinkbg, arc=0pt, boxrule=0pt,
     borderline west={2.2pt}{0pt}{thinkbar}, left=8pt, right=8pt, top=5pt, bottom=5pt,
     before skip=8pt, after skip=8pt]
  \textsf{\bfseries\color{thinkbar}Q\thethink.\ #1.}\enspace}{\end{tcolorbox}}
\newenvironment{takeaway}{%
  \begin{tcolorbox}[enhanced, breakable, colback=keybg, colframe=keybg, arc=0pt, boxrule=0pt,
     borderline west={2.2pt}{0pt}{keybar}, left=8pt, right=8pt, top=5pt, bottom=5pt,
     before skip=8pt, after skip=8pt]
  \textsf{\bfseries\color{keybar}Key takeaways.}\par\smallskip}{\end{tcolorbox}}
\newcommand{\plain}[1]{\par\smallskip\noindent{\sffamily\bfseries\color{thinkbar}In plain words:}\enspace\emph{#1}\par\smallskip}

\begin{document}

% ---------------------------------------------------------------------------
% Title block
% ---------------------------------------------------------------------------
\thispagestyle{plain}
\begin{center}
{\sffamily\footnotesize\bfseries\color{accent}\MakeUppercase{Teaching edition}}\par
\vspace{6pt}
{\sffamily\bfseries\fontsize{22}{26}\selectfont\color{accent}Plugging the Team In\par}
\vspace{6pt}
{\sffamily\large The \tool{} Plugin for \cc{}, its MCP Server and the Hosted Service,\\
Explained for Software Engineering Students\par}
\vspace{10pt}
{\small\itshape Companion to ``When the Team Writes Itself'' (the \auto{} teaching edition)~\cite{autom2mteaching}\\
and to ``\tool{}: M2M Transformations for Agentic Collaborations in Software Engineering''~\cite{agentm2m2026}\par}
\vspace{8pt}
{\color{accent!40}\rule{0.35\textwidth}{0.6pt}}
\end{center}
\vspace{4pt}

\begin{tcolorbox}[enhanced, colback=accentlight, colframe=accentlight, arc=0pt,
   boxrule=0pt, borderline west={2.2pt}{0pt}{accent}, left=10pt, right=10pt,
   top=7pt, bottom=7pt]
{\sffamily\bfseries\color{accent}Abstract.}\enspace
\tool{} runs a team of AI agents whose hand-offs are checked by a program:
a deterministic engine decides what each agent must produce, and a large
language model (LLM) only fills in individual values, each of which a
validator must accept. \auto{} lets an LLM design the whole team and admits
the design only if a checker approves it. This document explains how both
were turned from a research library into tools a developer uses every day:
a \emph{plugin} for the coding assistant \cc{}, a server that speaks the
\emph{Model Context Protocol} (MCP), and a \emph{hosted service} that offers
the same tools over the web. Its central idea is \emph{host mode}: \cc{} does
all the LLM work, but only by answering prompts the engine builds, and the
engine keeps every decision. We build all background from what a
software-engineering undergraduate knows, follow two complete sessions whose
outputs come from the real engine, and close with security, testing, and an
honest account of what is not solved.
\end{tcolorbox}

{\hypersetup{linkcolor=black}
\setcounter{tocdepth}{1}
\makeatletter
\renewcommand*\l@section[2]{\addpenalty{-\@highpenalty}\vskip 0.15em
  {\sffamily\small\@dottedtocline{1}{0em}{2em}{#1}{#2}}}
\renewcommand*\l@subsection[2]{\small\@dottedtocline{2}{2em}{2.6em}{#1}{#2}}
\makeatother
\renewcommand{\contentsname}{\sffamily\color{accent}Contents}
\tableofcontents}

% ===========================================================================
\section{How to read this document}\label{sec:read}
% ===========================================================================

\paragraph{Who this is for} An undergraduate who has taken a first software
engineering course: you know classes and interfaces, unit tests, JSON, and
roughly what a web API is. You do not need to know large language models,
the Model Context Protocol, \cc{}, or model-driven engineering;
\Cref{sec:course} teaches what you need. The companion \auto{} teaching
edition~\cite{autom2mteaching} explains \emph{why} agent teams need checked
interfaces; this document explains \emph{how} that idea reaches a
developer's editor. You can read it without the companion.

\paragraph{What you will be able to do} After reading this document you can:
\begin{enumerate}[itemsep=1pt]
\item explain, message by message, how an AI assistant discovers and calls
the tools of an MCP server;
\item name the four extension points of a \cc{} plugin and say what the
\tool{} plugin puts in each;
\item explain why \emph{host mode} lets \cc{} do all the LLM work without
weakening any guarantee of the engine, and where that argument has limits;
\item read a hand-off rule, predict which values a requirement change
obliges to be redone, and check your prediction with \mcptool{impact};
\item explain how the hosted service keeps different users' data and code
apart, and why it refuses to run code when it cannot isolate it.
\end{enumerate}

\paragraph{The route through the document}
\Cref{sec:course} is a short course on the background.
\Cref{sec:devteam} introduces the running example, a four-agent
development team, and teaches you to read its files.
\Cref{sec:problem} states the problem as seven requirements.
\Cref{sec:mcp,sec:plugins} explain the two technologies the solution is
built from: MCP and \cc{} plugins.
\Cref{sec:arch,sec:host,sec:parts} explain the plugin: its architecture, its
central idea (host mode), and how skills, subagents and hooks enforce the
rules. \Cref{sec:walk1,sec:walk2} walk through two complete sessions.
\Cref{sec:hosted,sec:security} explain the hosted service and the security
model, \Cref{sec:engineering} how the plugin is built and tested, and
\Cref{sec:limits} revisits the requirements and what is still open.
\Cref{sec:exercises} has graded exercises. The appendices hold a reference
of all 24 tools, a ``try it yourself'' guide, the full team of the second
walkthrough, a glossary, and answers.

\paragraph{Four reading paths}
\emph{In 30 minutes:} \Cref{sec:problem}, \Cref{sec:host} up to
\Cref{sec:host-why}, \Cref{fig:arch,fig:hostseq,fig:fptrace,fig:autoloop},
and the green takeaway boxes.
\emph{Hands-on:} \Cref{sec:devteam}, \Cref{app:try}, then
\Cref{sec:walk1} while you run it.
\emph{Security:} \Cref{sec:recap-sandbox,sec:hosted,sec:security}.
\emph{In full:} in order, answering each orange question before checking
\Cref{app:answers}.

\paragraph{Boxes} Blue boxes are \emph{remarks} that sharpen an idea.
Orange boxes (``Check your understanding'') ask you a question; answer it
before reading on. Green boxes summarise the \emph{key takeaways}.
``In plain words'' lines restate a hard idea without jargon.

\paragraph{Colours} Listings and figures use one colour code:

\begin{center}
\small
\begin{tabular}{@{}l p{11.3cm}@{}}
\toprule
\legend{rawbg}{rawbar}{Person} & A file or text a person authors: a task, a command, a configuration file, a skill. In figures: you, and files you edit. \\
\legend{modelbg}{modelbar}{Code} & Source code: Python, rule modules, typed teams. In figures: the plugin and server code. \\
\legend{llmbg}{llmbar}{LLM} & A prompt for the LLM, or its answer. In figures: \cc{} and its subagents. \\
\legend{enginebg}{enginebar}{Engine} & Deterministic output of the engine, checker or server. In figures: the engine and the files only it writes. \\
\legend{wirebg}{wirebar}{Wire} & A protocol message between two programs (line breaks added). \\
\legend{badbg}{badbar}{Refused} & In figures only: an escalated value or a refused request. \\
\bottomrule
\end{tabular}
\end{center}

White boxes in figures are external programs (a browser, \code{curl}).
In protocol listings, \code{->} marks a message from client to server and
\code{<-} one from server to client.

\paragraph{What kind of claim each statement is} Every listing caption
starts with its kind (\Cref{tab:status}).

\begin{table}[!htb]
\centering\small
\caption{Kinds of claims, and how captions mark them.}
\label{tab:status}
\begin{tabularx}{\textwidth}{@{}l Y l@{}}
\toprule
Caption starts with & What it means & Where \\
\midrule
\emph{Observed} & Produced by the real engine, version 0.3.0, in host mode, by the scripts in \file{docs/plugin-walkthroughs/}, which call the same functions the MCP tools call. Outputs are \emph{abridged}: we remove keys and shorten text (marked \ldots) but never change a value. Digests are the real ones. ``Rendered'' means we reformatted JSON as \code{key: value} lines. & \Cref{sec:mcp,sec:parts,sec:walk1,sec:walk2,sec:hosted} \\
\emph{Code excerpt} & Copied from the repository (file named in the caption), sometimes shortened. & throughout \\
\emph{Illustrative} & Made by us to explain: the \code{Wallet} task, the values we typed in place of an LLM, the shape of a protocol message. & marked \\
(no listing) & \emph{Design explanations} say why the code is the way it is. Where the reason is ours rather than the code's or its documentation's, we say so. & throughout \\
\bottomrule
\end{tabularx}
\end{table}

\begin{remark}[Who typed the answers]
In both walkthroughs the values an LLM would normally write (API signatures,
tests, code) were typed by us, some of them deliberately wrong, so that you
can see what the engine does with a wrong answer. Everything the
\emph{engine} printed in response is real, and you can reproduce it with
the scripts in \file{docs/plugin-walkthroughs/} (\Cref{app:try}).
\end{remark}

% ===========================================================================
\section{Background you need: a short course}\label{sec:course}
% ===========================================================================

Skip any subsection you already know.

\subsection{LLMs and agents in one page}\label{sec:llm}

A \emph{large language model} (LLM) is a program that, given some text (the
\emph{prompt}), produces a likely continuation. Asking it once is called
\emph{drawing a sample}. Four properties matter here.

\begin{itemize}
\item \textbf{It is not deterministic.} Two samples for the same prompt can
differ. A setting called \emph{temperature} controls how much: at low
temperature, the same prompt gives nearly the same answer.
\item \textbf{It can be confidently wrong.} A sample can look right and be
false, for example code that compiles but computes the wrong thing.
\item \textbf{It only knows its context.} Everything the LLM can see in one
call (instructions, conversation so far, tool results) is its
\emph{context}. Whatever is not in the context, or in its training, is
invisible to it.
\item \textbf{It can be steered by text only.} You cannot ``forbid'' an LLM
to use something it has seen; you can only avoid showing it.
\end{itemize}

An \emph{agent} is an LLM in a loop with \emph{tools}: the program shows the
LLM a goal and a list of tools (read a file, run a test); the LLM picks an
action; the program performs it and adds the result to the context; and so
on, until the LLM says it is finished. A \emph{team} is several agents whose
outputs become each other's inputs; one agent passing work to the next is a
\emph{hand-off}.

\subsection{A coding assistant is an agent in your terminal}

\cc{} is an agent that works in a software project: it reads and edits
files, runs commands and tests, and talks to you in a terminal or an
editor~\cite{claudecodedocs}. Its built-in tools include \code{Read},
\code{Edit}, \code{Write}, \code{Glob} (find files by name), \code{Grep}
(search file contents) and \code{Bash} (run a shell command). Before it uses
a tool that changes something or that comes from an extension, \cc{}
normally asks you for permission. Three more properties matter here.

\begin{itemize}
\item \textbf{It can be extended} without changing \cc{} itself, by
installing a \emph{plugin} (\Cref{sec:plugins}).
\item \textbf{It can start helpers.} \cc{} can hand a sub-task to a
\emph{subagent}: a second agent with its own instructions, its own (often
smaller) list of allowed tools, and a fresh context. The main agent sees
only the subagent's final report.
\item \textbf{It sees a lot.} By default it can read the whole project and
remembers the whole conversation. Usually a strength; for \tool{} a danger,
because an agent that fills one value must see \emph{only} that value's
inputs (\Cref{sec:recap-agentm2m}). \Cref{sec:parts} shows how the plugin
resolves this tension.
\end{itemize}

In the vocabulary of MCP (\Cref{sec:mcp}), \cc{} is the \emph{host}: the
program in which everything else runs and which owns the LLM.

\subsection{Processes, standard input and output}

When \cc{} starts a helper program, the operating system creates a new
\emph{process}. Every process has three standard streams: \emph{standard
input} (stdin), which it reads; \emph{standard output} (stdout), which it
writes; and \emph{standard error} (stderr), for diagnostics. The parent can
write into the child's stdin and read its stdout, like a pipe in
\code{cat file | grep x}. The \tool{} plugin uses exactly this: \cc{} starts
the \tool{} server as a child process and exchanges JSON messages with it,
one per line, over stdin and stdout.

\subsection{JSON, JSON Schema and remote procedure calls}

\emph{JSON} is a text format for structured data: objects
\code{\{"key": value\}}, arrays \code{[1, 2]}, strings, numbers, booleans
and \code{null}. A \emph{JSON Schema} is a JSON document that says which
JSON documents are allowed, for example ``an object with a required string
field \code{target\_key}''. It is to JSON what a class declaration is to an
object.

A \emph{remote procedure call} (RPC) lets one program call a function in
another program as if it were local: the caller sends the function's name
and arguments; the callee runs it and sends back the result or an error.
\emph{JSON-RPC 2.0}~\cite{jsonrpc2010} is a minimal RPC protocol in JSON
with three kinds of messages (\Cref{lst:jsonrpc}).

\begin{lstlisting}[style=wire, caption={Illustrative: the three kinds of JSON-RPC 2.0 messages.}, label={lst:jsonrpc}]
request       {"jsonrpc": "2.0", "id": 7, "method": "add", "params": {"a": 2, "b": 3}}
response      {"jsonrpc": "2.0", "id": 7, "result": 5}
   or         {"jsonrpc": "2.0", "id": 7, "error": {"code": -32601, "message": "Method not found"}}
notification  {"jsonrpc": "2.0", "method": "log", "params": {"text": "hello"}}    (no id: no answer)
\end{lstlisting}

The \code{id} pairs each response with its request, because several requests
may be in flight at once. A \emph{notification} has no \code{id} and gets no
response.

\begin{think}
Why must the response carry the same \code{id} as its request? What could
go wrong if two requests were outstanding and responses had no \code{id}?
\end{think}

\subsection{HTTP, REST, API keys and hashes}

Programs on \emph{different} machines usually talk over \emph{HTTP}: a
client sends a request (a method such as \code{GET} or \code{POST}, a path
such as \code{/api/auto/run}, \emph{headers} and an optional body) and the
server returns a status code (\code{200} OK, \code{401} unauthorized,
\code{403} forbidden, \code{404} not found, \code{409} conflict, \code{500}
server error, \ldots), headers and a body~\cite{rfc9110}. A \emph{REST
API}~\cite{fielding2000rest} exposes resources as paths and operations as
HTTP methods.

A \emph{cryptographic hash} such as SHA-256 turns any data into a short,
fixed-length string, its \emph{digest}. Equal data always gives the equal
digest; any change, however small, gives a different one; and you cannot
recover the data from the digest. We use digests twice. First, a server
that serves many users identifies them by \emph{API keys}: long random
secrets that the server gives each user once and that the user sends with
every request, usually as the header \code{Authorization: Bearer <key>}. A
careful server stores only the \emph{digest} of each key: when a request
arrives, it hashes the presented key and looks the digest up. If the
database leaks, the keys cannot be recovered. Second, the engine uses
digests to notice that the inputs of a value have changed
(\Cref{sec:recap-agentm2m}).

\subsection{Plugins and inversion of control}

A \emph{plugin} adds behaviour to a host application through
\emph{extension points} that the host defines; you know this from IDE
extensions. The host stays in charge: it decides when to load the plugin and
when to call it (\emph{inversion of control}: ``don't call us, we'll call
you''). A plugin can do nothing the extension points do not allow. \cc{}
offers four extension points (tools through MCP servers, skills, subagents,
hooks) and one distribution channel (marketplaces).

\subsection{\tool{} in one example}\label{sec:recap-agentm2m}

This subsection and the next condense the companion
documents~\cite{agentm2m2026,autom2mteaching}. \Cref{fig:recap} shows the
whole idea on one value.

\begin{figure}[!htb]
\centering
\begin{tikzpicture}[font=\sffamily\small, >={Stealth[length=4pt]},
  obj/.style={draw=modelbar, fill=modelbg, rounded corners=2pt, align=left, inner sep=4pt, font=\sffamily\scriptsize},
  eng/.style={draw=enginebar, fill=enginebg, rounded corners=2pt, align=center, inner sep=4pt, font=\sffamily\scriptsize},
  llm/.style={draw=llmbar, fill=llmbg, rounded corners=2pt, align=center, inner sep=4pt, font=\sffamily\scriptsize},
  lbl/.style={font=\sffamily\scriptsize, align=center}]
  \node[font=\sffamily\scriptsize\bfseries] at (0,1.65) {Analyst's view (\cls{Req})};
  \node[obj, text width=36mm] (s) at (0,0.5) {\textbf{UserStory} S2\\title = ``mark a task done''\\status = accepted\\criteria: S2.1 ``completing a\\done task is rejected''};
  \node[font=\sffamily\scriptsize\bfseries] at (9.6,1.65) {Architect's view (\cls{Arch})};
  \node[obj, text width=36mm] (op) at (9.6,0.5) {\textbf{Operation} op\_s2\\name = ``op\_s2'' \hfill\emph{(engine)}\\signature = ?\hfill\emph{(LLM)}};
  \node[eng, text width=30mm] (rule) at (4.8,1.15) {rule \code{Story2Operation}:\\for every accepted story,\\create one operation};
  \draw[->, thick, enginebar] (s.east |- rule) -- (rule); \draw[->, thick, enginebar] (rule) -- (op.west |- rule);
  \node[llm, text width=30mm] (llm) at (4.8,-1.0) {LLM sees \emph{only} the\\footprint (S2.1) and\\writes the signature};
  \draw[->, dashed, llmbar] (s.south east) -- node[below left, lbl] {footprint} (llm.west);
  \node[eng, text width=36mm] (chk) at (9.6,-1.55) {validator \code{parses()}:\\accept or reject;\\stamp = digest(footprint)};
  \draw[->, llmbar] (llm.east) -- (chk.west |- llm.east);
  \draw[->, enginebar] (chk) -- (op);
  \node[eng, text width=24mm] (tl) at (4.8,-2.55) {trace link\\S2 $\to$ op\_s2};
\end{tikzpicture}
\caption{\tool{} on one value. The engine (green) decides that \code{op\_s2}
exists and fills its name; the LLM (purple) writes only the signature, from
the footprint alone; a validator decides whether that answer is kept; the
engine records a trace link and a stamp.}
\label{fig:recap}
\end{figure}

\paragraph{Views} Each agent of a team owns one \emph{view}: a model that
conforms to a small class diagram, its \emph{metamodel}. Think of the
metamodel as a \emph{form} and the view as the \emph{filled-in form}. Only
the owner may change its view: these are its \emph{write rights} (written
$\omega$ in the \tool{} paper, and \code{omega} in error messages).

\paragraph{Hand-offs are rules} A hand-off from one view to another is a
\emph{model-to-model (M2M) transformation}: a program that reads one model and
produces another. It is written as rules in the style of the ATL
language~\cite{jouault2008atl}: ``for every accepted user story, create one
API operation''. A deterministic \emph{engine} runs the rules: it decides
which target objects exist and how they link. A field computed by a formula
is a \emph{structural binding}. A field no formula can compute (an API
signature, a test) is a \emph{stochastic binding}, written
\code{@llm(prompt, footprint)}: an LLM writes it. In this document the
\emph{binding} is the clause in the rule; the \emph{value} is what gets
written into one target object.

\paragraph{Footprint and validator} The \emph{footprint} is the exact data
the LLM may see when it writes one value. The value is accepted only if its
\emph{validator} (the \code{@check} clause) passes. A \emph{form} validator
checks shape (``does it parse?''); a \emph{behaviour} validator executes the
value (``do the tests pass?''). After $k$ rejections on the same footprint
(here $k=3$) the value is \emph{escalated}: the engine stops asking and
reports it, instead of looping forever.

\paragraph{Trace links, stamps and obligations} Every time a rule fires, the
engine records a \emph{trace link} (``this target came from that source,
through this rule''). With every accepted value it stores a \emph{stamp}:
the digest of the footprint the value was written from. When a source
changes, the values whose stamp no longer equals the digest of their
current footprint are \emph{stale}; redoing them are the
\emph{obligations} of the change. Values whose footprint did not change are
left alone.

\paragraph{Fixpoint and $\varphi$} The engine runs all hand-offs again and
again until a pass creates and changes nothing: that state is a
\emph{fixpoint}. The team is done when the \emph{acceptance predicate}
$\varphi$ holds: every rule match has its target, every value is accepted
for its current footprint, and nothing is pending or escalated. An agent
cannot end the run by claiming success.

\subsection{\auto{} in one example}\label{sec:recap-auto}

In \tool{} a person writes views and rules. In \auto{} a \emph{builder}
LLM writes the whole team design as one JSON document, a \emph{typed team}.
A deterministic \emph{checker} admits it only if six conditions hold
(\Cref{tab:w}); otherwise it returns exact \emph{diagnostics} and the
builder revises.

\begin{table}[!htb]
\centering\small
\caption{The six admission conditions, with an example violation each (all
but \W{3} are real diagnostics from this document).}
\label{tab:w}
\begin{tabularx}{\textwidth}{@{}l L{4.3cm} Y@{}}
\toprule
& Meaning & Example violation \\
\midrule
\W{1} & Well typed: rules read and write only fields that exist; every LLM value has a known validator. & \code{footprint path 'd.method.returnType': Goal!Method has no feature 'returnType'} \\
\W{2} & One writer: each view has one owning agent; each class is created by one rule. & \code{view Test written by ['Developer', 'Tester']} \\
\W{3} & Complete: every mandatory field is filled exactly once. & a rule that never fills the mandatory \code{name} \\
\W{4} & Anchored coverage: every goal flows into a value checked by a \emph{behaviour} validator; the deliverable is behaviour-checked; every view lies on such a flow. & \code{view Code reaches no behavioural check} \\
\W{5} & The engine decides ``done'': $\varphi$ is exactly the four engine clauses; no cycles between views. & \code{done-clause Tester.says('ALL TESTS PASS') is not engine-checkable} \\
\W{6} & Right tools: the owner of a value has the tools its validators need. & \code{Method2Test.code needs ['exec']; owner Tester has []} \\
\bottomrule
\end{tabularx}
\end{table}

The task is a Python class with unimplemented methods; it is turned
(\emph{lifted}) into a fixed \cls{Goal} view with one \cls{Method} object per
method to implement. For an \auto{} team, $\varphi$ is the conjunction of
four clauses:
\begin{itemize}[itemsep=1pt]
\item \code{cover(G)}: every goal object in $G$ (here every method) has a
validated descendant, and every method has an accepted deliverable (code);
\item \code{valid}: every accepted value still passes its validators on the
final models;
\item \code{fresh}: every stamp equals the digest of its current footprint;
\item \code{noObl}: no value is pending or escalated.
\end{itemize}
If $\varphi$ fails, the engine \emph{locates} the failing value by trace
lookup, and a repair, a revised team called a \emph{delta}, must pass the
same checker before it is applied.

\begin{remark}[The one sentence to remember]
In \tool{} \emph{an LLM writes values; the engine owns structure}. In
\auto{} \emph{an LLM proposes the team; the checker owns admission}. In the
plugin \emph{\cc{} does all the LLM work; the engine keeps every decision}.
\end{remark}

\begin{tcolorbox}[enhanced, breakable, colback=accentlight, colframe=accentlight, arc=0pt,
   boxrule=0pt, borderline west={2.2pt}{0pt}{accent}, left=8pt, right=8pt, top=5pt, bottom=5pt]
\textsf{\bfseries\color{accent}Words we use.}\enspace
A \textbf{binding} is a clause in a rule; a \textbf{value} is what it writes
into one object. A value is \textbf{pending} when it is ready to be written
(never accepted, or \textbf{stale}); \textbf{blocked} when its footprint
still contains an upstream value that was never accepted; an
\textbf{obligation} when it is pending because something changed.
\textbf{Engine LLM} means an LLM the engine calls itself through an API; in
\textbf{host mode} there is none, and \cc{} writes the values. An
\textbf{agent} is a member of the team (Analyst, Architect, \ldots); \cc{}
plays the agents.
\end{tcolorbox}

\subsection{Isolation and sandboxes}\label{sec:recap-sandbox}

Behaviour validators \emph{execute} code that an LLM wrote. Executing
untrusted code is dangerous: it could delete files, read secrets, or use the
network. A \emph{sandbox} runs code with restricted powers. A plain
\emph{child process} with resource limits (CPU time, memory, file size)
stops runaway code but can still read every file the user can. Linux
\emph{namespaces} let a process get its own private view of the network, the
process list or the file system; a namespace sandbox such as
bubblewrap~\cite{bubblewrap} uses them to give code no network and only a
few read-only directories, so other users' files simply are not there. The
principle behind both is \emph{least privilege}~\cite{saltzer1975protection}:
give each part of a system only the rights it needs.

\begin{takeaway}
\begin{itemize}
\item An LLM produces samples from its context; an agent is an LLM with
tools; \cc{} is an agent and the \emph{host} in which the plugin runs.
\item Programs talk through pipes (same machine) or HTTP (different
machines); JSON-RPC pairs requests and responses by \code{id}.
\item A digest changes whenever the data changes: servers store digests of
API keys, and the engine stores digests of footprints (stamps).
\item \tool{}: the engine builds structure, the LLM fills footprint-bounded
values, validators decide, stamps make change impact exact, $\varphi$
decides ``done''. \auto{}: a checker admits LLM-proposed teams
(\W{1}--\W{6}).
\end{itemize}
\end{takeaway}

% ===========================================================================
\section{The running example: a four-agent development team}\label{sec:devteam}
% ===========================================================================

The plugin ships three ready-made teams, called \emph{templates}. We use the
first, \devteam{}, the running example of the \tool{} paper, throughout. This
section introduces it and teaches you to read its two kinds of files: the
team description \file{team.yaml} and the rule modules.

\subsection{The team}

\devteam{} builds a small to-do web service (\Cref{fig:devteam}). The
\emph{Analyst} owns the view \cls{Req}: user stories with acceptance
criteria. The \emph{Architect} owns \cls{Arch}: API operations. The
\emph{Developer} owns \cls{Code}: code edits. The \emph{Tester} owns
\cls{Test}: test cases. The starting (\emph{seed}) model has three stories:
\code{S1} \emph{create a task} (criterion \code{S1.1} ``a title is
required''), \code{S2} \emph{mark a task done} (\code{S2.1} ``completing a
done task is rejected''), both accepted, and \code{S3} \emph{export tasks as
CSV} (\code{S3.1} ``the file has a header row''), still a draft.

\begin{figure}[!htb]
\centering
\begin{tikzpicture}[font=\sffamily\small, >={Stealth[length=4pt]},
  ag/.style={draw=llmbar, fill=llmbg, rounded corners=2pt, minimum width=22mm, minimum height=6.5mm, align=center, font=\sffamily\scriptsize},
  vw/.style={draw=modelbar, fill=modelbg, rounded corners=2pt, minimum width=22mm, minimum height=10mm, align=center, font=\sffamily\scriptsize},
  lbl/.style={font=\sffamily\scriptsize, align=center, fill=white, inner sep=1.5pt}]
  \node[ag] (a1) at (0,1.25) {Analyst};
  \node[vw] (v1) at (0,0) {\textbf{Req}\\stories, criteria};
  \node[ag] (a2) at (5.0,1.25) {Architect};
  \node[vw] (v2) at (5.0,0) {\textbf{Arch}\\operations};
  \node[ag] (a3) at (10.0,1.25) {Developer};
  \node[vw] (v3) at (10.0,0) {\textbf{Code}\\code edits};
  \node[ag] (a4) at (5.0,-3.6) {Tester};
  \node[vw] (v4) at (5.0,-2.6) {\textbf{Test}\\test cases};
  \draw[->, thick] (v1) -- node[lbl, above=1pt] {\code{Req2Arch}: signature\\\code{@check parses()}} (v2);
  \draw[->, thick] (v2) -- node[lbl, above=1pt] {\code{Arch2Code}: body\\{\color{badbar}\code{@check compiles()}}} (v3);
  \draw[->, thick] (v1) |- node[lbl, pos=0.25, left=1pt] {\code{Req2Test}: oracle\\\code{@check}\\\code{failsOnStub()}} (v4);
  \node[ag, dashed] (a5) at (10.0,-3.6) {Security Reviewer};
  \node[vw, dashed] (v5) at (10.0,-2.6) {\textbf{Sec}\\reviews};
  \draw[->, dashed] (v2.south east) -- node[lbl, sloped, above=1pt] {\code{Arch2Sec} (added at run time)} (v5.north west);
\end{tikzpicture}
\caption{\devteam{}. Each agent (purple) owns one view (blue). Each arrow is
a hand-off with one LLM-written value and its validator. The red validator
checks only form, a weakness we meet in \Cref{sec:walk1}. The dashed agent
is added while the team runs.}
\label{fig:devteam}
\end{figure}

\subsection{Reading \file{team.yaml}}

The file \file{.agentm2m/team.yaml} describes the views (their classes,
owner and seed model) and lists the hand-offs. \Cref{lst:teamyaml} shows the
part for the Analyst's view.

\begin{lstlisting}[style=raw, caption={Code excerpt: \file{templates/devteam/team.yaml} (abridged).}, label={lst:teamyaml}]
name: devteam
views:
  Req:
    owner: Analyst                       # the only agent that may edit this view
    classes:
      Epic:      {attributes: [name]}
      Criterion: {attributes: [id, text]}
      UserStory:
        attributes: [id, title, status]  # status: draft | accepted
        references:
          epic: Epic
          criteria: {type: Criterion, many: true, containment: true}
    root: {class: ReqModel, slots: {epics: Epic, stories: UserStory}}
    seed:                                # the initial model
      stories:
        - {id: S2, title: mark a task done, status: accepted, epic: "Epic#E1",
           criteria: [{id: S2.1, text: completing a done task is rejected}]}
        ...
  Arch: {owner: Architect, classes: {Operation: {attributes: [name, signature, criteria], ...}}, ...}
  ...
handoffs:
  - {name: Req2Arch,  rule: rules/Req2Arch.agentm2m}
  - {name: Arch2Code, rule: rules/Arch2Code.agentm2m}
  - {name: Req2Test,  rule: rules/Req2Test.agentm2m}
\end{lstlisting}

Every object has a \emph{key} of the form \code{Type\#id}, such as
\code{UserStory\#S2} or \code{Criterion\#S2.1}; the tools use these keys to
name objects. A reference such as \code{epic: "Epic\#E1"} points to another
object by its key. A \emph{containment} reference (\code{criteria}) means the
criteria are part of their story, like fields of an object.

\subsection{Reading a rule module}

\Cref{lst:req2arch} shows the first hand-off in full, and
\Cref{tab:syntax} explains its notation. The expression language inside the
rules is a small subset of OCL, the expression language of model-driven
engineering; read \code{s.criteria} like a field access in Java.

\begin{lstlisting}[language=ATL, style=model, caption={Code excerpt: \file{templates/devteam/rules/Req2Arch.agentm2m}.}, label={lst:req2arch}]
module Req2Arch;
create OUT : Arch from IN : Req;
uses 'helpers.py';

rule Epic2Component {
  from e : Req!Epic
  to  c : Arch!Component ( name <- e.name )
}

rule Story2Operation {
  from s : Req!UserStory
           (s.status = #accepted)                    -- guard: drafts are not handed off
  to  op : Arch!Operation (
    name      <- s.id.toOpName(),                    -- structural
    component <- s.epic,                             -- resolved via trace (Epic -> Component)
    criteria  <- s.criteria.criteriaText(),          -- structural copy, read later by Arch2Code
    signature <- @llm('Derive an API signature for this user story. Answer with exactly one
                       line of the form name(param: Type, ...) -> ReturnType and nothing else.',
                       s.criteria),                  -- footprint
    @check signature.parses() and signature.params()->notEmpty() )
}
\end{lstlisting}

\begin{table}[!htb]
\centering\small
\caption{The notation of a rule module.}
\label{tab:syntax}
\begin{tabularx}{\textwidth}{@{}L{4.6cm} Y@{}}
\toprule
Notation & Meaning \\
\midrule
\code{module Req2Arch;} & The name of the hand-off. \\
\code{create OUT : Arch from IN : Req;} & It writes view \cls{Arch} and reads view \cls{Req}. A hand-off may read several views. \\
\code{uses 'helpers.py';} & Load Python helper functions from this file. \\
\code{rule R \{ from s : Req!UserStory} & For every object \code{s} of class \cls{UserStory} in view \cls{Req} \ldots \\
\code{(s.status = \#accepted)} & \ldots that satisfies this \emph{guard} (\code{\#accepted} is an enumeration literal) \ldots \\
\code{to op : Arch!Operation (...)} & \ldots create one \cls{Operation} \code{op} in view \cls{Arch}, and fill its fields: \\
\code{name <- s.id.toOpName()} & a \emph{structural binding}: the engine computes the value. \code{toOpName} is a Python function from \file{helpers.py}, called as if it were a method of \code{s.id}. \\
\code{component <- s.epic} & a reference: the engine replaces the source epic by the component it created from it (found through the trace). \\
\code{signature <- @llm(p, f)} & a \emph{stochastic binding}: an LLM writes the value from prompt \code{p} and footprint \code{f}. \\
\code{@check e} & the validator of the binding just above it; \code{e} must be true. \\
\code{x->notEmpty()} & a collection operation (as in OCL): ``\code{x} is not empty''. \\
\bottomrule
\end{tabularx}
\end{table}

Validators are ordinary Python functions in \file{helpers.py}. A good
validator does not just return \code{False}; it returns a
\code{Rejected(reason)} object, so that whoever tries again learns what was
wrong (\Cref{lst:helpers}).

\begin{lstlisting}[language=Python, style=model, caption={Code excerpt: \file{templates/devteam/rules/helpers.py} (abridged).}, label={lst:helpers}]
def parses(signature: str):
    if signature_parses((signature or "").strip()):
        return True
    return Rejected("expected exactly one line: name(param: Type, ...) -> ReturnType")

def compiles(body: str):
    if python_compiles(_unfence(body)):
        return True
    return Rejected("not valid Python")
\end{lstlisting}

The other two hand-offs are similar. \code{Arch2Code} creates one
\cls{CodeEdit} per operation; its \code{body} is written by an LLM whose
footprint is the operation's signature \emph{and} the criteria that
\code{Req2Arch} copied onto it, and whose validator is \code{compiles()}.
\code{Req2Test} has no guard: it creates one \cls{TestCase} per criterion,
drafts included; its \code{oracle} is written from \code{c.text} and checked
by \code{failsOnStub()}. That validator runs the oracle against a stub
implementation that only raises \code{NotImplementedError} and accepts the
oracle if it raises. It thus rejects an oracle that never calls the
implementation; it does \emph{not} reject one that calls it but asserts
nothing, because the stub's exception already makes it fail.

\subsection{The other two templates}

\Cref{tab:templates} summarises all three templates. Two features of the
others are worth knowing. The \code{research} template has a hand-off that
reads \emph{two} views in one rule (\code{from p : Experiments!ExperimentPlan,
c : Literature!Claim (p.topic = c.topic)}), pairing every plan with every
claim of the same topic. The \code{incident} template uses a \emph{Lift}
binding, \code{self <- @llm(...)}: the LLM answers with one JSON object, and
the engine writes its keys onto several attributes at once, rejecting the
answer unless every key names a real attribute. It also has an
\emph{executable-oracle} validator, \code{passesDryRun}, which runs an
LLM-written Python script and accepts it only if it exits with code 0.

\begin{table}[!htb]
\centering\small
\caption{The three templates of \mcptool{team\_init}.}
\label{tab:templates}
\rowcolors{2}{black!3}{white}
\begin{tabularx}{\textwidth}{@{}l L{4.4cm} Y@{}}
\toprule
\rowcolor{white}Template & Agents (views) & What it teaches \\
\midrule
\code{devteam} & Analyst (Req), Architect (Arch), Developer (Code), Tester (Test); Security Reviewer added at run time & Guards, blocked values, exact change impact, team evolution; a form-only validator on code. \\
\code{research} & Literature Reviewer, Experiment Designer, Report Writer & A multi-source hand-off (one rule reads two views). \\
\code{incident} & Monitor, Triage, Remediation, Postmortem & A Lift binding (one JSON answer fills several fields); an executable-oracle validator that runs LLM-written scripts. \\
\bottomrule
\end{tabularx}
\end{table}

\begin{think}
In \devteam{}, \code{Story2Operation} will match 2 stories but
\code{Criterion2TestCase} 3 criteria. Using \Cref{lst:req2arch} and the
description of \code{Req2Test}, explain the difference.
\end{think}

\begin{takeaway}
\begin{itemize}
\item \file{team.yaml} declares views (classes, owner, seed) and hand-offs;
objects are named by keys such as \code{Criterion\#S2.1}.
\item A rule says: for every source object matching the pattern and guard,
create one target object; structural bindings are computed, \code{@llm}
bindings are written by an LLM from their footprint, \code{@check} decides.
\item Validators are Python functions that return \code{True} or
\code{Rejected(reason)}.
\end{itemize}
\end{takeaway}

% ===========================================================================
\section{The problem: from a research library to a tool}\label{sec:problem}
% ===========================================================================

\subsection{Where we started}

Before version 0.2.0, \tool{} could be used only as a Python library. To
run a team you wrote Python that built metamodels, created seed models,
registered hand-offs and called the runtime (the files
\file{examples/*/run.py} still show this). All state lived in memory and
vanished when the program ended. The engine called an LLM itself, through an
API key for OpenAI or Anthropic or through a local Ollama server (a program
that runs open LLMs on your own machine). None of this fits how a developer
works: in an editor, with an assistant, over days, on a project that
changes.

\subsection{Seven requirements}

We phrase the problem as requirements, as you would in a requirements
engineering course. \Cref{tab:req} says where the design answers each, and
\Cref{sec:limits} checks how well it did.

\begin{description}[leftmargin=1.2em, style=nextline, itemsep=3pt]
\item[R1 Usable from the assistant.] A developer must be able to start,
run, change and inspect a team by talking to \cc{}, without writing Python.
\item[R2 No second LLM bill.] The developer already pays for \cc{}'s LLM.
Requiring a second API key so that the engine can call an LLM would be a
barrier. Yet the engine's guarantees (validators, escalation, stamps,
$\varphi$) must not weaken.
\item[R3 Footprint discipline survives.] Change impact is exact only if
each value was written from its footprint and nothing else. An assistant
that has read the whole repository could quietly use more. The design must
make the footprint the \emph{only} information available, not just ask
politely.
\item[R4 State survives.] Models, trace links and stamps must persist
across sessions and survive several processes touching them.
\item[R5 Mistakes are caught early.] When the assistant edits a rule file
or a typed team, errors must be reported immediately, before anything runs.
\item[R6 Usable without installing the engine.] Teams that share a server,
or users who cannot install Python packages, should be able to use the same
tools remotely.
\item[R7 Safe.] The engine executes code: helpers that the project provides,
and code that the LLM writes. A shared server must not let one user's code
see another's data.
\end{description}

\begin{table}[!htb]
\centering\small
\caption{Requirements and where the design answers them.}
\label{tab:req}
\begin{tabularx}{\textwidth}{@{}l Y l@{}}
\toprule
Req. & Design answer & Section \\
\midrule
R1 & An MCP server exposes 24 tools; skills script the common workflows. & \ref{sec:mcp}, \ref{sec:parts}, \ref{app:tools} \\
R2 & \emph{Host mode}: the engine never samples; \cc{} answers engine-built prompts, and the engine judges the answers exactly as its own samples. & \ref{sec:host} \\
R3 & Subagents whose tool allowlist contains only ``get my next prompt'' and ``submit my answer''; values tied to a \code{footprint\_version}. & \ref{sec:host}, \ref{sec:parts} \\
R4 & A workspace directory \file{.agentm2m/} with atomic writes, a file lock and reload-on-change. & \ref{sec:arch} \\
R5 & A hook re-validates after every edit of a team file and feeds errors back to \cc{}. & \ref{sec:parts} \\
R6 & \code{agentm2m serve}: the same tools over HTTP, API keys, per-key projects. & \ref{sec:hosted} \\
R7 & Write rights, path checks, one sandbox function, and a service that refuses to execute code without an isolating sandbox. & \ref{sec:security} \\
\bottomrule
\end{tabularx}
\end{table}

\begin{takeaway}
\begin{itemize}
\item Turning the engine into a tool is a software engineering problem with
explicit requirements, not just packaging.
\item The two hardest requirements pull against each other: use \cc{}'s LLM
(R2) without letting \cc{}'s broad knowledge leak into values (R3).
\end{itemize}
\end{takeaway}

% ===========================================================================
\section{The Model Context Protocol (MCP)}\label{sec:mcp}
% ===========================================================================

\subsection{The problem MCP solves}

Suppose there are $N$ AI assistants and $M$ tools (a database, an issue
tracker, \tool{}). If every assistant integrates every tool in its own way,
someone writes $N\times M$ adapters. If all speak one protocol, each side
writes one adapter: $N+M$. The Language Server Protocol~\cite{lsp} did this
for editors and programming languages: any editor can use any language
server. The \emph{Model Context Protocol}
(MCP)~\cite{mcp2024,mcpspec20250618} does it for AI assistants: it
standardises how an assistant discovers what a tool server offers and how it
calls it. It is built on JSON-RPC 2.0.

\subsection{Roles: host, client, server}

MCP names three roles (\Cref{fig:mcproles}). The \textbf{host} is the
application the user works in; here, \cc{}. Inside the host, one
\textbf{client} manages the connection to one server. The \textbf{server}
is the program that offers capabilities; here, the program
\code{agentm2m-mcp} (\file{src/agentm2m/mcp_server.py}). The LLM in the host
sees the tools of all connected servers.

\begin{figure}[!htb]
\centering
\begin{tikzpicture}[font=\sffamily\small, >={Stealth[length=4pt]},
  llm/.style={draw=llmbar, fill=llmbg, rounded corners=2pt, align=center, minimum height=9mm, inner sep=4pt},
  cli/.style={draw=modelbar, fill=modelbg, rounded corners=2pt, align=center, minimum height=9mm, inner sep=4pt},
  ext/.style={draw=extbar, fill=extbg, rounded corners=2pt, align=center, minimum height=9mm, inner sep=4pt},
  lbl/.style={font=\sffamily\scriptsize, align=center, fill=white, inner sep=1.5pt}]
  \node[llm, minimum width=22mm] (llm) at (0,0) {LLM\\(Claude)};
  \node[cli, minimum width=22mm] (c1) at (3.2,0.8) {MCP client 1};
  \node[cli, minimum width=22mm] (c2) at (3.2,-0.8) {MCP client 2};
  \begin{scope}[on background layer]
    \node[draw=llmbar, dashed, rounded corners=3pt, fit=(llm)(c1)(c2), inner sep=7pt,
          label={[font=\sffamily\scriptsize\bfseries, text=llmbar]above:Host: \cc{}}] {};
  \end{scope}
  \node[cli, minimum width=32mm] (s1) at (10.0,0.8) {\code{agentm2m-mcp}\\(24 tools)};
  \node[ext, minimum width=32mm] (s2) at (10.0,-0.8) {another server\\(e.g.\ a database)};
  \draw[<->] (c1) -- node[lbl, above=2pt] {JSON-RPC over stdio} (s1);
  \draw[<->] (c2) -- node[lbl, below=2pt] {JSON-RPC over HTTP} (s2);
  \draw[<->] (llm) -- (c1); \draw[<->] (llm) -- (c2);
\end{tikzpicture}
\caption{MCP roles. The host contains one client per server. The LLM decides
which tool to call; the host performs the call through the right client.}
\label{fig:mcproles}
\end{figure}

\subsection{What a server can offer}

MCP defines three kinds of capabilities, called \emph{primitives}.
\textbf{Tools} are functions the LLM may decide to call, such as
\mcptool{run} or \mcptool{submit\_binding}. \textbf{Resources} are pieces of
data the host may read, such as a file. \textbf{Prompts} are reusable prompt
templates the user may pick. A server also sends free-text
\textbf{instructions} when the connection opens, telling the LLM how to use
it. The \tool{} server uses tools and instructions only (the SDK advertises
empty resource and prompt lists). Its instructions summarise the intended
order of calls and the one rule that matters most: ``Answer a binding ONLY
from the prompt \mcptool{next\_bindings} returns''.

\subsection{The life of a connection}

A connection goes through four steps. We captured them by starting the real
server and writing JSON lines into its stdin (script
\file{mcp_stdio_demo.py}).

\what{1. Initialize.} The client says which protocol version it speaks and
what it can do. The server answers with the version both will use (here it
accepts the client's), its capabilities and its instructions.

\begin{lstlisting}[style=wire, caption={Observed: the \code{initialize} handshake with \code{agentm2m-mcp} over stdio (instructions abridged).}, label={lst:init}]
-> {"jsonrpc":"2.0","id":1,"method":"initialize","params":{"protocolVersion":"2025-06-18",
    "capabilities":{},"clientInfo":{"name":"demo","version":"1"}}}
<- {"jsonrpc":"2.0","id":1,"result":{"protocolVersion":"2025-06-18",
    "capabilities":{"prompts":{"listChanged":false},"resources":{"listChanged":false,"subscribe":false},
                    "tools":{"listChanged":false}},
    "serverInfo":{"name":"agentm2m","version":""},
    "instructions":"agentm2m runs a team of agents whose hand-offs are model-to-model
      transformations. ... Typical loop: team_status -> run -> next_bindings ->
      submit_binding (repeat) -> acceptance. ... Answer a binding ONLY from the prompt
      next_bindings returns: it already contains the whole allowed footprint. ..."}}
-> {"jsonrpc":"2.0","method":"notifications/initialized"}
\end{lstlisting}

(The empty \code{serverInfo.version} is harmless: the server does not tell
the SDK its version. \code{agentm2m-mcp --version} prints it.)

\what{2. Initialized.} The client confirms with a \emph{notification}. Only
now may normal requests flow.

\what{3. Discover.} The client asks \code{tools/list}. For each tool, the
server returns its name, a description and a JSON Schema of its arguments
(\Cref{lst:list}).

\begin{lstlisting}[style=wire, caption={Observed: the \code{tools/list} entry for \code{submit\_binding}, one of 24.}, label={lst:list}]
{"name": "submit_binding",
 "description": "Host mode: submit a value for one pending binding, passing the footprint_version
   that came with its prompt. Status is accepted, rejected (with reason and a retry_prompt;
   resubmit), escalated (validator kept rejecting: needs a source change or a human),
   already_accepted, or stale (the footprint changed since the prompt was issued: fetch it again).",
 "inputSchema": {"type": "object", "title": "submit_bindingArguments",
   "properties": {"target_key": {"title": "Target Key", "type": "string"},
                  "binding": {"title": "Binding", "type": "string"},
                  "value": {"title": "Value", "type": "string"},
                  "footprint_version": {"title": "Footprint Version", "type": "string"}},
   "required": ["target_key", "binding", "value", "footprint_version"]}}
\end{lstlisting}

The description is the \emph{only} documentation the LLM gets about a tool,
so it is written for the LLM: it names every outcome and what to do next.
Nobody wrote the schema by hand. The Python MCP SDK derives it from the
function's type hints, and the description from its docstring
(\Cref{lst:tooldef}); change the signature and the schema follows. The
\code{@mcp.tool()} line is a Python \emph{decorator}: a function that wraps
the function below it, here to register it as a tool.

\begin{lstlisting}[language=Python, style=model, caption={Code excerpt: how a tool is declared (\file{src/agentm2m/mcp_server.py}).}, label={lst:tooldef}]
@mcp.tool()
@_forwardable
def submit_binding(target_key: str, binding: str, value: str, footprint_version: str,
                   ctx: Context | None = None) -> dict:
    """Host mode: submit a value for one pending binding, passing the
    footprint_version that came with its prompt. Status is accepted, rejected ..."""
    return _call(_ws(ctx).submit_binding, target_key, binding, value, footprint_version)
\end{lstlisting}

The SDK fills the parameter \code{ctx} with information about the current
request (for example HTTP headers) and leaves it out of the schema.
\code{@\_forwardable} is a second decorator, explained in
\Cref{sec:forward}.

\what{4. Call.} The LLM decides to call a tool; the client sends
\code{tools/call} with the name and arguments; the server runs the function
and returns its result as text content (\Cref{lst:call}).

\begin{lstlisting}[style=wire, caption={Observed: a successful call (JSON text inside \code{content} abridged).}, label={lst:call}]
-> {"jsonrpc":"2.0","id":3,"method":"tools/call",
    "params":{"name":"team_init","arguments":{"template":"devteam"}}}
<- {"jsonrpc":"2.0","id":3,"result":{"isError":false,"content":[{"type":"text",
    "text":"{\n  \"team\": \"devteam\",\n  \"backend\": \"host\",\n  \"agents\": {...} ..."}]}}
\end{lstlisting}

\subsection{Two kinds of errors}

\Cref{lst:errors} shows three failed requests side by side.

\begin{lstlisting}[style=wire, caption={Observed: a protocol error and two tool errors.}, label={lst:errors}]
-> {"jsonrpc":"2.0","id":2,"method":"tools/run","params":{}}
<- {"jsonrpc":"2.0","id":2,"error":{"code":-32601,"message":"Method not found","data":"tools/run"}}

-> {"jsonrpc":"2.0","id":3,"method":"tools/call","params":{"name":"no_such_tool","arguments":{}}}
<- {"jsonrpc":"2.0","id":3,"result":{"isError":true,"content":[{"type":"text",
    "text":"Unknown tool: no_such_tool"}]}}

-> {"jsonrpc":"2.0","id":4,"method":"tools/call","params":{"name":"impact","arguments":
    {"view":"Req","as_agent":"Nobody","ops":[{"op":"delete","key":"UserStory#S3"}]}}}
<- {"jsonrpc":"2.0","id":4,"result":{"isError":true,"content":[{"type":"text",
    "text":"Error executing tool impact: unknown agent 'Nobody'; agents: ['Analyst',
            'Architect', 'Developer', 'Tester']"}]}}
\end{lstlisting}

The first is a \emph{protocol error}: \code{tools/run} is not an MCP method,
so the answer is a JSON-RPC \code{error}. That is a bug in a program and goes
to the host. The other two are \emph{tool errors}: the protocol worked, the
\emph{tool} failed, which the server signals with a normal \code{result} that
has \code{"isError": true} and a readable message. A tool error is
information for the LLM, which can read it and correct itself, here by
naming a real agent. The \tool{} server turns every expected failure into
such a sentence (\Cref{lst:callwrap}). Unexpected exceptions are reported
with their Python type name; that reveals a little about the internals,
which is acceptable for a tool that runs on your own machine or behind an
API key.

\begin{lstlisting}[language=Python, style=model, caption={Code excerpt: expected failures become readable tool errors (\file{mcp_server.py}).}, label={lst:callwrap}]
def _call(fn, *args, **kwargs) -> dict:
    try:
        return fn(*args, **kwargs)
    except (WorkspaceError, SandboxRefused) as exc:
        # A deliberate, actionable tool error (shown to the model verbatim),
        # not a crash with a traceback.
        raise ToolError(str(exc)) from None
    except Exception as exc:  # e.g. an OCL error in a user's rule module
        raise ToolError(f"{type(exc).__name__}: {exc}") from exc
\end{lstlisting}

\subsection{Transports: stdio and Streamable HTTP}

MCP messages travel one of two ways.

\begin{itemize}
\item \textbf{stdio.} The host starts the server as a child process and
exchanges one JSON message per line over its stdin and stdout. Nothing goes
over the network; the server lives as long as the host session. This is how
the plugin runs by default.
\item \textbf{Streamable HTTP.} The server is a web service with one
endpoint (here \code{/mcp}). The client sends every message as an HTTP
\code{POST}; the server answers with plain JSON or with a short stream of
\emph{server-sent events} (a way for one HTTP response to deliver several
messages one after another). The server may assign a session id in the
header \code{Mcp-Session-Id} when the connection is initialised; the
\tool{} service does, and the client then repeats it on every request,
together with \code{MCP-Protocol-Version}. (The specification also uses
\code{GET} to open a stream from the server and \code{DELETE} to end a
session.) Authentication is ordinary HTTP, here an API key in the
\code{Authorization} header. The hosted service uses this transport
(\Cref{sec:hosted}).
\end{itemize}

\subsection{Tool names inside \cc{}, and things called ``agentm2m''}

\cc{} gives every MCP tool a globally unique name, so that tools of different
servers cannot clash: \code{mcp\_\_<server>\_\_<tool>}. A server that comes
with a plugin also carries the plugin's name. Many things in this project
are called \code{agentm2m}, which makes the resulting names look odd;
\Cref{tab:names} sorts them out.

\begin{table}[!htb]
\centering\small
\caption{Things called \code{agentm2m}.}
\label{tab:names}
\begin{tabularx}{\textwidth}{@{}l Y@{}}
\toprule
Name & What it is \\
\midrule
\code{agentm2m} (Python package) & The engine, version 0.3.0. \\
\code{agentm2m} (command) & Its command-line interface: \code{agentm2m workspace \ldots}, \code{agentm2m auto \ldots}, \code{agentm2m serve}. \\
\code{agentm2m-mcp} (command) & The program that runs the MCP server over stdio. \\
\code{agentm2m} (MCP server name) & The name under which the plugin's \file{.mcp.json} registers that server. \\
\code{agentm2m} (plugin) & The \cc{} plugin in \file{plugin/agentm2m/}; its skills are \code{/agentm2m:<skill>}. \\
\code{agentm2m-dev} (marketplace) & The development marketplace in \file{plugin/.claude-plugin/}. \\
\bottomrule
\end{tabularx}
\end{table}

So \mcptool{submit\_binding} appears to \cc{} as
\code{mcp\_\_plugin\_agentm2m\_agentm2m\_\_submit\_binding}: the prefix
\code{mcp}, then \code{plugin\_} + plugin name \code{agentm2m} + \code{\_} +
server name \code{agentm2m}, then the tool name, joined by double
underscores. If you connect the hosted service directly, without the plugin
(\Cref{sec:hosted}), the same tool is \code{mcp\_\_agentm2m\_\_submit\_binding}.
These names matter because subagent allowlists (\Cref{sec:parts}) must spell
them exactly.

\begin{remark}[What MCP does and does not give you]
MCP makes a tool \emph{discoverable} and \emph{callable} with well-formed
arguments. It says nothing about how a \emph{team} of agents fits together:
who owns which artefact, whether every requirement reaches a test, who
decides that the work is done. The \auto{} teaching edition lists MCP among
the approaches that make individual messages well formed but check no
team-level property~\cite{autom2mteaching}. Here MCP is the
\emph{transport}; the team-level guarantees come from the engine behind it.
\end{remark}

\begin{think}
The server's \code{instructions} and each tool's description are plain
text read by the LLM, not by a program. What could go wrong if a tool's
description said ``returns the accepted value'' while the tool actually
returns a status object? Who would notice?
\end{think}

\begin{takeaway}
\begin{itemize}
\item MCP standardises how a host discovers and calls the tools of a server,
over JSON-RPC 2.0: \code{initialize}, \code{initialized},
\code{tools/list}, \code{tools/call}.
\item Tool schemas are generated from Python type hints; descriptions are
written for the LLM and name every outcome.
\item Tool failures are results with \code{isError: true} that the LLM can
act on; protocol errors are JSON-RPC errors, i.e.\ bugs.
\item stdio for a local server; Streamable HTTP with an API key for a
remote one.
\end{itemize}
\end{takeaway}

% ===========================================================================
\section{\cc{} plugins}\label{sec:plugins}
% ===========================================================================

A \cc{} plugin is a directory with a fixed layout~\cite{claudecodedocs}.
\Cref{lst:tree} shows the \tool{} plugin.

\begin{lstlisting}[style=raw, caption={Code excerpt: the layout of \file{plugin/agentm2m/}, version 0.3.0.}, label={lst:tree}]
plugin/agentm2m/
  .claude-plugin/plugin.json      manifest: name, version, description, license, keywords
  .mcp.json                       which MCP server(s) to start, and how
  skills/<name>/SKILL.md          10 skills: init, run, change, evolve, status, author-handoff,
                                  auto-build, check, diagnose, agentm2m-concepts
  agents/<name>.md                4 subagents: binding-worker, handoff-architect,
                                  team-builder, auto-binding-worker
  hooks/hooks.json                2 hooks: SessionStart, PostToolUse (Edit|Write|MultiEdit)
  hooks/session-start.sh, hooks/validate-on-edit.sh
  README.md, CHANGELOG.md, LICENSE
plugin/.claude-plugin/marketplace.json   a development marketplace listing this plugin
\end{lstlisting}

(Plugins may contain more, for example \file{commands/}; this one does not.)

\subsection{The manifest}

\file{plugin.json} names the plugin and its version. The name
(\code{agentm2m}) becomes the namespace of everything the plugin adds: its
skills are invoked as \skill{run}, its tools are prefixed
\code{mcp\_\_plugin\_agentm2m\_\ldots}. The version, 0.3.0, matches the
engine's; \Cref{sec:engineering} explains why.

\subsection{The MCP server configuration}

\file{.mcp.json} tells \cc{} how to start the server (\Cref{lst:mcpjson}).

\begin{lstlisting}[style=raw, caption={Code excerpt: \file{plugin/agentm2m/.mcp.json}.}, label={lst:mcpjson}]
{"agentm2m": {
   "command": "uvx",
   "args": ["--from", "${AGENTM2M_ENGINE:-agentm2m==0.3.0}", "--with", "mcp>=1.2", "agentm2m-mcp"],
   "env": {"AGENTM2M_LLM":     "${AGENTM2M_LLM:-host}",
           "AGENTM2M_URL":     "${AGENTM2M_URL:-}",
           "AGENTM2M_API_KEY": "${AGENTM2M_API_KEY:-}",
           "AGENTM2M_PROJECT": "${AGENTM2M_PROJECT:-}"}}}
\end{lstlisting}

Read it in three parts.

\begin{itemize}
\item \code{uvx} is a command of the Python package manager
uv~\cite{uvdocs}. \code{uvx --from P C} installs package \code{P} into a
cached, isolated environment (once) and runs its command \code{C}. The
plugin therefore contains no engine code: the engine is an ordinary Python
package, and \code{--with mcp>=1.2} adds the MCP SDK to its environment.
\item \code{\$\{AGENTM2M\_ENGINE:-agentm2m==0.3.0\}} is shell-style
substitution: ``the value of the environment variable \code{AGENTM2M\_ENGINE},
or \code{agentm2m==0.3.0} if it is unset''. The default will fetch version
0.3.0 from the Python Package Index once it is published; until then (see
\Cref{sec:engineering}) you must set \code{AGENTM2M\_ENGINE} to the path of
a checkout of the repository.
\item \code{env} passes four settings to the server: who fills values
(\code{AGENTM2M\_LLM}, default \code{host}, \Cref{sec:host}) and, if set,
the URL, API key and project of a hosted service (\Cref{sec:hosted}).
\end{itemize}

\subsection{Skills}

A \emph{skill} is a Markdown file of instructions with a short YAML header
(\emph{frontmatter}). The user can invoke it as a slash command
(\skill{run}), and \cc{} also loads it by itself when the skill's
description matches the task. Whatever the user types after the command
replaces the placeholder \code{\$ARGUMENTS}. A skill is a \emph{procedure
for the LLM}, in plain language: which tools to call, in which order, what to
check, what to report (\Cref{lst:skill}).

\begin{lstlisting}[style=raw, caption={Code excerpt: the start of \file{skills/change/SKILL.md}.}, label={lst:skill}]
---
name: change
description: Apply a change to an agentm2m team's source view (e.g. a tightened acceptance
  criterion or a new user story), preview exactly which downstream artifacts it obliges to be
  redone, then propagate it. Use when requirements or any upstream artifact of the team changes.
argument-hint: "<describe the change>"
---
# Change and propagate
1. Call `team_status` and `model_show` for the view to change. Identify the owning agent ...
2. Turn `$ARGUMENTS` into `model_edit` ops (`create` / `set` / `delete`, element keys like
   `Criterion#S2.1`). Do not apply them yet.
3. Call `impact` with `view`, `ops`, and `as_agent` to preview. Show the user: ...
\end{lstlisting}

\code{argument-hint} is what \cc{} shows when you start typing the command.
One skill, \code{agentm2m-concepts}, sets \code{user-invocable: false}: it is
not a command but background knowledge (``\tool{} in one page'') that \cc{}
loads by itself in projects with a \file{.agentm2m/} directory.

\subsection{Subagents}

A file in \file{agents/} defines a subagent: a name, a description (when to
use it), a \code{model} (\code{inherit}: the same model as the main agent)
and a \code{tools} \emph{allowlist}. A subagent can call only the tools on
its list. \Cref{sec:parts} shows how the plugin uses this.

\subsection{Hooks}

A \emph{hook} is a shell command that \cc{} runs automatically at a fixed
moment, an \emph{event}. The plugin registers two (\Cref{lst:hooks}): one
when a session starts, and one after every use of the \code{Edit},
\code{Write} or \code{MultiEdit} tools. \code{\$\{CLAUDE\_PLUGIN\_ROOT\}} is
replaced by the directory where the plugin is installed. A hook talks back
through its output and exit code: for \code{SessionStart}, what it prints is
added to \cc{}'s context; for \code{PostToolUse}, exit code 2 means
``something is wrong'', and what it wrote to stderr is shown to \cc{}, which
can then fix it. Note that edits made through \code{Bash} (for example with
\code{sed}) do not match the hook.

\begin{lstlisting}[style=raw, caption={Code excerpt: \file{hooks/hooks.json} (abridged).}, label={lst:hooks}]
{"hooks": {
   "SessionStart": [{"hooks": [{"type": "command", "timeout": 10,
        "command": "bash \"${CLAUDE_PLUGIN_ROOT}/hooks/session-start.sh\""}]}],
   "PostToolUse":  [{"matcher": "Edit|Write|MultiEdit",
        "hooks": [{"type": "command", "timeout": 120,
        "command": "bash \"${CLAUDE_PLUGIN_ROOT}/hooks/validate-on-edit.sh\""}]}]}}
\end{lstlisting}

\subsection{Marketplaces and installation}

Plugins are distributed through \emph{marketplaces}: a file
\file{.claude-plugin/marketplace.json} that lists plugins and where their
directories are. A user adds a marketplace once and installs plugins from it
as \code{plugin@marketplace}. The repository ships a development
marketplace, \code{agentm2m-dev}; the script
\file{plugin/scripts/dev-install.sh} registers it and installs the plugin
from your checkout. \Cref{app:try} gives the complete steps to try it.

\begin{table}[!htb]
\centering\small
\caption{\cc{}'s extension points and what the \tool{} plugin puts in each.}
\label{tab:ext}
\begin{tabularx}{\textwidth}{@{}l Y Y@{}}
\toprule
Extension point & What it is & In the \tool{} plugin \\
\midrule
MCP server & Tools the LLM can call & \code{agentm2m-mcp}: 24 tools (12 for hand-written teams, 12 for \auto{}) \\
Skills & Procedures in plain language & 10 skills, 9 of them slash commands \\
Subagents & Helpers with their own tool allowlists & 4 subagents; two of them can only fetch and submit values \\
Hooks & Commands run at fixed events & announce a team at session start; re-validate after edits \\
\midrule
(Distribution) & Marketplace & \code{agentm2m-dev} (local); a public one at release \\
\bottomrule
\end{tabularx}
\end{table}

\begin{think}
Which extension point would you use for each of these? (a) After every
edit of a rule file, check that it still parses. (b) Teach \cc{} the steps
of propagating a requirement change. (c) Guarantee that the helper which
writes test oracles cannot open the project's source files. (d) Let the LLM
preview the impact of a change.
\end{think}

\begin{takeaway}
\begin{itemize}
\item A plugin is a directory: a manifest, an MCP server configuration,
skills, subagents and hooks; the engine itself is a separate Python package.
\item \code{uvx} installs the pinned engine and runs its MCP server;
\code{AGENTM2M\_ENGINE} points it at a local checkout instead.
\item Skills are procedures for the LLM; subagents have tool allowlists;
hooks run at fixed events and report back through exit code 2.
\end{itemize}
\end{takeaway}

% ===========================================================================
\section{Architecture of the plugin}\label{sec:arch}
% ===========================================================================

\subsection{Layers}

\Cref{fig:arch} shows the layers from the user down to the files on disk.
Each layer talks only to the one below it: a \emph{layered architecture},
as in networking or web applications.

\begin{figure}[!htb]
\centering
\begin{tikzpicture}[font=\sffamily\small, >={Stealth[length=4pt]},
  lay/.style={draw=#1bar, fill=#1bg, rounded corners=2pt, minimum height=9mm, text width=7.9cm, align=left, inner sep=5pt},
  side/.style={draw=#1bar, fill=#1bg, rounded corners=2pt, minimum height=8mm, text width=2.9cm, align=left, inner sep=4pt, font=\sffamily\scriptsize},
  lbl/.style={font=\sffamily\scriptsize, align=center, fill=white, inner sep=1.5pt}]
  \node[lay=raw] (u) at (0,0) {\textbf{You} type \skill{run} or ``tighten S2.1''.};
  \node[lay=llm] (cc) at (0,-1.3) {\textbf{\cc{}}: reads skills, dispatches subagents, runs hooks, calls MCP tools.};
  \node[lay=wire] (w) at (0,-2.6) {\textbf{MCP over stdio}: JSON-RPC lines to the child process \code{agentm2m-mcp}.};
  \node[lay=model] (srv) at (0,-3.9) {\textbf{MCP server} \file{mcp_server.py}: 24 thin tool functions.};
  \node[lay=model] (svc) at (0,-5.2) {\textbf{Service layer}: \code{Workspace}, \code{AutoWorkspace}; every operation returns a JSON-able dictionary.};
  \node[lay=engine] (eng) at (0,-6.5) {\textbf{Engine}: rules, matcher, executor, trace, runtime, checker, compiler, sandbox.};
  \node[draw=rawbar, fill=rawbg, rounded corners=2pt, minimum height=9mm, text width=4.0cm, align=left, inner sep=5pt] (f1) at (-2.3,-7.85) {\textbf{You edit}: \file{team.yaml}, \file{rules/}};
  \node[draw=enginebar, fill=enginebg, rounded corners=2pt, minimum height=9mm, text width=4.0cm, align=left, inner sep=5pt] (f2) at (2.3,-7.85) {\textbf{Engine only}: \file{state/}, \file{auto/*.json}};
  \foreach \a/\b in {u/cc, cc/w, w/srv, srv/svc, svc/eng} \draw[<->] (\a) -- (\b);
  \draw[<->] (eng) -- (f1); \draw[<->] (eng) -- (f2);
  \node[side=model] (cli) at (6.3,-4.4) {\textbf{CLI} \code{agentm2m workspace/auto \ldots}};
  \node[side=model] (rest) at (6.3,-5.5) {\textbf{REST API} (hosted service)};
  \node[side=raw] (hook) at (6.3,-2.6) {\textbf{validation hook} (runs the CLI)};
  \draw[->] (cli) -- (svc.east |- cli); \draw[->] (rest) -- (svc.east |- rest); \draw[->] (hook) -- (cli);
  \node[side=ext, draw=extbar, fill=white] (fw) at (-5.9,-1.3) {\textbf{hosted service}\\(if \code{AGENTM2M\_URL} is set)};
  \draw[->, dashed] (srv.west) -| node[lbl, pos=0.75, left] {forward} (fw.south);
\end{tikzpicture}
\caption{The layers of the plugin. Only \cc{} involves an LLM; everything
from the MCP server down is deterministic. The CLI, the validation hook and
the REST API of the hosted service reach the same service layer.}
\label{fig:arch}
\end{figure}

\subsection{The MCP layer is deliberately thin}

Every tool function in \file{mcp_server.py} is one or two lines: find the
workspace, call one method, return its dictionary (\Cref{lst:tooldef}). All
logic lives in the \emph{service layer}: \code{Workspace}
(\file{src/agentm2m/workspace.py}) for hand-written teams and
\code{AutoWorkspace} (\file{src/agentm2m/auto/workspace.py}) for \auto{}.
Three benefits follow, as in any layered design. \emph{Testability}: the
service layer is tested with plain Python calls, without an MCP server; most
of the plugin's tests do exactly that. \emph{Reuse}: the command-line
interface, the MCP server and the REST API of the hosted service all call
the same methods, so they cannot behave differently. \emph{Extensibility}:
the hosted service added a REST API without touching any business logic;
each endpoint calls the same method as the corresponding tool.

\subsection{Which project does a tool act on?}

A tool call carries no project path. The server finds the project directory
by a fixed rule: on the hosted service, from the authenticated user and the
requested project (\Cref{sec:hosted}); locally, from
\code{AGENTM2M\_PROJECT\_DIR}, else \code{CLAUDE\_PROJECT\_DIR} (the project
\cc{} was started in), else the current directory, which for a server
started by \cc{} is the project. One \code{Workspace} object is kept per
directory.

\subsection{What is stored where}

\begin{table}[!htb]
\centering\small
\caption{The workspace directory \file{.agentm2m/} of a project.}
\label{tab:files}
\rowcolors{2}{black!3}{white}
\begin{tabularx}{\textwidth}{@{}l Y l@{}}
\toprule
\rowcolor{white}File & Contents & Written by \\
\midrule
\file{team.yaml} & views and hand-offs (\Cref{lst:teamyaml}) & you \\
\file{rules/*.agentm2m} & rule modules (\Cref{lst:req2arch}) & you \\
\file{rules/helpers.py} & validators and helpers; \emph{executed} by the engine & you \\
\file{state/state.json} & view models, trace links, stamps, team evolutions & engine \\
\file{state/.lock} & the file lock & engine \\
\file{auto/task.json} & \auto{}: the lifted task & engine \\
\file{auto/team.json} & \auto{}: the admitted typed team & engine (via \mcptool{auto\_submit\_team}) \\
\file{auto/proposal.json} & \auto{}: the last rejected proposal & engine \\
\file{auto/state.json} & \auto{}: models and trace links of the running team & engine \\
\file{auto/history.json} & \auto{}: admissions, runs, attributions (last 200 events) & engine \\
\file{auto/faults.json} & \auto{}: the last fault reports & engine \\
\file{auto/compiled/}, \file{auto/runs/} & \auto{}: generated rule modules; work files of unattended runs & engine \\
\bottomrule
\end{tabularx}
\end{table}

Commit \file{team.yaml} and \file{rules/} to version control: you wrote
them. Leave \file{state/} and \file{auto/} to the engine (commit them only
if you want to share a run). The session-start hook repeats the warning
``never edit \file{state/} by hand'': a hand edit could give a value a stamp
it never earned.

\subsection{Several processes, one directory}

The same workspace can be touched by several processes: the MCP server, the
validation hook (which runs the CLI in a separate process), a second \cc{}
session, or you on the command line. Three mechanisms protect the files.

\begin{enumerate}
\item \textbf{Atomic writes.} State is written to \file{state.json.tmp} and
then renamed over \file{state.json} with \code{os.replace}. On POSIX file
systems (Linux, macOS) a rename is atomic: a reader sees the old file or the
new one, never half of one.
\item \textbf{A file lock.} Each load and each write of a hand-written
team's state takes an exclusive operating-system lock on
\file{state/.lock} (\code{fcntl.flock}), so no process reads a file while
another writes it. Inside one process, a re-entrant lock
(\code{threading.RLock}, a lock the same thread may take again) makes
concurrent tool calls run one at a time. The two locks answer different
questions: one is between processes, the other between threads.
\item \textbf{Reload on change.} Before each operation the workspace
computes a \emph{fingerprint}: modification time and size of
\file{team.yaml}, \file{state.json} and every rule file. If it differs from
the one remembered after the last load, another process changed something,
and the workspace reloads from disk.
\end{enumerate}

These mechanisms keep the \emph{files} consistent. They do not make a
whole operation atomic across processes, as this interleaving shows
(we ran it with two \code{Workspace} objects on one directory):

\begin{center}
\small
\begin{tabular}{@{}l l l@{}}
\toprule
Time & Process A & Process B \\
\midrule
1 & loads the state & \\
2 & & submits the signature of \code{S2}, saves \\
3 & submits the signature of \code{S1}: \code{accepted} & \\
4 & saves: notices the fingerprint changed, reloads, writes & \\
\midrule
result & \multicolumn{2}{l}{\code{S2} fresh; \code{S1} back to \emph{missing} although A reported \code{accepted}} \\
\bottomrule
\end{tabular}
\end{center}

The state stays consistent, and the lost value is simply offered again by
the next \mcptool{next\_bindings}; but an acceptance that was reported can
be lost. This is the classic \emph{lost update} problem of concurrent
programming. Within one server process it cannot happen (the re-entrant lock
serialises every operation); across processes it can. \auto{}'s files have
atomic writes and reload-on-change, but no file lock.

\begin{remark}[Errors leave no trace]
\mcptool{model\_edit} is atomic in the database sense: it takes a snapshot
of the state, applies all operations, and restores the snapshot if any of
them fails. Either every operation of an edit is applied, or none is.
\end{remark}

\begin{think}
In the interleaving above, why does process A's save not simply overwrite
B's change? And why is it acceptable that A's value is lost?
\end{think}

\begin{takeaway}
\begin{itemize}
\item Layers: \cc{} $\to$ MCP (stdio) $\to$ thin tool functions $\to$
service layer $\to$ engine $\to$ files. Only \cc{} uses an LLM.
\item The service layer is shared by MCP, CLI and REST, so all three behave
the same and can be tested without a server.
\item Atomic renames, a file lock and fingerprint reloads keep the files
consistent; overlapping operations of two processes can still lose a
reported acceptance, which is then re-offered.
\end{itemize}
\end{takeaway}

% ===========================================================================
\section{Host mode: \cc{} plays the agents, the engine decides}\label{sec:host}
% ===========================================================================

\subsection{Who calls the LLM?}

In the original library the engine called an LLM through a \emph{backend}: a
small class with one method, \code{generate(prompt) -> str}. The plugin
keeps those backends and adds one (\Cref{tab:backends}). The environment
variable \code{AGENTM2M\_LLM} chooses.

\begin{table}[!htb]
\centering\small
\caption{Backends: who produces the values of \code{@llm} bindings.}
\label{tab:backends}
\begin{tabularx}{\textwidth}{@{}l Y@{}}
\toprule
\code{AGENTM2M\_LLM} & Who fills the values \\
\midrule
\code{host} (default) & \cc{}, through \mcptool{next\_bindings} and \mcptool{submit\_binding}. No extra API key. \\
\code{anthropic}, \code{openai}, \code{ollama} & The engine calls that provider itself (an \emph{engine LLM}; needs an API key, or a local Ollama). Good for unattended runs. \\
\code{mock} & Canned values, for demos and tests. \\
\bottomrule
\end{tabularx}
\end{table}

\subsection{The trick: a backend that never answers}\label{sec:host-why}

How can the engine run its normal algorithm (match, create, bind, sample,
validate, stamp) without sampling? The host backend is marked
\emph{deferred}. When the engine reaches a value that needs (re)sampling, it
checks this mark, builds the complete prompt itself, and records a
\code{PendingBinding} instead of asking for a sample. (The host backend's
\code{generate} method throws an exception, \code{PendingSample}, in case
any code path ever calls it anyway.) The prompt is the binding's prompt
text, followed by ``Context (footprint only):'' and the footprint rendered
as text. \Cref{lst:pending} shows a real pending value as
\mcptool{next\_bindings} returns it.

\begin{lstlisting}[style=engine, caption={Observed: one pending value, as returned by \code{next\_bindings} (from \Cref{sec:walk1}, Step 8).}, label={lst:pending}]
{"handoff": "Arch2Sec", "rule": "Operation2Review",
 "target_key": "Operation2Review::sr::op=Operation#op_s1",
 "binding": "notes",
 "prompt": "Summarize any authentication or authorization concerns for this operation signature
   in at most three sentences.\n\nContext (footprint only):\ncreate_task(title: str) -> Task",
 "attempts": 0, "kind": "new", "footprint_version": "216262ec9373d8d9",
 "agent": "SecurityReviewer", "max_attempts": 3}
\end{lstlisting}

\begin{itemize}[itemsep=1pt]
\item \code{handoff}, \code{rule}, \code{binding}: where the value comes from.
\item \code{target\_key}: the value's place, \code{Rule::variable::match}; read
``rule \code{Operation2Review}, target variable \code{sr}, for the match
\code{op = Operation\#op\_s1}''.
\item \code{prompt}: everything the writer of the value may know.
\item \code{attempts}, \code{max\_attempts}: rejections so far, and $k$.
\item \code{kind}: \code{new} (never accepted) or \code{stale} (accepted
once, footprint changed since).
\item \code{footprint\_version}: the digest of the footprint, which becomes
the stamp if the value is accepted.
\item \code{agent}: the owner of the target view, who should write it.
\end{itemize}

The structural part of each hand-off runs normally: every target object,
reference and trace link exists before any value is asked for.

\plain{the engine writes the question and keeps the answer key; \cc{} only
fills in the blank.}

\begin{think}
Why can't the host backend simply call \cc{}'s LLM from inside
\code{generate()} and wait for the answer? Think about who calls whom when a
tool runs.
\end{think}

\subsection{One value, from prompt to stamp}

\Cref{fig:hostseq} follows one value through the system.

\begin{figure}[!htb]
\centering
\begin{tikzpicture}[font=\sffamily\scriptsize, >={Stealth[length=4pt]},
  hd/.style={draw=#1bar, fill=#1bg, rounded corners=2pt, minimum width=24mm, minimum height=10mm, align=center, font=\sffamily\small, anchor=north},
  msg/.style={->, thick}, ret/.style={->, dashed},
  lb/.style={fill=white, inner sep=1.5pt},
  note/.style={align=left, font=\sffamily\scriptsize, inner sep=3pt}]
  \node[hd=llm] (m) at (0,0) {\cc{}\\(main agent)};
  \node[hd=llm] (w) at (3.5,0) {subagent\\\code{binding-worker}};
  \node[hd=model] (s) at (7.1,0) {MCP server};
  \node[hd=engine] (e) at (10.8,0) {engine};
  \foreach \x in {m,s,e} \draw[black!35] (\x.south) -- ++(0,-10.6);
  \draw[black!35] (3.5,-2.9) -- (3.5,-11.6);
  \draw[black!35, dotted] (w.south) -- (3.5,-2.9);
  \draw[msg] (0,-1.4) -- node[lb, pos=0.22, above] {\mcptool{run}} (7.1,-1.4);
  \draw[msg] (7.1,-1.6) -- node[lb, above] {run to fixpoint} (10.8,-1.6);
  \draw[ret] (10.8,-2.1) -- node[lb, pos=0.5, above] {structure built: 5 pending, 2 blocked} (0,-2.1);
  \draw[msg] (0,-2.9) -- node[lb, above] {dispatch} (3.5,-2.9);
  \draw[enginebar!60!black, rounded corners=2pt] (2.9,-3.4) rectangle (12.3,-10.3);
  \node[anchor=north west, fill=white, draw=enginebar!60!black, font=\sffamily\scriptsize\bfseries] at (2.9,-3.4) {loop [values left]};
  \draw[msg] (3.5,-4.1) -- node[lb, above] {\mcptool{next\_bindings}} (7.1,-4.1);
  \draw[msg] (7.1,-4.3) -- (10.8,-4.3);
  \draw[ret] (10.8,-4.8) -- node[lb, above] {prompt + \code{footprint\_version} $d$} (3.5,-4.8);
  \node[note, fill=llmbg, draw=llmbar] at (3.5,-5.7) {writes value $v$\\from the prompt only};
  \draw[msg] (3.5,-6.6) -- node[lb, above] {\mcptool{submit\_binding}($v$, $d$)} (7.1,-6.6);
  \draw[msg] (7.1,-6.8) -- (10.8,-6.8);
  \node[note, fill=enginebg, draw=enginebar, text width=30mm] at (10.8,-8.0) {digest of current footprint $\neq d$: \code{stale}\\else run \code{@check}($v$):\\pass: write $v$, stamp $d$};
  \draw[ret] (10.8,-9.3) -- node[lb, above] {\code{accepted} | \code{rejected}(reason, retry prompt) | \code{escalated} | \code{stale}} (3.5,-9.3);
  \draw[ret] (3.5,-11.0) -- node[lb, above] {report} (0,-11.0);
\end{tikzpicture}
\caption{Host mode for one value. \cc{} and its subagent (purple) only
write text from a prompt; the server (blue) and the engine (green) decide
everything else. The subagent exists only from its dispatch on, and can
call only the two tools shown. Return messages pass through the server.}
\label{fig:hostseq}
\end{figure}

When a value is submitted, the engine does exactly what it would do with a
sample it drew itself, through the same function (\code{accept\_sample}): it
runs the \code{@check} validator, and only if it passes writes the value into
the target object and records the stamp. Validators, the budget $k$,
escalation, stamps and $\varphi$ are therefore identical with and without an
engine LLM. This is how the design meets R2.

\subsection{The five answers to a submission}

\begin{table}[!htb]
\centering\small
\caption{The statuses \mcptool{submit\_binding} returns, and what the caller should do.}
\label{tab:status-submit}
\begin{tabularx}{\textwidth}{@{}l Y Y@{}}
\toprule
Status & Meaning & Next step \\
\midrule
\code{accepted} & The validator passed; the value is written and stamped. & Next value. \\
\code{rejected} & The validator failed. The answer carries the \code{reason}, \code{attempts\_left} and a \code{retry\_prompt}. & Fix exactly what the reason says; resubmit. \\
\code{escalated} & The $k$-th rejection on the same footprint (or an earlier escalation). The value is not offered again for this footprint. & Leave it: it needs a change of the source or the team, or a human. \\
\code{stale} & The footprint changed since the prompt was issued, or the target no longer exists. & Fetch the value again; it comes back with a new prompt. \\
\code{already\_accepted} & A value is already accepted for this footprint. & Skip it. \\
\bottomrule
\end{tabularx}
\end{table}

The \code{retry\_prompt} is the original prompt plus ``Your previous answer
was rejected (\emph{reason}): \emph{your answer} (first 400 characters).
Answer again \ldots''. It adds only the model's \emph{own} previous answer and
the validator's reason, never source data outside the footprint. The code's
comments explain why it exists: re-sending the identical prompt at low
temperature tends to reproduce the same rejected answer.

\subsection{Why \code{footprint\_version} exists}

The development log (\file{plugin/DEVELOPMENT_PLAN.md}) records a bug found
during live testing: ``a code body written against the old signature was
stamped against the new one''. The following reconstruction shows how that
can happen. A Developer subagent receives the prompt for a code body, which
contains the operation's signature. While it writes, the Architect's new
signature is accepted. The Developer then submits code written against the
\emph{old} signature, and the engine stamps it with the digest of the
\emph{new} footprint, as if the code had been derived from data it never
saw. Change impact is then silently wrong. The fix: every prompt carries the
digest of its footprint, the submission must return it, and the engine
refuses a value whose digest no longer matches (\code{stale}). This is
\emph{optimistic concurrency control}~\cite{kung1981optimistic}, the idea
behind a version number on a database row or the \code{If-Match} header in
HTTP~\cite{rfc9110}: work without locking, and check at the end that nothing
changed underneath.

\plain{\code{footprint\_version} is a receipt that says which version of the
inputs you answered; an answer with an old receipt is sent back.}

\subsection{Blocked values}\label{sec:blocked}

Some footprints read values that an LLM writes upstream: the Developer's
code body reads the Architect's signature. With an engine LLM, hand-offs run
in order and the signature always exists in time. In host mode values arrive
whenever \cc{} submits them, so the engine could offer the Developer a
prompt with an empty signature. It does not: a value whose footprint reads
an upstream value that has \emph{never} been accepted is \emph{blocked} and
offered only once that value is accepted. A \emph{stale} upstream value does
not block: it exists, so it is shown, and if it is later replaced, the
\code{footprint\_version} check catches the outdated answer
(\Cref{sec:walk1}, Step 6). That is why the skills say ``call \mcptool{run}
again: accepted upstream values unblock downstream bindings''.

\subsection{The life of one value}

\Cref{fig:states} summarises the states of one value as
\mcptool{team\_status} reports them.

\begin{figure}[!htb]
\centering
\begin{tikzpicture}[font=\sffamily\small, >={Stealth[length=4pt]},
  st/.style={draw=modelbar, fill=modelbg, rounded corners=6pt, minimum width=22mm, minimum height=8mm, align=center},
  ok/.style={draw=enginebar, fill=enginebg, rounded corners=6pt, minimum width=22mm, minimum height=8mm, align=center},
  bad/.style={draw=badbar, fill=badbg, rounded corners=6pt, minimum width=22mm, minimum height=8mm, align=center},
  lbl/.style={font=\sffamily\scriptsize, align=center, fill=white, inner sep=1pt}]
  \node[st] (unc) at (0,0) {uncovered};
  \node[st] (mis) at (4.4,0) {missing};
  \node[ok] (fre) at (9.4,0) {fresh};
  \node[st] (sta) at (9.4,-3.0) {stale};
  \node[bad] (esc) at (4.4,-3.0) {escalated};
  \draw[->] (unc) -- node[lbl, above=2pt] {rule fires:\\target + trace link} (mis);
  \draw[->] (mis) -- node[lbl, above=2pt] {value passes\\\code{@check}: stamp} (fre);
  \draw[->] ([xshift=5mm]fre.south) -- node[lbl, right=1pt] {footprint\\changes} ([xshift=5mm]sta.north);
  \draw[->] ([xshift=-5mm]sta.north) -- node[lbl, left=1pt] {new value\\passes} ([xshift=-5mm]fre.south);
  \draw[->] ([xshift=5mm]mis.south) -- node[lbl, right=1pt] {$k$-th\\rejection} ([xshift=5mm]esc.north);
  \draw[->] ([xshift=-5mm]esc.north) -- node[lbl, left=1pt] {footprint changes\\(never accepted)} ([xshift=-5mm]mis.south);
  \draw[->] ([yshift=2mm]sta.west) -- node[lbl, above=2pt] {$k$-th rejection} ([yshift=2mm]esc.east);
  \draw[->] ([yshift=-2mm]esc.east) -- node[lbl, below=2pt] {footprint changes\\(accepted before)} ([yshift=-2mm]sta.west);
  \draw[->] (mis.north west) .. controls +(-0.3,0.9) and +(0.3,0.9) .. node[lbl, above] {rejected, attempts left} (mis.north east);
  \draw[->] (sta.north east) .. controls +(1.1,0.5) and +(1.1,-0.5) .. node[lbl, right] {rejected,\\attempts left} (sta.south east);
\end{tikzpicture}
\caption{States of one value for a hand-written team. ``Missing'' and
``stale'' values are \emph{pending}: offered by \mcptool{next\_bindings},
unless they are \emph{blocked} (\Cref{sec:blocked}). An escalated value gets a
fresh budget only when its footprint changes. For an \auto{} team, an
admitted repair also clears every escalation (\Cref{sec:walk2}).}
\label{fig:states}
\end{figure}

\subsection{What host mode keeps, and what it gives up}

\what{Kept.} Every structural decision (which objects exist, how they link,
which trace links are recorded) is made by the engine. Every value passes
the same validators, under the same budget $k$, and gets the same stamp.
``Done'' is still $\varphi$, computed by the engine.

\what{Given up.} The engine no longer controls \emph{which} model writes a
value, at what temperature, or what else that model has seen. When the
engine calls an LLM provider itself, that LLM receives the prompt and
nothing else. \cc{}, in contrast, has seen the whole conversation.
Footprint discipline therefore has to be enforced from the outside, by the
plugin (\Cref{sec:parts}).

\begin{think}
The Developer subagent fetched the prompt for the code body of
\code{op\_s2}. Before it submitted, the Analyst changed criterion
\code{S2.1}, and the Architect's new signature for \code{op\_s2} was
accepted. The Developer now submits its code with the old
\code{footprint\_version}. Which status does it get, and why is that right?
\end{think}

\begin{takeaway}
\begin{itemize}
\item In host mode the engine never samples: it builds the complete,
footprint-bounded prompt and waits for \cc{} to submit a value.
\item A submitted value goes through exactly the same acceptance path as an
engine sample: validator, budget $k$, escalation, stamp.
\item \code{footprint\_version} refuses answers to outdated prompts;
values whose upstream input was never accepted are blocked.
\item What the engine cannot control (which model, what else it saw) the
plugin must control with subagents.
\end{itemize}
\end{takeaway}

% ===========================================================================
\section{Skills, subagents and hooks}\label{sec:parts}
% ===========================================================================

The tools are the \emph{what}; skills, subagents and hooks are the
\emph{how}. This section shows how they cooperate, and how they close the gap
host mode leaves open (R3) and catch mistakes early (R5).

\subsection{The skills}

\begin{table}[!htb]
\centering\small
\caption{The ten skills of the plugin.}
\label{tab:skills}
\rowcolors{2}{black!3}{white}
\begin{tabularx}{\textwidth}{@{}L{4.0cm} Y@{}}
\toprule
\rowcolor{white}Skill & What it makes \cc{} do \\
\midrule
\skill{init} [template] & Check for an existing workspace (ask before replacing it), \mcptool{team\_init}, \mcptool{team\_validate}, report agents, views, hand-offs and the number of values waiting. For a custom team, hand the design to the \code{handoff-architect} subagent. \\
\skill{run} [agent] & \mcptool{run}; while values are pending, dispatch one \code{binding-worker} per owning agent (in parallel), then \mcptool{run} again; finally \mcptool{acceptance}, reporting every escalation with a suggested source change. \\
\skill{change} <change> & Turn the request into \mcptool{model\_edit} operations, preview them with \mcptool{impact} and show the obligations, apply, run, and show with \mcptool{trace\_query} that untouched artefacts were not redone. \\
\skill{evolve} <role> & Design a minimal view and hand-off for a new agent, \mcptool{team\_evolve}, run, report the obligations the new agent received. \\
\skill{status} [key] & Report the hand-written and/or \auto{} team, open items per agent, $\varphi$, and the trace closure of an element. \\
\skill{author-handoff} & A reference for writing views and rules, with rules of thumb (keep footprints minimal; validators return reasons). \\
\skill{auto-build} <file> & The \auto{} loop: task, propose (via \code{team-builder}), admit (at most 3 rejected rounds), run with \code{auto-binding-worker}s, $\varphi$, deliverable. \\
\skill{check} [file] & \W{1}--\W{6} on a typed team; group diagnostics by condition; propose the smallest JSON fix for each. \\
\skill{diagnose} & $\varphi$ failed: \mcptool{auto\_status}, \mcptool{auto\_attribute}, propose a minimal delta, submit it (checked), run again. \\
\code{agentm2m-concepts} & Not a command: background knowledge loaded automatically in projects with \file{.agentm2m/}. \\
\bottomrule
\end{tabularx}
\end{table}

Two habits run through \Cref{tab:skills}. The skills that change a team
\emph{preview or check first}: \mcptool{impact} before \mcptool{model\_edit}
in \skill{change}, the checker inside \mcptool{auto\_submit\_team} before
anything runs in \skill{auto-build}. And reports are tied to tool results;
\skill{change} even says ``Name only elements that appear in tool
results'', which guards against the LLM inventing a plausible but false
summary.

\subsection{The subagents and least privilege}

\Cref{tab:agents} lists the four subagents and their tool allowlists (tool
names abbreviated; in the files they are the full names of
\Cref{tab:names}).

\begin{table}[!htb]
\centering\footnotesize
\caption{Subagents and their tool allowlists.}
\label{tab:agents}
\rowcolors{2}{black!3}{white}
\begin{tabularx}{\textwidth}{@{}l Y Y@{}}
\toprule
\rowcolor{white}Subagent & Job & May call \\
\midrule
\code{binding-worker} & Fill pending values of one agent of a hand-written team. & \mcptool{next\_bindings}, \mcptool{submit\_binding}. Nothing else. \\
\code{auto-binding-worker} & The same for an \auto{} team. & \mcptool{auto\_next\_bindings}, \mcptool{auto\_submit\_binding}. \\
\code{team-builder} & Write a typed team for a task and iterate until admitted. & \code{Read}, \code{Grep}, \code{Glob}, \mcptool{auto\_status}, \mcptool{auto\_propose}, \mcptool{auto\_check}, \mcptool{auto\_submit\_team}. \\
\code{handoff-architect} & Design views and rules for a custom team. & \code{Read}, \code{Grep}, \code{Glob}, \code{Edit}, \code{Write}, \mcptool{team\_status}, \mcptool{team\_validate}, \mcptool{impact}, \mcptool{model\_show}, \mcptool{auto\_check}. \\
\bottomrule
\end{tabularx}
\end{table}

\begin{figure}[!htb]
\centering
\begin{tikzpicture}[font=\sffamily\scriptsize, >={Stealth[length=4pt]},
  box/.style={draw=llmbar, fill=llmbg, rounded corners=3pt, text width=55mm, align=left, inner sep=6pt}]
  \node[box] (main) at (0,0) {\textbf{\small Main agent's context}\\[2pt]
   the whole conversation\\files it has read\\all tools: \code{Read}, \code{Grep}, \code{Bash}, \code{Edit}, \ldots\\all 24 \tool{} tools};
  \node[box] (work) at (8.2,0) {\textbf{\small \code{binding-worker}'s context}\\[2pt]
   the dispatch sentence:\\``You are the Architect. Fill all pending\\agentm2m bindings for agent Architect.''\\its instructions (\file{agents/binding-worker.md})\\the prompts \mcptool{next\_bindings} returns\\tools: \mcptool{next\_bindings}, \mcptool{submit\_binding}};
  \draw[very thick, badbar] (4.1,1.35) -- (4.1,-1.35) node[below, align=center, text=badbar] {allowlist\\(least privilege)};
  \draw[->, thick] (main.east) -- node[above, fill=white, inner sep=1pt] {dispatch} (work.west);
\end{tikzpicture}
\caption{Context isolation. The worker starts with a fresh context; its
allowlist stops it from fetching anything but its prompts. Only the dispatch
sentence crosses the wall, and the skill fixes its wording.}
\label{fig:context}
\end{figure}

The two binding workers are the plugin's answer to R3
(\Cref{fig:context}). A \code{binding-worker} cannot read a file, search the
repository or run a command: those tools are not on its list, so \cc{} will
not let it call them. It starts with a fresh context, and its only source of
facts about the team is the prompt \mcptool{next\_bindings} returns, which is
the footprint. By the principle of least privilege, the worker cannot
\emph{fetch} anything beyond its prompts. What it is \emph{told} at dispatch
is still written by the main agent, which has seen everything; the skill
therefore fixes that message to one sentence. The guarantee is strong but not
absolute (\Cref{sec:limits}).

The two designers get the opposite rights: they \emph{should} read the
project's documentation to choose roles, and the architect may write rule
files. But neither can fill values, and the \code{team-builder} cannot run
the team: its instructions say ``Never call run tools'', and its allowlist
makes that true. Each subagent can do its job and little more.

\begin{remark}[Parallel workers, one engine]
\skill{run} dispatches one worker per agent \emph{in parallel}. Inside one
server process this is safe: each worker fetches and submits only its own
agent's values, the server runs one operation at a time, and a value
submitted against an outdated footprint is refused as \code{stale}. A worker
that finds only blocked values simply stops; \skill{run} calls \mcptool{run}
again and re-dispatches. Parallelism can cost a retry, but not a wrong stamp.
\end{remark}

\begin{think}\label{q:read}
A colleague proposes to add \code{Read} to \code{binding-worker}'s allowlist
``so it can follow the project's code style''. Explain, using the definition
of a stamp, what would break.
\end{think}

\subsection{The hooks}

\what{Session start.} \file{session-start.sh} checks whether the project's
\file{.agentm2m/} directory contains \file{team.yaml}, \file{auto/team.json}
or \file{auto/task.json}, and if so prints a short note that \cc{}
adds to its context. It deliberately does \emph{not} start the engine, so
session start stays instant and works offline (\Cref{lst:hookstart}).

\begin{lstlisting}[style=engine, caption={Observed: output of the \code{SessionStart} hook in a fresh \devteam{} project.}, label={lst:hookstart}]
agentm2m workspace detected at .agentm2m/ (team: devteam). Hand-offs between its agents are M2M
transformations run by the agentm2m MCP server: use /agentm2m:status to inspect, /agentm2m:run
to execute, /agentm2m:change for requirement changes. Never edit .agentm2m/state/ by hand.
\end{lstlisting}

\what{After every edit.} \cc{} passes the hook, on stdin, a JSON description
of the tool call that just happened (\Cref{lst:hookin}).

\begin{lstlisting}[style=wire, caption={Illustrative: the input of a \code{PostToolUse} hook (abridged; the fields used here are real).}, label={lst:hookin}]
{"hook_event_name": "PostToolUse", "tool_name": "Edit",
 "tool_input": {"file_path": "/home/you/play/.agentm2m/rules/Req2Arch.agentm2m", ...}, ...}
\end{lstlisting}

The script extracts \code{tool\_input.file\_path} with a three-line inline
Python program and then decides by the path:

\begin{itemize}[itemsep=1pt]
\item a file under \file{.agentm2m/state/}, or another file under
\file{.agentm2m/auto/}: do nothing (engine-owned);
\item \file{.agentm2m/auto/team.json} or any \file{*.typed-team.json}: run
the \auto{} checker (\code{agentm2m auto check});
\item a rule module, \file{team.yaml}, \file{helpers.py} or a view spec in
\file{.agentm2m/}: run \code{agentm2m workspace validate --brief};
\item anything else: do nothing.
\end{itemize}

If validation fails, the script writes the errors to stderr and exits with
code 2; \cc{} shows them to the LLM, which can fix the file at once.
\Cref{lst:hookerr} shows what it reports for a misspelt class name and for a
typed team whose Tester lost its tools.

\begin{lstlisting}[style=engine, caption={Observed: what the validation hook reports (we ran the commands the hook runs; file path shortened).}, label={lst:hookerr}]
# a rule module with the typo  to op : Arch!Operatoin (
agentm2m: workspace INVALID
  - Req2Arch: unknown class or missing root slot 'Operatoin' (declare it in team.yaml)

# a *.typed-team.json in which the Tester has no tools
AutoM2M checker rejected bad.typed-team.json:
W6  Story2TestSuite.code needs ['exec']; owner Tester has []
REJECTED: 1 violation(s)
Fix the diagnostics (or submit through auto_submit_team).
\end{lstlisting}

Two robustness choices are worth copying. The hook runs the engine through
\code{uvx} with the same pinned version as the MCP server. And if \code{uvx}
is missing (or \code{python3}, which it needs to read its input), the hook
exits silently with code 0: a missing tool must never block the user's
editing. When we ran the hook on a machine without \code{uvx} in its
\code{PATH}, this is exactly what happened; \Cref{lst:hookerr} was produced
by running the same commands directly.

\begin{takeaway}
\begin{itemize}
\item Skills that change a team preview or check first, and their reports
are tied to tool results.
\item Binding workers can call two tools only: they cannot fetch anything
beyond their prompts. Designers can read and write files but cannot fill
values or run the team.
\item Hooks announce a team at start (offline) and re-check every edited
team file, feeding errors back through exit code 2; they fail soft.
\end{itemize}
\end{takeaway}

% ===========================================================================
\section{Walkthrough 1: a hand-written team in the plugin}\label{sec:walk1}
% ===========================================================================

We now follow a complete session with \devteam{} (\Cref{sec:devteam}). Each
step names the skill that triggers it; the listings show the tool outputs.
To reproduce them without \cc{}, run
\file{docs/plugin-walkthroughs/walkthrough1_devteam.py}.

\begin{tcolorbox}[enhanced, breakable, colback=rawbg, colframe=rawbg, arc=0pt,
   boxrule=0pt, borderline west={2.2pt}{0pt}{rawbar}, left=8pt, right=8pt, top=5pt, bottom=5pt]
\textsf{\bfseries\color{rawbar}What you see in the terminal (illustrative).}\enspace
You type \skill{init devteam}. \cc{} loads the \code{init} skill and asks
your permission before calling the plugin's MCP tools (you can allow them
once or for the session). It calls \mcptool{team\_status},
\mcptool{team\_init} and \mcptool{team\_validate}, and summarises: four
agents, their views, three hand-offs, values waiting. You type \skill{run}.
\cc{} calls \mcptool{run}, then starts one \code{binding-worker} subagent per
agent that has work; you see each subagent's progress as a collapsed line,
and each returns a short report (``accepted 2, escalations: none''). \cc{}
runs again until nothing is pending and ends with $\varphi$ and any
escalations. Below we show what happens \emph{inside} these steps.
\end{tcolorbox}

\subsection{Step 1 (\skill{init}): create, validate, look}

\mcptool{team\_init} copies the template into \file{.agentm2m/};
\mcptool{team\_validate} parses every rule module and counts matches
(\Cref{lst:w1validate}).

\begin{lstlisting}[style=engine, caption={Observed: \code{team\_validate} on the fresh workspace (abridged).}, label={lst:w1validate}]
{"ok": true, "errors": [],
 "handoffs": [
  {"name": "Req2Arch",  "rules": [{"rule": "Epic2Component",     "matches": 1, "stochastic_bindings": 0},
                                  {"rule": "Story2Operation",    "matches": 2, "stochastic_bindings": 1}]},
  {"name": "Arch2Code", "rules": [{"rule": "Operation2CodeEdit", "matches": 0, "stochastic_bindings": 1}]},
  {"name": "Req2Test",  "rules": [{"rule": "Criterion2TestCase", "matches": 3, "stochastic_bindings": 1}]}]}
\end{lstlisting}

\code{Story2Operation} matches 2 stories (the guard excludes the draft);
\code{Operation2CodeEdit} matches nothing yet, because no operation exists
yet; \code{Criterion2TestCase} matches all 3 criteria. \mcptool{team\_status}
gives the overview a student should look at first (\Cref{lst:w1status}).

\begin{lstlisting}[style=engine, caption={Observed: \code{team\_status} before the first run (abridged).}, label={lst:w1status}]
{"team": "devteam", "backend": "host",
 "agents": {"Analyst": ["Req"], "Architect": ["Arch"], "Developer": ["Code"], "Tester": ["Test"]},
 "views": {"Req": {"owner": "Analyst", "elements": 7}, "Arch": {"owner": "Architect", "elements": 0}, ...},
 "handoffs": [{"name": "Req2Arch", "trace_links": 0, "bindings": {"uncovered": 2}}, ...],
 "bindings": {"uncovered": 5}, "phi": false, "open_total": 5}
\end{lstlisting}

The five values are \code{uncovered} (\Cref{fig:states}): the rules would
match, but no target object exists yet.

\subsection{Step 2 (\skill{run}): build the structure}

\begin{lstlisting}[style=engine, caption={Observed: the first \code{run} (abridged).}, label={lst:w1run}]
{"backend": "host", "passes": 2,
 "handoffs": {"Req2Arch":  {"created": 3, "pending": 2, "blocked": 0},
              "Arch2Code": {"created": 2, "pending": 0, "blocked": 2},
              "Req2Test":  {"created": 3, "pending": 3, "blocked": 0}},
 "pending": 5, "blocked": 2, "phi": false,
 "pending_preview": [
   {"target_key": "Story2Operation::op::s=UserStory#S1",       "binding": "signature", "agent": "Architect"},
   {"target_key": "Story2Operation::op::s=UserStory#S2",       "binding": "signature", "agent": "Architect"},
   {"target_key": "Criterion2TestCase::tc::c=Criterion#S1.1",  "binding": "oracle",    "agent": "Tester"}, ...]}
\end{lstlisting}

Without any LLM call the engine created 8 objects (1 component, 2
operations, 2 code edits, 3 test cases) with their references and trace
links. Within one pass the hand-offs run in the order of \file{team.yaml},
so \code{Arch2Code} already saw the operations that \code{Req2Arch} had just
created; the second pass changed nothing, which tells the engine it has
reached the fixpoint. Five
values are pending. The two code bodies are \emph{blocked}: their footprint
contains a signature that has never been accepted.

\subsection{Step 3 (inside a \code{binding-worker}): the prompts}

\mcptool{next\_bindings} returns the ready values with their prompts, a
\code{remaining} count, a \code{blocked\_on\_upstream} count, and an
\code{instructions} field that repeats the rule ``answer ONLY from its
prompt''. \Cref{lst:w1prompts} renders two of the five values.

\begin{lstlisting}[style=llm, caption={Observed: two of the five values returned by \code{next\_bindings} (rendered).}, label={lst:w1prompts}]
target_key: Story2Operation::op::s=UserStory#S1   binding: signature   agent: Architect
footprint_version: d550b926cae2fdb8               attempts: 0   kind: new   max_attempts: 3
prompt:
  Derive an API signature for this user story. Answer with exactly one line of the form
  name(param: Type, ...) -> ReturnType and nothing else.

  Context (footprint only):
  - Criterion({'id': 'S1.1', 'text': 'a title is required'})

target_key: Criterion2TestCase::tc::c=Criterion#S2.1   binding: oracle   agent: Tester
footprint_version: 123ccdc996aff09e
prompt:
  Write a test oracle for this acceptance criterion: Python code defining a function
  test_oracle() that calls implementation() and asserts the criterion holds. Answer with
  Python code only.

  Context (footprint only):
  completing a done task is rejected
\end{lstlisting}

The prompt is \emph{everything} the writer of the value may know. Notice what
is missing: the signature's footprint is \code{s.criteria}, so the Architect
does not even see the story's title. That is a choice of this template, and
it is visible to everyone who reads the rule. Whatever is not in the
footprint cannot influence the value, and a later change to it will not make
the value stale.

\subsection{Step 4: a wrong answer}

We first submit a sloppy signature, the way an LLM might
(\Cref{lst:w1reject}).

\begin{lstlisting}[style=engine, caption={Observed: \code{submit\_binding} with the value ``Here is the signature: create\_task''.}, label={lst:w1reject}]
{"target_key": "Story2Operation::op::s=UserStory#S1", "binding": "signature", "handoff": "Req2Arch",
 "status": "rejected",
 "reason": "expected exactly one line: name(param: Type, ...) -> ReturnType",
 "attempts": 1, "attempts_left": 2,
 "retry_prompt": "Derive an API signature for this user story. ... Context (footprint only):
   - Criterion({'id': 'S1.1', 'text': 'a title is required'})
   Your previous answer was rejected (expected exactly one line: name(param: Type, ...) ->
   ReturnType):
   Here is the signature: create_task
   Answer again, following the requested output format exactly."}
\end{lstlisting}

The reason comes from \code{parses()} in \Cref{lst:helpers}. Because the
validator explains itself, the retry can be precise.

\subsection{Step 5: fill everything; $\varphi$ holds}

We then submitted acceptable values, in this order:
\begin{enumerate}[itemsep=1pt]
\item the signatures \code{create\_task(title: str) -> Task} and
\code{mark\_done(task\_id: int) -> Task}: accepted;
\item three oracles, for example for \code{S1.1}:
\begin{lstlisting}[language=Python, style=llm, caption={Illustrative: the oracle we submitted for \code{S1.1} (accepted).}, label={lst:w1oracle}]
def test_oracle():
    try:
        implementation('')
    except ValueError:
        return
    raise AssertionError('empty title accepted')
\end{lstlisting}
all accepted by \code{failsOnStub()};
\item a new \mcptool{next\_bindings}: the signatures being accepted, the two
code bodies are now unblocked and offered, with the signature and the copied
criteria in their prompts;
\item two code bodies, both the deliberately useless
\code{def f(x): return x}: both accepted.
\end{enumerate}
The next \mcptool{run} reported \code{"pending": 0, "phi": true}; all seven
values are \code{fresh}.

\begin{remark}[A weakness made visible]
Both code bodies were accepted because the template checks them only with
\code{compiles()}, a \emph{form} validator (red in \Cref{fig:devteam}).
$\varphi$ holds, yet nothing checked that the code does what \code{S2.1}
asks. If this team were written as a typed team, the \auto{} checker would
reject it under \W{4} (``view Code reaches no behavioural check''); the
companion reports the same finding for the hand-written pilot team, together
with \code{Arch}~\cite{autom2mteaching}. $\varphi$ is only as meaningful as
the validators it counts.
\end{remark}

\begin{think}
Which validator let \code{def f(x): return x} through, and which admission
condition would reject this team design?
\end{think}

\subsection{Step 6 (\skill{change}): a change, previewed before it costs anything}

The Analyst tightens \code{S2.1} to ``completing a done task returns HTTP
409''. The \skill{change} skill first looks at the source
(\mcptool{model\_show}), turns the request into an edit operation, and
previews it with \mcptool{impact}, which applies the edit to a scratch copy,
runs the structural phase and calls no LLM (\Cref{lst:w1impact}).

\begin{lstlisting}[style=engine, caption={Observed: \code{impact} for the edit \code{\{"op": "set", "key": "Criterion\#S2.1", "values": \{"text": ...\}\}} of view \code{Req} as the Analyst.}, label={lst:w1impact}]
{"obligations": [
   {"handoff": "Req2Arch",  "target_key": "Story2Operation::op::s=UserStory#S2",
    "binding": "signature", "agent": "Architect"},
   {"handoff": "Arch2Code", "target_key": "Operation2CodeEdit::ce::op=Operation#op_s2",
    "binding": "body", "agent": "Developer"},
   {"handoff": "Req2Test",  "target_key": "Criterion2TestCase::tc::c=Criterion#S2.1",
    "binding": "oracle", "agent": "Tester"}],
 "new_bindings": [], "created": [], "deleted": [],
 "llm_calls_upper_bound": 9,
 "note": "direct obligations only: a re-sampled value that changes may oblige downstream bindings on the next run"}
\end{lstlisting}

Exactly three values must be redone: \code{S2}'s signature (its footprint is
\code{S2}'s criteria), \code{S2}'s code body (its footprint contains the
criteria copied onto the operation) and \code{S2.1}'s oracle. Nothing of
\code{S1} is touched. The upper bound of 9 LLM calls is 3 obligations times
$k=3$. The note matters: if the Architect's new signature differs from the
old one, the code body's footprint changes \emph{again}, one hop further; if
it is identical, propagation stops. (The plugin's tests check both cases.)
\Cref{fig:fptrace} shows why the engine finds exactly these three.

\begin{figure}[!htb]
\centering
\begin{tikzpicture}[font=\sffamily\scriptsize, >={Stealth[length=4pt]},
  o/.style={draw=modelbar, fill=modelbg, rounded corners=2pt, align=center, inner sep=4pt, minimum height=9mm},
  hot/.style={draw=rawbar, fill=rawbg, rounded corners=2pt, align=center, inner sep=4pt, minimum height=9mm, very thick},
  lb/.style={fill=white, inner sep=1.5pt, align=center}]
  \node[o] (s2) at (0,0) {\textbf{UserStory\#S2}};
  \node[hot] (c) at (0,-3.0) {\textbf{Criterion\#S2.1}\\text changed};
  \node[o] (op) at (5.2,0) {\textbf{Operation\#op\_s2}\\signature $\star$\\criteria (copy)};
  \node[o] (ce) at (10.4,0) {\textbf{CodeEdit (op\_s2)}\\body $\star$};
  \node[o] (tc) at (5.2,-3.0) {\textbf{TestCase\#S2.1}\\oracle $\star$};
  \draw[thick] (s2) -- node[lb, left] {contains} (c);
  \draw[->, very thick, enginebar] (s2) -- node[lb, above] {trace (\code{Req2Arch})} (op);
  \draw[->, very thick, enginebar] (op) -- node[lb, above] {trace (\code{Arch2Code})} (ce);
  \draw[->, very thick, enginebar] ([yshift=2mm]c.east) -- node[lb, above] {trace (\code{Req2Test})} ([yshift=2mm]tc.west);
  \draw[->, dashed, thick, llmbar] ([yshift=-2mm]c.east) -- node[lb, below] {footprint of oracle (\code{c.text})} ([yshift=-2mm]tc.west);
  \draw[->, dashed, thick, llmbar] (c.north east) -- node[lb, pos=0.45] {footprint of signature\\(\code{s.criteria})} (op.south west);
  \draw[->, dashed, thick, llmbar] (op.south east) to[bend right=30] node[lb, below=2pt] {footprint of body (signature + criteria)} (ce.south);
\end{tikzpicture}
\caption{Trace links versus footprints for the change of \code{S2.1}. Solid
green arrows are trace links (``created from''): \mcptool{trace\_query}
follows them. Dashed purple arrows are footprint reads (``written from''):
\mcptool{impact} compares their stamps. The three starred values are the
obligations.}
\label{fig:fptrace}
\end{figure}

Before applying the edit, two mistakes the engine refuses
(\Cref{lst:w1refuse}), and the reason for the second (\Cref{lst:w1show}).

\begin{lstlisting}[style=engine, caption={Observed: \code{model\_edit} refusals.}, label={lst:w1refuse}]
model_edit(view="Req", ..., as_agent="Architect")
  write rights (omega): agent Architect may not write view Req; its owner is Analyst
model_edit(view="Arch", ops=[set Operation#op_s2 name], as_agent="Architect")
  cannot modify Operation#op_s2: it is engine-owned (created by hand-off rule Story2Operation);
  change its source view instead
\end{lstlisting}

\begin{lstlisting}[style=engine, caption={Observed: \code{model\_show(view="Arch", key="Operation\#op\_s2")}.}, label={lst:w1show}]
{"view": "Arch", "element": {"type": "Operation", "key": "Operation#op_s2",
  "engine_owned": "Story2Operation::op::s=UserStory#S2",
  "name": "op_s2", "signature": "mark_done(task_id: int) -> Task",
  "criteria": "S2.1: completing a done task is rejected", "component": "Arch:Component#E1"}}
\end{lstlisting}

The Architect owns view \cls{Arch}, but \code{op\_s2} was \emph{created by a
rule}: it is \emph{engine-owned}, as \mcptool{model\_show} says. If an agent
could edit it, the next run would either overwrite the edit or contradict
the trace. To change it, change its source.

Applied as the Analyst, the edit reports \code{"open\_obligations": 3}, and
\mcptool{next\_bindings} offers the three values with \code{"kind":
"stale"} and new prompts. The Developer's prompt shows the \emph{old}
signature next to the \emph{new} criterion (\Cref{lst:w1stale}). It is not
blocked, because the old signature was once accepted (\Cref{sec:blocked}). If
the Developer answered now and the Architect's new signature then differed,
the Developer's answer would be refused as \code{stale} or become an
obligation again; either way no wrong stamp results.

\begin{lstlisting}[style=llm, caption={Observed: the Developer's stale prompt after the change (rendered).}, label={lst:w1stale}]
target_key: Operation2CodeEdit::ce::op=Operation#op_s2   binding: body   kind: stale
prompt:
  Implement this operation as one self-contained Python function: first the signature, then
  the acceptance criteria it must satisfy. Answer with Python code only.

  Context (footprint only):
  - mark_done(task_id: int) -> Task
  - S2.1: completing a done task returns HTTP 409
\end{lstlisting}

\subsection{Step 7 (\skill{status}): what came from what}

\mcptool{trace\_query} answers ``what was derived from this element?''. For
\code{Criterion\#S2.1} it returns one element; for the story, two
(\Cref{lst:w1trace}).

\begin{lstlisting}[style=engine, caption={Observed: \code{trace\_query} for the criterion and for its story (abridged).}, label={lst:w1trace}]
trace_query("Criterion#S2.1")
  downstream: [{"from": "Criterion#S2.1", "handoff": "Req2Test", "target": "Test:TestCase#S2.1"}]
trace_query("UserStory#S2")
  downstream: [{"from": "UserStory#S2",    "handoff": "Req2Arch",  "target": "Arch:Operation#op_s2"},
               {"from": "Operation#op_s2", "handoff": "Arch2Code", "target": "Code:CodeEdit#op_s2"}]
\end{lstlisting}

\begin{think}
\mcptool{impact} found three obligations for a change of
\code{Criterion\#S2.1}, but \mcptool{trace\_query} on the criterion lists only
one derived element. Using \Cref{fig:fptrace}, explain why both answers are
right.
\end{think}

\subsection{Step 8 (\skill{evolve}): grow the team while it runs}

\skill{evolve} adds a Security Reviewer with the ready-made view \cls{Sec}
and hand-off \code{Arch2Sec} from the template's \file{rules/extra/}. This is
a \emph{higher-order transformation}~\cite{tisi2009hot}: a change of the
team itself, while it runs. \mcptool{team\_evolve} reported
\code{"existing\_matches": 2}; the next \mcptool{run} created two
\cls{SecurityReview} objects and offered four new values (\code{notes} and
\code{risk} for each operation; \Cref{lst:pending} is one of them) next to
the three stale ones from Step 6. The new agent received work for everything
that already existed, without glue code. The evolution is stored in
\file{state.json}, so it survives restarts, and from now on
\mcptool{trace\_query} on \code{UserStory\#S2} also lists the review of
\code{op\_s2}.

\subsection{Step 9: back to $\varphi$}

We filled the seven pending values: the same signature for \code{op\_s2}
(so the code body did not become stale a second time), a new oracle and code body for the 409
behaviour, and two notes and two risk ratings. All were accepted; the final
\mcptool{run} reported \code{"phi": true} and \mcptool{team\_status}
\code{"bindings": \{"fresh": 11\}} with one evolution (\code{SecurityReviewer}).

\begin{takeaway}
\begin{itemize}
\item One \mcptool{run} builds all structure for free; values are filled one
prompt at a time, and blocked values wait for their inputs.
\item Validators that return reasons turn a rejection into a precise retry.
\item \mcptool{impact} previews a change exactly, before any token is spent:
it follows footprints, while \mcptool{trace\_query} follows trace links.
\item Write rights and engine ownership stop edits in the wrong place.
\item $\varphi$ is only as strong as the validators: a compile-only check
lets wrong code through.
\end{itemize}
\end{takeaway}

% ===========================================================================
\section{Walkthrough 2: \auto{} in the plugin}\label{sec:walk2}
% ===========================================================================

In the first walkthrough a person had designed the team. Now \cc{} designs
it, as the \emph{builder}, and the checker decides whether it may run. In
\cc{} you type \skill{auto-build wallet.py}. \Cref{fig:autoloop} shows the
loop the skill drives. The task is illustrative; all engine outputs are
observed (script \file{walkthrough2_autom2m.py}).

\begin{figure}[!htb]
\centering
\begin{tikzpicture}[font=\sffamily\scriptsize, >={Stealth[length=4pt]},
  t/.style={draw=modelbar, fill=modelbg, rounded corners=2pt, align=center, inner sep=3pt, minimum height=9mm, text width=24mm},
  l/.style={draw=llmbar, fill=llmbg, rounded corners=2pt, align=center, inner sep=3pt, minimum height=9mm, text width=24mm},
  e/.style={draw=enginebar, fill=enginebg, rounded corners=2pt, align=center, inner sep=3pt, minimum height=9mm, text width=24mm},
  d/.style={draw=enginebar, fill=enginebg, diamond, aspect=2, inner sep=1pt, align=center},
  lb/.style={fill=white, inner sep=1.5pt, align=center}]
  \node[t] (task) at (0,0) {\mcptool{auto\_task\_set}\\lift the skeleton};
  \node[l] (prop) at (3.4,0) {\mcptool{auto\_propose}\\\code{team-builder} writes\\a typed team};
  \node[e] (sub) at (6.8,0) {\mcptool{auto\_submit\_team}\\checker \W{1}--\W{6}};
  \node[l] (fill) at (10.6,0) {\mcptool{auto\_run} + fill\\values (\code{auto-}\\\code{binding-worker}s)};
  \node[d] (phi) at (10.6,-2.2) {$\varphi$?};
  \node[e] (del) at (13.6,-2.2) {\mcptool{auto\_deliverable}\\the code};
  \node[e] (att) at (6.8,-2.2) {\mcptool{auto\_attribute}\\locate the fault};
  \node[l] (delta) at (3.4,-2.2) {\mcptool{auto\_propose}\\(\code{delta})};
  \draw[->] (task) -- (prop); \draw[->] (prop) -- (sub);
  \draw[->] (sub) -- node[lb, above] {admitted} (fill);
  \draw[->] (sub.north) to[bend right=35] node[lb, above] {rejected: diagnostics (\code{revise}, $\le 3$ rounds)} (prop.north);
  \draw[->] (fill) -- (phi);
  \draw[->] (phi) -- node[lb, above] {yes} (del);
  \draw[->] (phi) -- node[lb, above] {no} (att);
  \draw[->] (att) -- (delta);
  \draw[->] (delta) -- node[lb, right, pos=0.4] {checked delta} (sub.south west);
\end{tikzpicture}
\caption{The \auto{} loop as \skill{auto-build} and \skill{diagnose} drive
it. Purple steps are written by \cc{} or its subagents; green steps are
decided by the checker or the engine.}
\label{fig:autoloop}
\end{figure}

\subsection{The task}

\auto{} currently handles Python tasks: a class whose methods are
\emph{stubs}, or one stub function. A stub has a docstring and a body of
\code{pass}, \code{...} or \code{raise NotImplementedError}. Docstring lines
that start with \code{>>>} are \emph{doctest examples}: a call and, on the
next line, its expected output.

\begin{lstlisting}[language=Python, style=raw, caption={Illustrative: the task \file{wallet.py}.}, label={lst:wallet}]
class Wallet:
    """A wallet that holds money in cents."""

    def __init__(self):
        self.cents = 0

    def deposit(self, cents):
        """
        Add money to the wallet.
        :param cents: int, must be positive
        :return: int, the new balance
        >>> w = Wallet()
        >>> w.deposit(250)
        250
        """
        pass

    def withdraw(self, cents):
        """
        Remove money; raise ValueError if the balance would go below zero.
        :return: int, the new balance
        >>> w = Wallet(); _ = w.deposit(300)
        >>> w.withdraw(120)
        180
        """
        pass
\end{lstlisting}

\mcptool{auto\_task\_set} parses the source and \emph{lifts} it into the
fixed \cls{Goal} view: one \cls{Task} object (name, description, skeleton)
and one \cls{Method} object per stub (name, signature, docstring, examples).
It found \code{deposit} and \code{withdraw}, both with examples;
\code{\_\_init\_\_} is already implemented and stays part of the frame into
which the methods will later be assembled.

\subsection{Step A: the builder prompt}

\mcptool{auto\_propose} chose mode \code{propose} (there is no team yet) and
returned a prompt of 8{,}587 characters. It contains, in order: the JSON
format of a typed team; the fixed \cls{Goal} view; the \emph{validator
library}, the only validators a team may use (\Cref{lst:vlib}); the six
admission conditions in plain words; a worked example (a three-agent team
for a parking meter); and the task.

\begin{lstlisting}[style=llm, caption={Observed: the validator library as the builder prompt lists it.}, label={lst:vlib}]
- nonempty() [form]: value is non-empty text
- json() [form]: value is a JSON document
- compiles() [form]: value contains Python code that parses
- defines(method: a Goal object) [form]: value defines a function with the method's name
- runs() [behaviour] needs tools ['exec']: value is Python code that executes without raising
- examples_run(method: a Goal object) [behaviour] needs tools ['exec']: the method, spliced
  into the class, runs the method's public docstring examples without raising
- examples_match(method: a Goal object) [behaviour] needs tools ['exec']: as examples_run,
  and every printed result equals the documented one
- passes_tests(method: a Goal object, tests: string) [behaviour] needs tools ['exec']: the
  method, spliced into the class, passes the given test code
- test_valid(method: a Goal object) [behaviour] needs tools ['exec']: value is test code that
  calls the method and fails on a stub implementation
\end{lstlisting}

Each validator declares its \emph{strength} (\code{form} or
\code{behaviour}) and the tools its owner needs. These declarations make
\W{4} (``every goal reaches a behavioural check'') and \W{6} (``the owner has
the tools'') decidable from the JSON alone. Only behaviour validators execute
code.

\subsection{Steps B and C: a rejected proposal, then an admitted one}

Our first proposal had the shape of the worked example: a Designer who plans
each method, a Tester who writes tests, a Developer who implements. We made
four deliberate mistakes (\Cref{lst:w2bad}).

\begin{lstlisting}[style=model, caption={Illustrative: the four faulty fragments of the first proposal.}, label={lst:w2bad}]
"agents": [..., {"name": "Tester", "role": "...", "tools": []}, ...]          (no exec)
"writes": {"Designer": ["Design"], "Tester": ["Test"], "Developer": ["Code", "Test"]}
... "footprint": ["d.method.task.skeleton", ..., "t.code", "d.method.returnType"] ...
"done": ["cover(G)", "valid", "fresh", "noObl", "Tester.says('ALL TESTS PASS')"]
\end{lstlisting}

\begin{lstlisting}[style=engine, caption={Observed: \code{auto\_submit\_team} on the first proposal.}, label={lst:w2reject}]
W1  Method2Impl.code: footprint path 'd.method.returnType': Goal!Method has no feature
    'returnType' (its features are ['docstring', 'examples', 'name', 'signature', 'task'])
W2  view Test written by ['Developer', 'Tester']; needs exactly one writer
W5  done-clause Tester.says('ALL TESTS PASS') is not engine-checkable
W6  Method2Test.code needs ['exec']; owner Tester has []
REJECTED: 4 violation(s)
next: fix every diagnostic and resubmit (auto_propose gives a revise prompt)
\end{lstlisting}

Each diagnostic names the condition, the exact element and what is wrong;
the \W{1} message even lists the features that \emph{do} exist. The rejected
proposal is saved in \file{auto/proposal.json}, and the next
\mcptool{auto\_propose} switches to mode \code{revise}: the same format and
rules, the rejected JSON, and ``Checker diagnostics (each must be fixed)''.
\skill{auto-build} allows three rejected rounds before it stops and shows the
diagnostics to you.

Our corrected team was admitted in round 2, with \code{"mode": "admitted"},
agents \code{Designer}, \code{Tester}, \code{Developer} and hand-offs
\code{Goal2Design}, \code{Goal2Test}, \code{Build2Code}. In fairness: apart
from its name, it is exactly the worked example from the builder prompt.
That is not a design achievement; it shows how strongly builders copy the
example they are given, which the project's preliminary experiments also
report (100 of 103 final teams had the example's shape; \file{README.md}). \Cref{app:team} prints the whole team;
\Cref{lst:w2team} shows its most interesting rule.

\begin{lstlisting}[style=model, caption={Code excerpt: the rule that writes the deliverable (from \Cref{app:team}).}, label={lst:w2team}]
{"name": "Build2Code", "sources": ["Goal", "Design", "Test"], "target": "Code",
 "rules": [{"name": "Method2Impl",
   "from":  [{"var": "d", "type": "Design!MethodDesign"}, {"var": "t", "type": "Test!MethodTest"}],
   "guard": "t.method = d.method",
   "to":    {"var": "i", "type": "Code!MethodImpl"},
   "bind":  {"name": "d.method.name", "method": "d.method"},
   "llm":   [{"feature": "code", "prompt": "Implement this method of the class.",
              "footprint": ["d.method.task.skeleton", "d.method.signature", "d.method.docstring",
                            "d.plan", "t.code"],
              "validator": [{"id": "examples_run", "args": {"method": "d.method"}},
                            {"id": "passes_tests", "args": {"method": "d.method", "tests": "t.code"}}]}]}]}
\end{lstlisting}

Read it like a hand-off rule: for every pair of a design and a test of the
same method, create one implementation; its \code{name} and \code{method} are
copied by the engine; its \code{code} is written by an LLM that sees the
skeleton, the signature, the docstring, the plan and the tests, and is
accepted only if the documented examples run \emph{and} the Tester's tests
pass.

\subsection{Step D: run and fill}

The first \mcptool{auto\_run} created the structure and reported 4 values
pending (two plans, two test suites) and 2 \emph{blocked} (the two
implementations, whose footprints contain a plan and a test that have never
been accepted). $\varphi$ was false with \code{cover(G)} and \code{fresh}
open, which simply means ``nothing is done yet''.

\auto{} prompts are richer than hand-written ones: the agent's role, then
the task, then a precise output format, then the footprint, each part
labelled by its path (\Cref{lst:w2prompt}).

\begin{lstlisting}[style=llm, caption={Observed: the Developer's prompt for \code{withdraw} (skeleton and tests abridged).}, label={lst:w2prompt}]
You are the Developer. You implement one method following its design so that its tests and
examples pass.

Task: Implement this method of the class.

Answer with the complete Python definition of the method named in the context (signature line
and body) in ONE ```python code block. Do not repeat the class.

Context (footprint only):
- [d.method.task.skeleton]
class Wallet: ...
- [d.method.signature]
def withdraw(self, cents):
- [d.method.docstring]
Remove money; raise ValueError if the balance would go below zero. ...
- [d.plan]
- if cents > self.cents raise ValueError
- subtract and return the new balance
- [t.code]
```python
import unittest
class TestWithdraw(unittest.TestCase):
    def test_subtracts(self): ...
    def test_overdraw_rejected(self):
        w = Wallet(); w.deposit(100)
        with self.assertRaises(ValueError):
            w.withdraw(1000)
```
\end{lstlisting}

We submitted the plans (validator \code{nonempty}) and the test suites
(validator \code{test\_valid}: the tests must call the method and \emph{fail}
on a stub). Tests must be \code{unittest.TestCase} classes (as the Tester's
prompt asks) or functions named \code{test\_*}; in an earlier attempt we
wrote plain \code{assert} lines at module level, which run while the file
loads, against the stub, and were rejected (``test code does not load'').
All four were accepted. The next \mcptool{auto\_run} unblocked the two
implementations. The \code{deposit} implementation was accepted. For
\code{withdraw} we submitted, three times, code that forgets the overdraft
check:

\begin{lstlisting}[language=Python, style=llm, caption={Illustrative: the wrong \code{withdraw} we submitted.}, label={lst:w2wrong}]
def withdraw(self, cents):
    self.cents -= cents
    return self.cents
\end{lstlisting}

\begin{lstlisting}[style=engine, caption={Observed: a behaviour validator rejects it (first of three attempts; the third returned \code{escalated}).}, label={lst:w2fail}]
{"target_key": "Method2Impl::i::d=MethodDesign#...#withdraw|t=MethodTest#...#withdraw",
 "binding": "code", "handoff": "Build2Code", "status": "rejected",
 "reason": "failing tests: AssertionError: ValueError not raised",
 "attempts": 1, "attempts_left": 2, "retry_prompt": "You are the Developer. ..."}
\end{lstlisting}

Compare this with Step 5 of the first walkthrough: there, a compile-only
validator let wrong code through. Here the code is \emph{executed} against
the documented examples and the Tester's tests, in the sandbox
(\Cref{sec:security}), and the reason names the failing assertion.

\subsection{Step E (\skill{diagnose}): $\varphi$ fails; locate the fault}

With nothing pending, \mcptool{auto\_run} reported $\varphi$ false and
\code{"deliverables\_accepted": \{"deposit": true, "withdraw": false\}}.
\mcptool{auto\_attribute} turned each failed clause into a located fault
(\Cref{lst:w2attr}).

\begin{lstlisting}[style=engine, caption={Observed: \code{auto\_attribute} in host mode (abridged).}, label={lst:w2attr}]
{"phi": false,
 "faults": [
  {"clause": "noObl", "handoff": "Build2Code", "rule": "Method2Impl", "binding": "code",
   "agent": "Developer", "fault_class": "unclassified",
   "summary": "noObl at Build2Code/Method2Impl.code (agent Developer): failing tests:
               AssertionError: ValueError not raised"},
  {"clause": "fresh", ..., "reason": "a stamp does not match its current footprint"},
  {"clause": "cover(G)", "goal_element": "Method#withdraw",
   "reason": "Method#withdraw has no accepted deliverable"}],
 "note": "host mode: located only; classification replays need an engine LLM"}
\end{lstlisting}

The location is exact and costs no LLM call: hand-off \code{Build2Code}, rule
\code{Method2Impl}, value \code{code}, agent \code{Developer}, with the
validator's reason. The three faults are one problem seen from three clauses:
an escalated value (\code{noObl}), a value without a valid stamp
(\code{fresh}), and a method without a deliverable (\code{cover(G)}).

What host mode cannot do is \emph{classify} the fault. With an engine LLM,
\mcptool{auto\_attribute} replays the failing value in a fixed order and
reports the first class that fits (\Cref{tab:faults}). Replays need an LLM
the engine can call by itself; in host mode the LLM is \cc{}, so faults stay
\code{unclassified} and \skill{diagnose} tells \cc{} to decide from the
reason.

\begin{table}[!htb]
\centering\small
\caption{Fault classes of \auto{} attribution (\file{auto/attribution.py}), checked in this order.}
\label{tab:faults}
\begin{tabularx}{\textwidth}{@{}l Y Y@{}}
\toprule
Class & What the replay shows & Remedy \\
\midrule
validator & the failing validator gives different verdicts on the same value & fix the validator (it is not deterministic) \\
sampling & a replay on the \emph{same} footprint passes & retry, or a stronger model \\
footprint & a replay on a footprint widened by one hop passes & widen the footprint in the team \\
upstream & a footprint value was written by an LLM upstream, and redoing it lets this value pass & redo the upstream value \\
specification & none of the above & change the prompt or the validator \\
coverage & (a \code{cover(G)} failure with no single value to blame) & change the team's shape \\
\bottomrule
\end{tabularx}
\end{table}

\subsection{Step F: a checked repair}

Here the Developer had everything it needed in its footprint: docstring,
plan and test all say ``raise \code{ValueError}''. This looks like a
\emph{sampling} fault, whose remedy is simply another try. But the value is
escalated, and an escalated value is not offered again on an unchanged
footprint. \skill{diagnose} therefore asks for a \emph{delta}:
\mcptool{auto\_propose} in mode \code{delta} returns the current team and the
fault summaries and asks for ``a minimal change to the team that removes
them''. Our delta added the doctest examples (\code{d.method.examples}) to
the footprint of \code{Method2Impl.code}.

\begin{lstlisting}[style=engine, caption={Observed: \code{auto\_submit\_team} with the delta.}, label={lst:w2delta}]
{"admitted": true, "report": "ADMITTED", "mode": "in-place", "kept_values": 5,
 "agents": ["Designer", "Tester", "Developer"], "handoffs": ["Goal2Design", "Goal2Test", "Build2Code"]}
\end{lstlisting}

The delta passed the same checker before anything changed. How it is
applied depends on what it changes (\Cref{tab:repair}). Ours only edited a
rule of an existing hand-off, so it was applied \emph{in place}.

\begin{table}[!htb]
\centering\small
\caption{How an admitted delta is applied to the running team (\file{auto/repair.py}).}
\label{tab:repair}
\begin{tabularx}{\textwidth}{@{}l Y Y@{}}
\toprule
Mode & When & What happens to accepted values \\
\midrule
\code{noop} & the team is unchanged & nothing \\
\code{in-place} & only rules of existing hand-offs change (prompts, footprints, validators, rules) & kept; values whose prompt or validator changed lose their stamp; values whose footprint changed become stale; \emph{all escalations are cleared} \\
\code{hot} & views, agents or hand-offs are only added & kept; the additions get fresh trace models; all escalations are cleared \\
\code{rebuild} & views, classes or write rights change, or a hand-off is removed or retargeted & discarded: the team starts from scratch \\
\bottomrule
\end{tabularx}
\end{table}

\code{"kept\_values": 5} counts the stamps still stored after the delta: two
plans, two test suites and the \code{deposit} code. Only values written by
the changed rule have a new footprint, so the plans and tests were not
redone; but both implementations, including the accepted \code{deposit}
code, came back as pending, because their footprint now also contains the
examples. This time both implementations were accepted, the next
\mcptool{auto\_run} reported \code{"phi": true}, and
\mcptool{auto\_deliverable} assembled the class (\Cref{lst:w2code}).

\begin{lstlisting}[language=Python, style=engine, caption={Observed: \code{auto\_deliverable} with \code{"complete": true} (assembly by the engine; method bodies typed by us).}, label={lst:w2code}]
class Wallet:
    """A wallet that holds money in cents."""

    def __init__(self):
        self.cents = 0

    def deposit(self, cents):
        self.cents += cents
        return self.cents

    def withdraw(self, cents):
        if cents > self.cents:
            raise ValueError('insufficient funds')
        self.cents -= cents
        return self.cents
\end{lstlisting}

\file{auto/history.json} recorded the whole story: the task; admission round
1 rejected (\W{1}, \W{2}, \W{5}, \W{6}); round 2 admitted; two runs; one
attribution; the delta as round 3, admitted \code{in-place} with 5 values
kept; the final run.

\begin{remark}[Be honest about what the repair did]
Widening the footprint added no information the Developer lacked. And it was
not even what re-opened \code{withdraw}: every admitted in-place or hot
delta clears \emph{all} escalations (\Cref{tab:repair}), so a delta that only
reworded a prompt would have done the same. In this example the correct
answer came from us. With an engine LLM, \mcptool{auto\_attribute} would have
replayed the value on its unchanged footprint and, if a replay passed,
reported a sampling fault, for which the right remedy is a retry or a
stronger model rather than a team change. In host mode that judgement is
left to \cc{} and to you.
\end{remark}

\begin{think}
(a) Why were the two implementations blocked after the first
\mcptool{auto\_run}? (b) After the delta, why were the plans and tests not
redone, while the accepted \code{deposit} code was? (c) Name a delta that
changes no footprint at all but would still have re-opened \code{withdraw}.
\end{think}

\begin{takeaway}
\begin{itemize}
\item In the plugin, \cc{} is the \auto{} builder; the checker returns exact
diagnostics and the revise prompt carries them back.
\item Behaviour validators execute the code against examples and tests, so
a wrong implementation is rejected with the failing assertion.
\item In host mode faults are located exactly but not classified; repairs
are deltas that pass the same checker, keep accepted values where they can,
and clear escalations.
\end{itemize}
\end{takeaway}

% ===========================================================================
\section{The hosted service}\label{sec:hosted}
% ===========================================================================

The plugin runs the engine on the developer's machine. Version 0.3.0 adds a
second way: one server, started with \code{agentm2m serve}, offers the same
tools to many users over the web (requirement R6).

\subsection{One process, three faces}

\begin{figure}[!htb]
\centering
\begin{tikzpicture}[font=\sffamily\small, >={Stealth[length=4pt]},
  c/.style={draw=extbar, fill=white, rounded corners=2pt, align=center, minimum height=8mm, minimum width=24mm, inner sep=3pt},
  cl/.style={draw=llmbar, fill=llmbg, rounded corners=2pt, align=center, minimum height=8mm, minimum width=24mm, inner sep=3pt},
  s/.style={draw=modelbar, fill=modelbg, rounded corners=2pt, align=center, minimum height=8mm, inner sep=3pt},
  e/.style={draw=enginebar, fill=enginebg, rounded corners=2pt, align=center, minimum height=8mm, inner sep=3pt},
  lb/.style={font=\sffamily\scriptsize, align=left, fill=white, inner sep=1.5pt}]
  \node[c] (br) at (0,1.4) {Browser};
  \node[c] (rest) at (0,0) {Script / \code{curl}};
  \node[cl] (cc) at (0,-1.4) {\cc{}};
  \node[s, minimum width=22mm, minimum height=40mm] (front) at (3.9,0) {\textbf{Front}\\(router)};
  \node[s, minimum width=28mm] (web) at (7.7,1.4) {\code{/}\\web app};
  \node[s, minimum width=28mm] (api) at (7.7,0) {\code{/api/...}\\REST};
  \node[s, minimum width=28mm] (mcp) at (7.7,-1.4) {\code{/mcp}\\MCP (HTTP)};
  \node[e, text width=31mm] (ws) at (11.9,-0.7) {per-tenant workspace\\{\scriptsize\code{tenants/<id>/<project>}}};
  \node[e, text width=29mm] (sbx) at (11.9,-2.4) {sandbox (\code{bwrap}),\\checked before use};
  \node[e, text width=22mm] (db) at (3.9,-3.0) {\file{service.db}\\(key digests)};
  \draw[->] (br) -- (front.west |- br); \draw[->] (rest) -- (front); \draw[->] (cc) -- (front.west |- cc);
  \draw[->] (front.east |- web) -- (web); \draw[->] (front) -- (api); \draw[->] (front.east |- mcp) -- (mcp);
  \draw[->] (api.east) -- (ws); \draw[->] (mcp.east) -- (ws);
  \draw[->] (ws) -- node[lb, right] {behaviour validators} (sbx);
  \draw[<->] (front) -- node[lb, right] {resolve key} (db);
\end{tikzpicture}
\caption{The hosted service. One process serves the web site, the REST API
and the MCP endpoint; REST and MCP reach the same per-tenant workspaces, and
all code execution goes through the sandbox. (The web app's Playground calls
the REST endpoint \code{/api/auto/check}.)}
\label{fig:hosted}
\end{figure}

\Cref{fig:hosted} shows the structure (\file{src/agentm2m/server/app.py}).
\code{/} serves the project's web site, including a \emph{Services} tab with
five pages: Status, API keys, a checker Playground (paste a typed team, see
its \W{1}--\W{6} diagnostics), API docs and Pricing. \code{/api/...} is a REST
API whose documentation is generated at \code{/api/docs}. \code{/mcp} is the
MCP endpoint, mounted from the \emph{same} server object as the stdio server,
so it offers the same 24 tools. A small router in front, \code{Front}, sends
each request to the right part. It is an ASGI application: ASGI is Python's
standard interface between a web server and a web application~\cite{asgi},
so one ASGI application can forward to others. Starting the service is one
command, or a container:

\begin{lstlisting}[language=bash, style=raw, caption={Code excerpt: starting the hosted service (from \file{README.md}).}, label={lst:serve}]
pip install -e ".[serve]" && (cd web && npm ci && npm run build)
agentm2m serve --port 8765        # http://localhost:8765, MCP at /mcp, API docs at /api/docs
# or
docker compose up --build
\end{lstlisting}

The health endpoint tells an operator what the service can do. We started
the service in-process with a throw-away data directory (script
\file{rest_demo.py}) on a machine without bubblewrap (\Cref{lst:health}).

\begin{lstlisting}[style=engine, caption={Observed: \code{GET /api/health} (abridged).}, label={lst:health}]
{"status": "ok", "version": "0.3.0",
 "mcp": {"endpoint": "/mcp", "transport": "streamable-http", "auth": "Authorization: Bearer <api key>",
         "tools": ["team_init", "team_status", ...], "tool_count": 24},
 "llm": {"mode": "host", "solve_backend": null},
 "execution": {"sandbox": "process", "isolated": false, "enabled": false,
   "reason": "code execution is disabled: this server has no isolating sandbox (set
              AGENTM2M_SANDBOX=bwrap, or AGENTM2M_SANDBOX=command with AGENTM2M_SANDBOX_CMD)"},
 "keys": {...}, "metrics": {...}, ...}
\end{lstlisting}

\subsection{API keys}

Anyone can create a key with \code{POST /api/keys} (or on the web page),
giving an optional label and an optional, unverified e-mail address; there is
no login. The key is \code{am2m\_} followed by 43 random URL-safe characters.
It is returned \emph{once}; the server stores only its SHA-256 digest, its
first characters (so you can recognise it in a list), the label, a tier and
usage counters, in SQLite. A key is presented as
\code{Authorization: Bearer am2m\_\ldots} (or \code{X-API-Key}) and revoked
with \code{DELETE /api/keys/me}. Every key is on the free tier; the function
that would restrict features by tier, \code{enforce}, exists but
deliberately does nothing yet.

\subsection{Tenants and projects}

Each key is a \emph{tenant}: a separate customer whose data must never mix
with another's. A tenant can have several \emph{projects}, chosen with the
header \code{X-AgentM2M-Project} (default \code{default}). Each project is an
ordinary project directory on the server,
\file{\$AGENTM2M\_DATA\_DIR/tenants/<key id>/<project>/}, holding a
\file{.agentm2m/} workspace exactly like a local one. Project names must match
\code{[A-Za-z0-9\_.-]\{1,64\}} and may not be \code{.} or \code{..}, so a name
cannot climb out of the tenant's directory.

\subsection{A REST session}

\Cref{lst:rest} drives the Wallet task of \Cref{sec:walk2} through the REST
API. Look at the last three answers: they are the service protecting itself.

\begin{lstlisting}[style=engine, caption={Observed: a REST session against the service (abridged).}, label={lst:rest}]
POST /api/keys {"label": "student"}                 -> 201 {"key": "am2m_EKwI...Upu0",
                                                        "record": {"id": "084e1ffa7a89d459", "tier": "free", ...},
                                                        "note": "Store the key now: it is shown only once ..."}
POST /api/auto/task (no key)                        -> 401 {"detail": "missing or invalid API key"}
POST /api/auto/task {"source": "<wallet.py>"}       -> 200 {"task": {"entry": "Wallet", "methods": [...]}}
GET  /api/auto/propose                              -> 200 {"mode": "propose", "prompt": "<8587 chars>"}
POST /api/auto/team {"team": {...}}                 -> 200 {"admitted": true, "mode": "admitted", ...}
POST /api/auto/run                                  -> 200 {"phi": false, "pending": 4, "blocked_on_upstream": 2}
POST /api/auto/bindings {<a test suite>}            -> 403 {"detail": "code execution is disabled: this server
                                                        has no isolating sandbox (set AGENTM2M_SANDBOX=bwrap, ...)"}
POST /api/auto/solve                                -> 501 {"detail": "no engine LLM configured on this server
                                                        (AGENTM2M_SOLVE_LLM); use host mode"}
POST /mcp (no key)                                  -> 401 {"error": "unauthorized", "detail": "send an API key ..."}
\end{lstlisting}

Everything that needs no code execution worked: creating a key, setting the
task, the builder prompt, the checker, the run. The first submission whose
validator would have \emph{executed} code was refused with \code{403},
because this server had no isolating sandbox (\Cref{sec:security}). The
directory \file{tenants/084e1ffa7a89d459/} now holds the project.

Errors map to HTTP status codes: \code{401} missing or invalid key;
\code{400} invalid project name; \code{422} a malformed request body;
\code{409} a workspace error (``no task yet''); \code{403} code execution
refused by the sandbox policy; \code{501} \mcptool{auto\_solve} without a
server LLM; \code{500} anything unexpected. \Cref{tab:rest} lists the
endpoints. The REST API covers the \auto{} tools; the hand-written-team tools
are available over MCP only.

\begin{table}[!htb]
\centering\small
\caption{REST endpoints of the hosted service.}
\label{tab:rest}
\rowcolors{2}{black!3}{white}
\begin{tabularx}{\textwidth}{@{}l l Y@{}}
\toprule
\rowcolor{white}Method and path & Key? & Purpose / equivalent tool \\
\midrule
\code{GET /api/health} & no & status, tool list, sandbox status, counters \\
\code{GET /api/pricing} & no & the tiers \\
\code{POST /api/keys} & no & create a key (shown once) \\
\code{GET}, \code{DELETE /api/keys/me} & yes & show or revoke the presented key \\
\code{POST /api/auto/check} & no & \mcptool{auto\_check} (stateless; used by the playground) \\
\code{GET /api/teams/samples} & no & sample typed teams \\
\code{POST /api/auto/task} & yes & \mcptool{auto\_task\_set} \\
\code{GET /api/auto/propose} & yes & \mcptool{auto\_propose} \\
\code{POST /api/auto/team} & yes & \mcptool{auto\_submit\_team} \\
\code{POST /api/auto/run} & yes & \mcptool{auto\_run} \\
\code{GET}, \code{POST /api/auto/bindings} & yes & \mcptool{auto\_next\_bindings}, \mcptool{auto\_submit\_binding} \\
\code{GET /api/auto/status} & yes & \mcptool{auto\_status} \\
\code{POST /api/auto/attribute} & yes & \mcptool{auto\_attribute} \\
\code{GET /api/auto/deliverable} & yes & \mcptool{auto\_deliverable} \\
\code{POST /api/auto/solve} & yes & \mcptool{auto\_solve} \\
\code{DELETE /api/auto} & yes & \mcptool{auto\_reset} \\
\bottomrule
\end{tabularx}
\end{table}

\subsection{How an MCP call learns who is calling}

The MCP tools are the same functions as in the stdio server. How does a tool
know which tenant called? \Cref{lst:front} shows the answer: the router
authenticates the key and adds a \emph{trusted} header naming the tenant; the
tool reads it through \code{ctx.headers}.

\begin{lstlisting}[language=Python, style=model, caption={Code excerpt: \code{Front.\_\_call\_\_} (\file{server/app.py}, abridged).}, label={lst:front}]
# never trust a client-supplied tenant header
headers = [(k, v) for k, v in scope.get("headers", []) if k.lower() != b"x-agentm2m-tenant"]
...
rec = await asyncio.to_thread(self.keys.resolve, key_from_headers(hmap))
if rec is None:          # -> 401 {"error": "unauthorized", ...}, WWW-Authenticate: Bearer
    ...
headers.append((b"x-agentm2m-tenant", rec.id.encode()))
await self.mcp_app(scope, receive, send)
\end{lstlisting}

The first line is the security-critical one. If the router did not remove a
tenant header sent by the \emph{client}, anybody with any valid key could
name another tenant's id and work in that tenant's projects. Rule of thumb: a
header that carries authority must be set by your own trusted code, and
anything a client sends under that name must be thrown away.

\begin{think}
Describe the attack that would succeed without the first line of
\Cref{lst:front}. What would the attacker need?
\end{think}

\subsection{Two ways to connect \cc{}}\label{sec:forward}

\what{Directly, without the plugin.} Register the endpoint as an MCP server:

\begin{lstlisting}[language=bash, style=raw, caption={Code excerpt: connecting \cc{} directly (from the plugin's \file{README.md}).}, label={lst:mcpadd}]
claude mcp add --transport http agentm2m https://<service-url>/mcp --header "Authorization: Bearer <api key>"
\end{lstlisting}

Nothing is installed locally, and the tools appear as
\code{mcp\_\_agentm2m\_\_*}. You get the tools, but not the skills, subagents
and hooks, whose tool names refer to the plugin's server.

\what{Through the plugin.} Start \cc{} with
\code{AGENTM2M\_URL=https://<service-url>} and \code{AGENTM2M\_API\_KEY=<key>}
(and optionally \code{AGENTM2M\_PROJECT}). The plugin still starts its local
\code{agentm2m-mcp} through \code{uvx}, but that server now \emph{forwards
every tool call} to the service under the same tool name
(\Cref{fig:forward}). The tool-calling skills and subagents work unchanged,
and the team's state lives on the service. Everything that edits or checks
\emph{local files} does not reach the service: the hooks validate local files
and announce only a local workspace; \code{handoff-architect} and
\skill{author-handoff} write rule files on your disk, which the service never
sees. On the service, a hand-written team therefore starts from a template
(\mcptool{team\_init}) and grows through \mcptool{team\_evolve}, whose
\code{rule\_text} creates rule files remotely.

\begin{figure}[!htb]
\centering
\begin{tikzpicture}[font=\sffamily\small, >={Stealth[length=4pt]},
  c/.style={draw=llmbar, fill=llmbg, rounded corners=2pt, align=center, minimum height=10mm, inner sep=4pt},
  s/.style={draw=modelbar, fill=modelbg, rounded corners=2pt, align=center, minimum height=10mm, inner sep=4pt},
  e/.style={draw=enginebar, fill=enginebg, rounded corners=2pt, align=center, minimum height=10mm, inner sep=4pt},
  lbl/.style={font=\sffamily\scriptsize, align=center, fill=white, inner sep=1.5pt}]
  \node[c] (cc) at (0,0) {\cc{}\\+ plugin};
  \node[s] (proxy) at (4.4,0) {local \code{agentm2m-mcp}\\(forwards)};
  \node[s] (front) at (9.2,0) {service \code{/mcp}\\key $\to$ tenant};
  \node[e] (ws) at (13.2,0) {tool runs on\\tenant's project};
  \draw[->] ([yshift=1.5mm]cc.east) -- node[lbl, above=2pt] {stdio: \mcptool{run}} ([yshift=1.5mm]proxy.west);
  \draw[->] ([yshift=1.5mm]proxy.east) -- node[lbl, above=2pt] {HTTPS + key: \code{tools/call run}} ([yshift=1.5mm]front.west);
  \draw[->] ([yshift=1.5mm]front.east) -- ([yshift=1.5mm]ws.west);
  \draw[->, dashed] ([yshift=-1.5mm]ws.west) -- ([yshift=-1.5mm]front.east);
  \draw[->, dashed] ([yshift=-1.5mm]front.west) -- node[lbl, below=2pt] {result} ([yshift=-1.5mm]proxy.east);
  \draw[->, dashed] ([yshift=-1.5mm]proxy.west) -- node[lbl, below=2pt] {same result} ([yshift=-1.5mm]cc.east);
\end{tikzpicture}
\caption{The plugin as a proxy for the hosted service.}
\label{fig:forward}
\end{figure}

This is the \emph{Proxy} design pattern~\cite{gamma1994patterns}: an object
with the same interface as the real one, which forwards each call. One
decorator on every tool implements it (\Cref{lst:forward}).

\begin{lstlisting}[language=Python, style=model, caption={Code excerpt: \code{\_forwardable} (\file{mcp_server.py}, abridged).}, label={lst:forward}]
def _forwardable(fn):
    """Run the tool here, or, when the plugin is pointed at a hosted service
    (AGENTM2M_URL + AGENTM2M_API_KEY), forward the same call there. Calls
    arriving over HTTP (i.e. on the hosted service itself) always run here."""
    def wrapper(*args, **kwargs):
        over_http = ...                          # does the request carry HTTP headers?
        rc = None if over_http else _remote_client()
        if rc is None:
            return fn(*args, **kwargs)           # local mode
        arguments = {k: v for k, v in bound.arguments.items() if k != "ctx"}
        return rc.call_tool(fn.__name__, arguments)   # same name, same arguments
    return wrapper
\end{lstlisting}

The same module runs on the server, which is why calls that arrive over
HTTP always run locally (the guard in the docstring). The forwarding client
(\file{src/agentm2m/remote.py}) is a minimal MCP client over Streamable
HTTP: it initialises once, remembers the session id, re-initialises if the
server answers \code{404} (it has forgotten the session, for example after a
restart), and turns a \code{401} into the actionable message ``create a new
one on its web page (Services $\to$ API keys)''.

\begin{think}
What would happen without the guard ``calls arriving over HTTP always run
here'' if the server's own environment contained \code{AGENTM2M\_URL}
pointing at itself?
\end{think}

\begin{remark}[Three hosts]
``Host'' means three different things in this document. The \emph{MCP host}
is the application that owns the LLM (\cc{}). \emph{Host mode} means that
host writes the values. The \emph{hosted service} is the web server. On the
hosted service, host mode still means that \emph{your} \cc{}, the client,
writes the values: by default the service never calls an LLM. An operator can
set \code{AGENTM2M\_SOLVE\_LLM} (for example \code{anthropic}, with an API
key) to enable \mcptool{auto\_solve}, the unattended loop with at most 3
admission rounds and 2 repairs. \mcptool{auto\_solve} also works locally when
\code{AGENTM2M\_LLM} names an engine LLM.
\end{remark}

\begin{takeaway}
\begin{itemize}
\item \code{agentm2m serve}: web app, REST API and MCP endpoint in one
process, over the same service layer.
\item API keys are shown once and stored as digests; each key is a tenant
with its own project directories.
\item Authority (the tenant) comes from a header only the server sets; a
spoofed one is stripped.
\item The plugin can act as a proxy: the tools and skills run against the
service, but anything that works on local files stays local.
\end{itemize}
\end{takeaway}

% ===========================================================================
\section{Security and trust}\label{sec:security}
% ===========================================================================

The plugin and the service execute code they did not write. This section
collects what can go wrong and what defends against it (requirement R7).

\subsection{Two kinds of foreign code}

\what{Project helpers.} A rule module may say \code{uses 'helpers.py'}; the
engine then imports that Python file to get its helpers and validators.
Importing a file \emph{runs} it. Opening a workspace is therefore like running
\code{make} in a repository: only do it for projects you trust. The engine
refuses paths that leave \file{.agentm2m/} (a \code{uses} clause or rule path
such as \code{../../x.py}), and \mcptool{team\_evolve}'s \code{rule\_text}
can only create \file{.agentm2m} rule files, never Python files: otherwise an
LLM could write a helper file and a rule that imports it. On the hosted
service this matters most: a tenant can create workspaces only from the
built-in templates and extend them with rule files, so a tenant cannot place
its own \file{helpers.py} on the server.

\what{LLM-written values.} Behaviour validators execute code that an LLM
wrote: the \auto{} library validators (\code{examples\_run},
\code{passes\_tests}, \code{test\_valid}, \code{runs}) and template
validators such as \code{failsOnStub} (\devteam{}) and \code{passesDryRun}
(\code{incident}). All of them go through one function, \code{run\_python}
in \file{src/agentm2m/auto/sandbox.py}, which supports four modes
(\Cref{tab:sandbox}).

\begin{table}[!htb]
\centering\small
\caption{Sandbox modes (\code{AGENTM2M\_SANDBOX}).}
\label{tab:sandbox}
\begin{tabularx}{\textwidth}{@{}l Y l@{}}
\toprule
Mode & What it does & Isolating? \\
\midrule
\code{process} (default) & A child process in a fresh temporary directory, with limits on CPU time, memory (2\,GB), file size (64\,MB) and wall-clock time, an isolated interpreter (\code{python -I}) and a minimal environment. It can still read what the user can read and use the network. & no \\
\code{bwrap} & bubblewrap: no network, its own process, IPC and user namespaces, a fresh \file{/proc}, only system directories and the Python installation mounted read-only, plus the run's temporary directory. The data directory and other tenants are invisible. (If the Python environment lives under your home directory, that part of home is visible, read-only.) & yes \\
\code{command} & Prefix every run with \code{AGENTM2M\_SANDBOX\_CMD} (for example an nsjail or firejail call). & assumed (the wrapper's job) \\
\code{off} & Never execute. A submission whose validator would execute code is refused (a tool error over MCP, \code{403} over REST) without spending an attempt. An unknown mode also means \code{off}. & -- \\
\bottomrule
\end{tabularx}
\end{table}

\subsection{Fail closed on a shared server}

On your own machine, \code{process} is acceptable: you run your own team's
code with your own rights. On a shared server it is not: one tenant's code
could read another tenant's files. The service therefore works like this:

\begin{enumerate}[itemsep=1pt]
\item At start-up it declares that isolation is required; from then on the
\code{process} mode is refused.
\item The first time execution status is needed, it \emph{probes} the
configured sandbox by running a tiny program in it, and remembers the result
for the life of the process.
\item If the probe fails or there is no isolating sandbox, code execution is
disabled: the health endpoint and the Status tab say why (\Cref{lst:health}),
and submissions that would execute code are refused (\Cref{lst:rest}).
Checking, task and team submission and status keep working.
\item Only an operator who trusts every key holder can override this with
\code{AGENTM2M\_ALLOW\_UNSANDBOXED\_EXEC=1}.
\end{enumerate}

This is \emph{failing closed}, Saltzer and Schroeder's \emph{fail-safe
defaults}~\cite{saltzer1975protection}: when a safety mechanism is
unavailable, the dangerous feature switches off, not the safety. The
container image installs bubblewrap, sets \code{AGENTM2M\_SANDBOX=bwrap}, and
runs the server as an unprivileged user.

\plain{if the server cannot lock the code in a room, it does not run the
code at all.}

\begin{remark}[A refusal is not a verdict]
When the sandbox refuses to run a validator, the engine does \emph{not} count
it as a rejected attempt: the refusal is passed on as an error, not turned
into a rejection. Otherwise a server without a sandbox would escalate every
correct value after $k$ ``failures'' that were never judged.
\end{remark}

\subsection{A threat table}

\begin{table}[!htb]
\centering\small
\caption{Threats and defences in version 0.3.0 (\cmark{} enforced by code; \pmark{} by documentation or configuration only).}
\label{tab:threats}
\rowcolors{2}{black!3}{white}
\begin{tabularx}{\textwidth}{@{}Y Y c@{}}
\toprule
\rowcolor{white}Threat & Defence & \\
\midrule
A workspace's \file{helpers.py} is malicious. & Only open trusted workspaces. & \pmark \\
Paths in rules or \code{uses} escape \file{.agentm2m/}. & Paths are checked. & \cmark \\
An LLM writes Python through \mcptool{team\_evolve}. & \code{rule\_text} only creates \file{.agentm2m} files. & \cmark \\
A rule expression reaches Python internals through double-underscore names such as \code{\_\_class\_\_}. & Expressions cannot reach private or double-underscore attributes. & \cmark \\
LLM-written code harms the machine. & One sandbox function; resource limits; \code{bwrap} on the server. & \cmark \\
One tenant reads another's projects. & Fail-closed isolation; per-tenant directories; project-name check. & \cmark \\
A client claims another tenant's identity. & The tenant header is stripped and set only by the router. & \cmark \\
The key database leaks. & Only SHA-256 digests are stored; keys are shown once. & \cmark \\
A web page in the user's browser makes requests to a local service under a forged host name (DNS rebinding~\cite{jackson2007dnsrebinding}). & Opt-in host checking with \code{AGENTM2M\_ALLOWED\_HOSTS}; otherwise the endpoint relies on API keys. & \pmark \\
Source data leaks to an LLM provider. & Host mode: nothing leaves \cc{}'s session. Engine LLMs receive only binding prompts (instruction plus footprint). & \cmark \\
\bottomrule
\end{tabularx}
\end{table}

\begin{think}
Why is the \code{process} sandbox good enough when you run the plugin on your
own laptop, but refused on the hosted service? What exactly is the difference
in whose code runs?
\end{think}

\begin{takeaway}
\begin{itemize}
\item Opening a workspace runs its \file{helpers.py}; treat it like any code
in the repository.
\item LLM-written code always runs through one sandbox function; the shared
server refuses to run it without real isolation (fail closed), and a refusal
never counts as a rejection.
\item Authority (the tenant) is set only by trusted code; secrets (keys) are
stored only as digests.
\end{itemize}
\end{takeaway}

% ===========================================================================
\section{Building, testing and releasing the plugin}\label{sec:engineering}
% ===========================================================================

\subsection{One version, six places}

The version number \code{0.3.0} appears in six files: the engine's
\file{pyproject.toml} and \file{src/agentm2m/\_\_init\_\_.py}
(\code{\_\_version\_\_}), the plugin manifest, the development marketplace
entry, the \code{uvx} pin in \file{.mcp.json}, and the pin in the validation
hook (and in a seventh, the public marketplace, once it exists). They must
agree: the skills and subagents name tools that a specific engine version
provides, and the hook must check files with the same engine that runs them.
The script \file{plugin/scripts/bump_version.py} changes all of them at once
(\code{--check} prints ``ok: 0.3.0 in 6 places''), and a test fails if they
differ.

\subsection{Scripts}

\begin{itemize}
\item \file{validate.sh} runs everything continuous integration (CI) checks:
\code{claude plugin validate --strict} on the plugin and the marketplace,
JSON syntax, version consistency, that every template creates a valid
workspace, and both test suites.
\item \file{build.sh} builds the engine's \emph{wheel} (an installable
archive) and \emph{source distribution}, checks them with \code{twine} (the
Python upload tool), verifies that the wheel contains what the plugin needs at
run time (the grammar file, the templates, the \code{agentm2m-mcp} entry
point), installs it in a fresh \code{uvx} environment as a smoke test, and
zips the plugin.
\item \file{dev-install.sh} validates the plugin, checks that the local
engine starts through \code{uvx}, registers the development marketplace and
installs the plugin from the checkout.
\item \file{release.sh} defaults to a \emph{dry run}: it builds everything
and uploads nothing. Publishing (first to TestPyPI, a test copy of the Python
Package Index, then to PyPI, then a git tag and a public marketplace entry)
must be switched on explicitly with \code{AGENTM2M\_ALLOW\_PUBLIC=1}.\footnote{As
of version 0.3.0 the reason is that the paper describing the tool is under
double-anonymous review, where reviewers must not learn who the authors are;
a public package under the authors' names would reveal them.} Listing in
Anthropic's plugin directory is a manual form; the script only prepares the
entry.
\end{itemize}

\file{pyproject.toml} also contains two settings for uv.
\code{exclude-newer = "7 days"} makes uv refuse dependency versions
published in the last week when developers lock or sync \emph{this
repository}: a cooldown against freshly compromised packages, a supply-chain
defence. It does not affect what \code{uvx --from agentm2m==\ldots} resolves on
a user's machine (we checked: uv does not apply a package's own resolver
settings to \code{uvx}); users who want the same cooldown set it in their own
uv configuration. \code{cache-keys} makes \code{uvx} rebuild a local checkout
when its source files change, so \code{AGENTM2M\_ENGINE=<checkout>} always
runs your latest code.

\subsection{Tests}

The plugin has its own test suite (\file{plugin/tests/}) next to the
engine's (\file{tests/}). Its five files test different layers:

\begin{itemize}
\item \file{test_plugin_package.py}: the plugin as installed. Manifests are
valid and versions agree; every skill has frontmatter and names only real
tools; every subagent allowlists only tools the server exposes (the test reads
the tool names from \file{mcp_server.py}); hook scripts are executable and
behave on real input, including the silent exit without \code{uvx}.
\item \file{test_host_mode.py}: the binding protocol. The first run builds
structure without values; a prompt contains only the footprint; a rejection
gives a reason and a retry prompt; escalation after $k$ rejections, cleared by
a footprint change; a value for an outdated footprint is refused; blocked
values unblock.
\item \file{test_change_and_persistence.py}: state survives a reload; impact
is exactly the affected values; an identical re-derivation stops
propagation, a changed one propagates one hop further; a workflow spread over
several processes; instances see each other's writes.
\item \file{test_edit_trace_hot.py}: write rights, engine ownership, atomic
edits, trace queries, and team evolution with retroactive obligations.
\item \file{test_mcp_server.py}: the real server over stdio, as \cc{} uses
it: the full host workflow, an engine backend, the \auto{} workflow, and
forwarding to a hosted service (including a bad key).
\end{itemize}

When we ran both suites for this document, the result was \emph{117 passed, 3
skipped}, in about 23 seconds. Testing an MCP server ``over stdio'' means
starting it as a child process and speaking JSON-RPC to it, exactly as we did
for \Cref{lst:init,lst:list,lst:call,lst:errors}.

\begin{think}
Suppose someone bumps the engine to 0.4.0 in \file{.mcp.json} but forgets
the hook. What can go wrong for a user, and which test would catch it before
release?
\end{think}

\begin{takeaway}
\begin{itemize}
\item The plugin and the engine are released together; one script keeps six
version numbers in step and a test enforces it.
\item Tests target each layer: the installed package, the binding protocol,
persistence, edits, and the real server over stdio.
\item The release script is safe by default (dry run).
\end{itemize}
\end{takeaway}

% ===========================================================================
\section{Requirements revisited, and what is still open}\label{sec:limits}
% ===========================================================================

\Cref{sec:problem} set seven requirements. \Cref{tab:revisit} checks each
against what we have seen.

\begin{table}[!htb]
\centering\small
\caption{The requirements of \Cref{sec:problem}, revisited.}
\label{tab:revisit}
\rowcolors{2}{black!3}{white}
\begin{tabularx}{\textwidth}{@{}l l Y Y@{}}
\toprule
\rowcolor{white}Req. & Status & Evidence & What remains \\
\midrule
R1 & met & 24 tools, 9 slash commands; \Cref{sec:walk1,sec:walk2}; \file{test_mcp_server.py} & -- \\
R2 & met & host mode uses the engine's own acceptance path (\Cref{sec:host}); \file{test_host_mode.py} & no control over which model writes values \\
R3 & partly & workers can call only two tools (\Cref{fig:context}) & the dispatch message and values the main agent fills itself rest on instructions \\
R4 & partly & state survives reloads and processes; \file{test_change_and_persistence.py} & overlapping operations of two processes can lose a reported acceptance (\Cref{sec:arch}); \auto{} files have no lock \\
R5 & partly & the hook re-checks edited team files (\Cref{lst:hookerr}) & edits through \code{Bash} are not seen; the hook is silent without \code{uvx} or \code{python3} \\
R6 & partly & \code{agentm2m serve}; direct \code{claude mcp add} needs nothing installed & the plugin route still runs a local proxy; REST covers \auto{} only; local-file skills and hooks do not reach the service \\
R7 & partly & fail-closed server, one sandbox function, path and header checks (\Cref{tab:threats}) & locally, \code{process} mode does not isolate; \file{helpers.py} is trusted \\
\bottomrule
\end{tabularx}
\end{table}

Beyond the requirements, these limits matter:

\begin{itemize}
\item \textbf{Host-mode attribution only locates.} Classifying a fault
(\Cref{tab:faults}) needs replays by an engine LLM.
\item \textbf{Escalation semantics differ.} For a hand-written team an
escalated value is re-offered only when its footprint changes, so a retry may
need a source edit; for an \auto{} team any admitted repair clears all
escalations.
\item \textbf{Validator strength is not checked.} The hand-written templates
use form validators on code, and \code{failsOnStub} accepts an oracle that
asserts nothing; $\varphi$ can hold for wrong code.
\item \textbf{\auto{} is narrow.} Python classes and functions only; one tool
(\code{exec}); acyclic hand-offs; one writer per view.
\item \textbf{The service is a research preview.} No rate limits, tiers not
enforced, in-memory metrics, no verified e-mail addresses, no login.
\item \textbf{Not yet published.} Until the engine is on PyPI,
\code{AGENTM2M\_ENGINE} must point at a checkout.
\item \textbf{Scale is unmeasured.} State is one JSON file per workspace and
every run re-evaluates all matches; we have not measured large teams in the
plugin.
\end{itemize}

\subsection{Conclusion}

A research engine became a tool by keeping one idea at every layer:
\emph{an LLM writes, a program decides}. MCP gives \cc{} a precise, typed way
to call the engine; host mode lets \cc{} do all the LLM work while the engine
keeps structure, acceptance and completion; subagent allowlists and
\code{footprint\_version} protect the one property host mode puts at risk,
that every value is written from its footprint; hooks catch broken designs
the moment they are written; and the hosted service offers the same tools to
many users without letting their data or code meet. Where the design falls
short, the shortfalls are known, named, and mostly measurable, which is
itself what a checked composition is meant to give.

% ===========================================================================
\section{Exercises}\label{sec:exercises}
% ===========================================================================

Difficulty: $\star$ check your reading, $\star\star$ apply, $\star\star\star$
design. Hands-on exercises are marked (H); \Cref{app:try} shows how to run
the engine. Hints are in \Cref{app:hints}.

\begin{enumerate}[leftmargin=1.6em, itemsep=4pt]
\item $\star$ \textbf{(Protocol)} Write the JSON-RPC request that calls
\mcptool{impact} for the change of \Cref{sec:walk1}, Step 6. Then write the
response for the case where \code{as\_agent} names an agent that does not
exist. Is it a protocol error or a tool error?
\item $\star$ \textbf{(States)} For each of these situations, give the state
of the value in \Cref{fig:states} and the status \mcptool{submit\_binding}
would return: (a) a value submitted with an outdated
\code{footprint\_version}; (b) the third wrong answer on one footprint; (c) a
correct answer for a value that is already fresh.
\item $\star\star$ \textbf{(Schema)} \mcptool{model\_edit} declares \code{ops}
as \code{list[dict[str, Any]]}; its docstring describes three kinds of
operation. Sketch the JSON Schema the SDK generates. What would you gain and
lose with a precise schema for each kind of operation?
\item $\star\star$ \textbf{(Plugin anatomy)} Write a new skill
\skill{explain} that, given an element key, explains in plain words what was
derived from it and why. Give its frontmatter (\code{name},
\code{description}, \code{argument-hint}) and the tools it may call.
\item $\star\star$ \textbf{(Least privilege)} Exercise the colleague's
proposal of Q\ref{q:read} further: design a way to give the worker the
project's style guide without breaking change impact.
\item $\star\star$ \textbf{(Concurrency)} Two \cc{} sessions run \skill{run}
on the same project at the same time. List every mechanism of
\Cref{sec:arch,sec:host} that keeps the state consistent, and describe one
interleaving that ends with a \code{stale} status and one that loses a
reported acceptance.
\item $\star\star$ (H) \textbf{(Templates)} Initialise the \code{incident}
template with \code{AGENTM2M\_LLM=mock} and run it. Explain every
escalation from the template's comments and validators.
\item $\star\star$ (H) \textbf{(Validators)} Add a validator to
\file{helpers.py} that rejects an oracle without an \code{assert}
statement, with a \code{Rejected} reason. Use it in \code{Req2Test} and
observe the \code{retry\_prompt} for a bad oracle.
\item $\star\star\star$ \textbf{(Design)} Modify \code{Operation2CodeEdit} of
\devteam{} so that a code body is accepted only if it passes the oracles of
its story's criteria. What must the footprint or the validator's inputs
contain, and which \auto{} condition does the modified team now satisfy?
\item $\star\star$ (H) \textbf{(Service)} Start \code{agentm2m serve}
locally, create a key with \code{curl}, and drive the Wallet task through the
REST API up to the first \code{403}. Then set the sandbox so that execution
is enabled, explain why your choice is or is not safe for a shared server,
and finish the task.
\item $\star\star$ \textbf{(Security)} Show a project name that the rule
\code{[A-Za-z0-9\_.-]\{1,64\}} without \code{.} and \code{..} blocks and that
could otherwise escape the tenant's directory. Is \code{..x} dangerous?
\item $\star\star\star$ \textbf{(Repair)} In \Cref{sec:walk2} the delta
widened a footprint. Propose a delta that would help a Developer who really
lacked information, and say whether it is applied in place, hot, or by
rebuilding.
\item $\star\star\star$ \textbf{(Design critique)} The REST API covers only
the \auto{} tools. Write a one-page design note: should the hand-written-team
tools get REST endpoints too? Consider testability, security
(\file{helpers.py}) and the users of each interface.
\end{enumerate}

\appendix
% ===========================================================================
\section{Reference: the 24 tools}\label{app:tools}
% ===========================================================================

Arguments in square brackets are optional. ``Changes state'' means the tool
may write to the workspace.

\begin{table}[!htb]
\centering\small
\caption{\tool{} tools (hand-written teams, \file{.agentm2m/}).}
\label{tab:tools-a}
\rowcolors{2}{black!3}{white}
\begin{tabularx}{\textwidth}{@{}L{3.1cm} L{3.4cm} Y c@{}}
\toprule
\rowcolor{white}Tool & Arguments & What it does & Changes state \\
\midrule
\mcptool{team\_init} & [template], [force] & Copies a template (\code{devteam}, \code{research}, \code{incident}) into \file{.agentm2m/}; \code{force} replaces an existing workspace. & yes \\
\mcptool{team\_status} & -- & Agents and views, hand-offs, trace links, binding states, open items, $\varphi$. & no \\
\mcptool{team\_validate} & -- & Parses the spec and every rule module; checks classes, root slots, helper paths; counts matches. & no \\
\mcptool{model\_show} & view, [key], [depth] & Shows a view or one element; marks engine-owned elements. & no \\
\mcptool{model\_edit} & view, ops, as\_agent & Atomic \code{create}/\code{set}/\code{delete} as the owning agent; rebuilds structure at once. & yes \\
\mcptool{impact} & [view, ops, as\_agent] & Previews the obligations of an edit on a scratch copy. No LLM call. & no \\
\mcptool{run} & [max\_passes] & Runs every hand-off to a fixpoint; reports created, pending, blocked, escalated, $\varphi$. & yes \\
\mcptool{next\_bindings} & [agent], [limit] & The next values to fill, each with its prompt and \code{footprint\_version}. Host mode only; runs the structural phase first. & yes \\
\mcptool{submit\_binding} & target\_key, binding, value, footprint\_version & Judges one value with its validator (\Cref{tab:status-submit}). & yes \\
\mcptool{trace\_query} & key, [transitive] & Upstream sources of an element and everything derived from it. & no \\
\mcptool{team\_evolve} & agent, view, handoff, rule, view\_spec or view\_spec\_file, [rule\_text] & Adds an agent, its view and a hand-off while the team runs. The view is given inline or as a YAML file in \file{.agentm2m/}; \code{rule\_text} creates the rule file. & yes \\
\mcptool{acceptance} & -- & $\varphi$ and the list of what is still open. & no \\
\bottomrule
\end{tabularx}
\end{table}

\begin{table}[!htb]
\centering\small
\caption{\auto{} tools (builder-proposed teams for a Python task, \file{.agentm2m/auto/}).}
\label{tab:tools-b}
\rowcolors{2}{black!3}{white}
\begin{tabularx}{\textwidth}{@{}L{3.1cm} L{3.4cm} Y c@{}}
\toprule
\rowcolor{white}Tool & Arguments & What it does & Changes state \\
\midrule
\mcptool{auto\_task\_set} & source, [description] & Lifts a Python skeleton into the \cls{Goal} view; discards the run state. & yes \\
\mcptool{auto\_check} & [team\_json], [naive] & \W{1}--\W{6} on any typed team, without changing anything. \code{naive} uses a weaker, class-level \W{4} (for comparison only). & no \\
\mcptool{auto\_propose} & [mode] & The builder prompt: \code{propose}, \code{revise}, \code{delta}, or \code{auto} (picks one). & no$^{*}$ \\
\mcptool{auto\_submit\_team} & team\_json & Checks a typed team; admits it, or applies it as a checked delta (\Cref{tab:repair}). & yes \\
\mcptool{auto\_run} & -- & Runs to a fixpoint and evaluates $\varphi$. & yes \\
\mcptool{auto\_next\_bindings} & [agent], [limit] & The next values to fill (role, task, format, footprint). & yes \\
\mcptool{auto\_submit\_binding} & target\_key, binding, value, footprint\_version & The library validators decide; behaviour validators execute the code. & yes \\
\mcptool{auto\_status} & -- & Task, team, last diagnostics, $\varphi$ and failing clauses, pending values, history. & no \\
\mcptool{auto\_attribute} & [limit] & Locates every failed clause; with an engine LLM also classifies it (\Cref{tab:faults}). & yes \\
\mcptool{auto\_deliverable} & -- & The assembled Python code and which methods are accepted. & no \\
\mcptool{auto\_solve} & [source] & The whole loop unattended (at most 3 admission rounds, 2 repairs); needs an engine LLM. & yes \\
\mcptool{auto\_reset} & -- & Deletes \file{.agentm2m/auto/}. & yes \\
\bottomrule
\end{tabularx}
\par\smallskip{\footnotesize $^{*}$ In mode \code{delta} without stored fault reports, it first runs attribution.}
\end{table}

% ===========================================================================
\section{Try it yourself}\label{app:try}
% ===========================================================================

\paragraph{With \cc{}} You need \cc{}, uv (which provides \code{uvx}),
Python 3.11 or newer, and a checkout of the repository.

\begin{lstlisting}[language=bash, style=raw, caption={Code excerpt: installing the plugin from a checkout (adapted from \file{plugin/scripts/dev-install.sh}).}, label={lst:install}]
git clone <repository> agentm2m-repo && cd agentm2m-repo
python -m venv .venv && . .venv/bin/activate && pip install -e ".[dev]"
plugin/scripts/dev-install.sh            # validates and installs agentm2m@agentm2m-dev
export AGENTM2M_ENGINE=$PWD              # the engine is not on PyPI yet
mkdir /tmp/play && cd /tmp/play && git init
claude                                   # a scratch project, NOT the engine repository
\end{lstlisting}

In \cc{}, \code{/plugin} should list \code{agentm2m}, and \code{/mcp} should
show the \code{agentm2m} server as connected. Then type \skill{init devteam}
and \skill{run}, and follow \Cref{sec:walk1}. If the server does not start,
check that \code{uvx --version} works and that \code{AGENTM2M\_ENGINE} is set
in the shell that started \cc{}.

\paragraph{Without \cc{}} Everything except the LLM is available from the
command line and from Python. The scripts in
\file{docs/plugin-walkthroughs/} reproduce every observed output of this
document; \code{AGENTM2M\_LLM=mock} lets the engine fill values with canned
text.

\begin{lstlisting}[language=bash, style=raw, caption={Code excerpt: the engine without \cc{}.}, label={lst:nocc}]
python docs/plugin-walkthroughs/walkthrough1_devteam.py      # Section 10
python docs/plugin-walkthroughs/walkthrough2_autom2m.py      # Section 11
python docs/plugin-walkthroughs/mcp_stdio_demo.py /tmp/p     # Section 5 (empty dir /tmp/p)
python docs/plugin-walkthroughs/rest_demo.py                 # Section 12
agentm2m workspace --dir /tmp/q --llm mock init incident && agentm2m workspace --dir /tmp/q --llm mock run
agentm2m auto check teams/devteam_proposal.json
agentm2m auto --dir /tmp/w task docs/plugin-walkthroughs/wallet.py && agentm2m auto --dir /tmp/w propose --prompt-only
\end{lstlisting}

\begin{table}[!htb]
\centering\small
\caption{Environment variables.}
\label{tab:env}
\rowcolors{2}{black!3}{white}
\begin{tabularx}{\textwidth}{@{}L{5.3cm} Y@{}}
\toprule
\rowcolor{white}Variable & Meaning \\
\midrule
\code{AGENTM2M\_LLM} & Who fills values: \code{host} (default), \code{anthropic}, \code{openai}, \code{ollama}, \code{mock}. \\
\code{AGENTM2M\_MODEL} & Model for engine LLMs. \\
\code{AGENTM2M\_MAX\_RESAMPLES} & $k$ for hand-written teams in every mode (default 3); \auto{} uses a fixed $k=3$. \\
\code{AGENTM2M\_ENGINE} & Engine for \code{uvx} (default \code{agentm2m==0.3.0}); a path for a checkout. \\
\code{AGENTM2M\_PROJECT\_DIR} & Project the local server opens (default: \code{CLAUDE\_PROJECT\_DIR}, then the current directory). \\
\code{AGENTM2M\_URL}, \code{AGENTM2M\_API\_KEY}, \code{AGENTM2M\_PROJECT} & Forward every tool call to a hosted service. \\
\code{AGENTM2M\_SANDBOX}, \code{AGENTM2M\_SANDBOX\_CMD} & \code{process}, \code{bwrap}, \code{command} (with its wrapper), \code{off}. \\
\code{AGENTM2M\_DATA\_DIR} & Server: keys database and tenant projects (default \file{\textasciitilde/.agentm2m-service}). \\
\code{AGENTM2M\_SOLVE\_LLM} & Server: engine LLM for \mcptool{auto\_solve}. \\
\code{AGENTM2M\_ALLOW\_UNSANDBOXED\_EXEC} & Server: \code{1} allows execution without isolation (trusted users only). \\
\code{AGENTM2M\_ALLOWED\_HOSTS} & Server: opt-in DNS-rebinding protection for \code{/mcp}. \\
\bottomrule
\end{tabularx}
\end{table}

% ===========================================================================
\section{The admitted team of Walkthrough 2}\label{app:team}
% ===========================================================================

\begin{lstlisting}[style=model, caption={Code excerpt: the admitted typed team (the builder prompt's worked example, renamed; the fixed \cls{Goal} view elided).}, label={lst:fullteam}]
{"name": "wallet_team", "goal_view": "Goal",
 "agents": [
   {"name": "Designer",  "tools": [],
    "role": "You decide, for one method, the algorithm, edge cases and the fields it reads or updates."},
   {"name": "Tester",    "tools": ["exec"],
    "role": "You turn a method's documentation into executable unit tests with concrete expected values."},
   {"name": "Developer", "tools": ["exec"],
    "role": "You implement one method following its design so that its tests and examples pass."}],
 "views": {"Goal": "<fixed: Task, Method>",
   "Design": {"classes": {"MethodDesign": {"attributes": {"name": "string", "plan": "string"},
                          "references": {"method": {"type": "Goal.Method", "required": true}}}}},
   "Test":   {"classes": {"MethodTest":   {"attributes": {"name": "string", "code": "string"},
                          "references": {"method": {"type": "Goal.Method", "required": true}}}}},
   "Code":   {"classes": {"MethodImpl":   {"attributes": {"name": "string", "code": "string"},
                          "references": {"method": {"type": "Goal.Method", "required": true}}}}}},
 "writes": {"Designer": ["Design"], "Tester": ["Test"], "Developer": ["Code"]},
 "handoffs": [
   {"name": "Goal2Design", "sources": ["Goal"], "target": "Design", "rules": [{"name": "Method2Design",
     "from": [{"var": "m", "type": "Goal!Method"}], "guard": null, "to": {"var": "d", "type": "Design!MethodDesign"},
     "bind": {"name": "m.name", "method": "m"},
     "llm": [{"feature": "plan", "prompt": "Describe in a few bullet points how to implement this method: ...",
              "footprint": ["m.task.skeleton", "m.signature", "m.docstring"], "validator": [{"id": "nonempty"}]}]}]},
   {"name": "Goal2Test", "sources": ["Goal"], "target": "Test", "rules": [{"name": "Method2Test",
     "from": [{"var": "m", "type": "Goal!Method"}], "guard": null, "to": {"var": "t", "type": "Test!MethodTest"},
     "bind": {"name": "m.name", "method": "m"},
     "llm": [{"feature": "code", "prompt": "Write unit tests for this method.",
              "footprint": ["m.task.skeleton", "m.name", "m.docstring", "m.examples"],
              "validator": [{"id": "test_valid", "args": {"method": "m"}}]}]}]},
   {"name": "Build2Code", ... see the rule in Section 11 ...}],
 "goal": [{"class": "Method", "scope": "all", "kind": "checked"},
          {"class": "Method", "scope": "all", "kind": "delivered"}],
 "deliverable": {"view": "Code", "class": "MethodImpl", "feature": "code", "for": "method"},
 "done": ["cover(G)", "valid", "fresh", "noObl"]}
\end{lstlisting}

Read it top to bottom: three agents (only those that own behaviour-checked
values need \code{exec}, \W{6}); one view per agent and one writer per view
(\W{2}); three hand-offs from the \cls{Goal} view; two goal obligations, one
saying every method must reach a behaviour check (\code{checked}) and one
saying every method must be delivered (\code{delivered}); the deliverable,
which tells the engine where the code of each method is; and the four engine
clauses of $\varphi$ (\W{5}).

% ===========================================================================
\section{Glossary}\label{app:glossary}
% ===========================================================================

\paragraph{Protocols and the web}
\begin{description}[leftmargin=1.2em, style=nextline, itemsep=1pt]
\item[API key] A secret that identifies a caller of a web service; stored by the server only as a digest.
\item[ASGI] Python's standard interface between web servers and web applications.
\item[Digest / hash] A short fingerprint of data; equal data gives an equal digest, any change a different one.
\item[JSON-RPC] Remote procedure calls in JSON: requests (with \code{id}), responses, notifications.
\item[MCP] Model Context Protocol: how AI hosts discover and call the tools of servers.
\item[MCP host, client, server] The application that owns the LLM; its connection to one server; the program offering tools.
\item[Primitive] One of MCP's kinds of capability: tools, resources, prompts.
\item[REST] An API style: resources as paths, operations as HTTP methods.
\item[Server-sent events] A way for one HTTP response to deliver several messages one after another.
\item[Session id] The \code{Mcp-Session-Id} a Streamable HTTP server may assign; the client repeats it on every request.
\item[stdio / Streamable HTTP] MCP's transports: pipes to a child process; HTTP POSTs to one endpoint.
\item[Tool error vs.\ protocol error] A result with \code{isError: true} that the LLM can act on, vs.\ a JSON-RPC \code{error} (a bug).
\end{description}

\paragraph{\cc{} and the plugin}
\begin{description}[leftmargin=1.2em, style=nextline, itemsep=1pt]
\item[Allowlist] The tools a subagent may call; all others are refused.
\item[Context] Everything the LLM can see in one call.
\item[Extension point] A place where a plugin may add behaviour: MCP tools, skills, subagents, hooks.
\item[Frontmatter] The YAML header of a skill or subagent file.
\item[Hook] A command \cc{} runs at an event (session start, after a tool use); exit code 2 reports a problem.
\item[Inversion of control] The host calls the plugin, not the other way round.
\item[Marketplace] A list of plugins from which \cc{} installs them.
\item[Skill] A plain-language procedure for \cc{}, invoked as a slash command or loaded automatically.
\item[Subagent] A helper agent with its own instructions, fresh context and tool allowlist.
\item[\code{uvx}] Runs a Python package's program in a cached, isolated environment.
\end{description}

\paragraph{The engine (\tool{})}
\begin{description}[leftmargin=1.2em, style=nextline, itemsep=1pt]
\item[Binding (structural / stochastic)] A clause in a rule that fills a field: by a formula, or by an LLM from a footprint.
\item[Blocked] A pending value whose footprint contains an upstream value that was never accepted.
\item[Deferred backend] The host backend: the engine records pending values instead of sampling.
\item[Engine-owned] An element created by a rule; agents may not edit it.
\item[Escalation] Giving up on a value after $k$ rejections on the same footprint.
\item[Fixpoint] The state in which another pass of all hand-offs changes nothing.
\item[Footprint] The exact data an LLM may see when writing one value.
\item[\code{footprint\_version}] The digest of a footprint, sent with each prompt and required with each submission.
\item[Guard] The condition in parentheses after a rule's source pattern.
\item[Host mode] \cc{}, not the engine, writes the values (\code{AGENTM2M\_LLM=host}).
\item[HOT] Higher-order transformation: a change of the team (new agent, view, hand-off) while it runs.
\item[Lift binding] \code{self <- @llm(...)}: one JSON answer fills several fields.
\item[Metamodel / model / view] A form (class diagram); a filled-in form; the model an agent owns.
\item[Obligation] A value that must be redone because its footprint changed.
\item[Pending] A value ready to be written: never accepted (missing) or stale.
\item[$\varphi$] The acceptance predicate: the engine's definition of ``done''.
\item[\code{PendingBinding}] The record of one pending value: place, prompt, attempts, \code{footprint\_version}.
\item[Rule] ``For every source object matching this pattern and guard, create one target object and fill its fields.''
\item[Stamp] The digest of the footprint a value was accepted for.
\item[Stale] A value whose stamp no longer matches its current footprint.
\item[Target key] The name of a value's place: \code{Rule::variable::match}.
\item[Trace link] A record that a target object was created from a source object by a rule.
\item[Validator (form / behaviour)] A check a value must pass: ``does it parse?'' vs.\ ``does it pass the tests?''.
\item[Write rights ($\omega$)] Which agent may write which view.
\end{description}

\paragraph{\auto{}}
\begin{description}[leftmargin=1.2em, style=nextline, itemsep=1pt]
\item[Builder] The LLM that writes a typed team; in the plugin, \cc{} (or the \code{team-builder} subagent).
\item[Checker, \W{1}--\W{6}] The program that admits a typed team only if the six conditions of \Cref{tab:w} hold.
\item[\code{cover(G)}, \code{valid}, \code{fresh}, \code{noObl}] The four clauses of $\varphi$ (\Cref{sec:recap-auto}).
\item[Delta] A revised typed team that repairs a running one; applied in place, hot, or by rebuilding.
\item[Deliverable] Where the result of each goal (the code of each method) is stored.
\item[Diagnostic] One checker message: condition, element, what is wrong.
\item[Doctest example] A \code{>>>} line in a docstring with its expected output.
\item[Fault class] Sampling, footprint, upstream, specification, validator or coverage (\Cref{tab:faults}).
\item[Goal view / lift] The fixed view (\cls{Task}, \cls{Method}) into which the task's source is turned.
\item[Stub] A method with a docstring and a body of \code{pass}, \code{...} or \code{raise NotImplementedError}.
\item[Typed team] A team design in a fixed JSON format that the checker can read.
\end{description}

\paragraph{Security and engineering}
\begin{description}[leftmargin=1.2em, style=nextline, itemsep=1pt]
\item[Atomic write] Write a temporary file, then rename it over the old one; readers never see half a file.
\item[DNS rebinding] An attack in which a web page reaches a local service through a forged host name.
\item[Dry run] Do everything except the irreversible step (here: uploading).
\item[Fail closed] When a safety mechanism is unavailable, switch off the risky feature, not the safety.
\item[Least privilege] Give each component only the rights it needs.
\item[Lost update] Two overlapping read-modify-write operations; one write silently undoes the other.
\item[Namespace (Linux)] A private view of the network, process list or file system for one process.
\item[Optimistic concurrency] Work without locking; check at the end that nothing changed underneath.
\item[Proxy] An object with the same interface as another, which forwards calls to it.
\item[Sandbox] An environment that restricts what executed code can do.
\item[Tenant] One customer (here: one API key) of a shared service, whose data is kept apart.
\end{description}

% ===========================================================================
\section{Answers to the questions}\label{app:answers}
% ===========================================================================

\what{Q1 (JSON-RPC ids).} Responses may arrive in a different order than the
requests were sent, for example when a quick call overtakes a slow one.
Without the \code{id} the client could attach the result of one call to the
other, for instance treat the result of \mcptool{impact} as that of
\mcptool{run}.

\what{Q2 (2 versus 3 matches).} \code{Story2Operation} has the guard
\code{s.status = \#accepted}, which excludes the draft \code{S3}.
\code{Criterion2TestCase} has no guard, so it matches all three criteria,
including the draft's \code{S3.1}.

\what{Q3 (descriptions).} The LLM would expect a value and might report ``the
value is accepted'' without reading the status, or treat a status object as
the value. No program would notice: the schema describes only the
arguments, and the description is never checked against the code. That is
why descriptions list every outcome and live next to the code as docstrings.

\what{Q4 (extension points).} (a) A \code{PostToolUse} hook. (b) A skill.
(c) A subagent whose allowlist lacks \code{Read}, \code{Grep} and
\code{Bash}. (d) An MCP tool (\mcptool{impact}), which a skill can call.

\what{Q5 (lost update).} Before writing, A's save notices that the
fingerprint changed and reloads the state from disk, which contains B's
change but not A's in-memory acceptance; it then writes that. So B's change
survives and A's is lost. It is acceptable because the state stays
consistent: A's value is simply \code{missing} again and is offered by the
next \mcptool{next\_bindings}. What is lost is work, not correctness.

\what{Q6 (why not call \cc{} from \code{generate}).} The engine runs
\emph{inside} a tool call that \cc{} made: \cc{} is waiting for the tool to
return. The server cannot ask \cc{}'s LLM a question in the middle of that
call (and MCP tools are not given a way to do so here); and the value should
be written by a subagent with a fresh, restricted context, which only \cc{}
can start. So the tool returns the prompt, and \cc{} answers in a later call.

\what{Q7 (outdated submission).} \code{stale}. The engine recomputes the
digest of the current footprint, which now contains the new signature and
criterion; it differs from the submitted \code{footprint\_version}.
Accepting the code would stamp it as derived from information it never saw,
and a later change would not know that it must be redone.

\what{Q8 (\code{Read} for the worker).} A stamp records the digest of the
footprint and promises ``this value depends on this data and nothing else''.
A worker that also reads files could let a style rule, or any other file,
influence the value. When that file changes, no stamp changes, so the value
is not redone: change impact is no longer exact.

\what{Q9 (compile-only validator).} \code{compiles()} on \code{CodeEdit.body}
checks only form. \W{4} would reject the design: the \cls{Code} view reaches
no behaviour check.

\what{Q10 (trace links versus impact).} \code{Story2Operation} matches a
\cls{UserStory}, not a \cls{Criterion}, so its trace link names
\code{UserStory\#S2}; the criterion enters only through the
\emph{footprint} \code{s.criteria} (\Cref{fig:fptrace}). Obligations are
found by comparing stamps with current footprints, not by following trace
links, so the signature and the code body are found anyway.
\mcptool{trace\_query} on \code{UserStory\#S2} shows \code{op\_s2} and its code
edit.

\what{Q11 (\auto{} walkthrough).} (a) Their footprints contain \code{d.plan}
and \code{t.code}, which had never been accepted. (b) The delta changed only
the rule \code{Method2Impl}; the footprints of the plans and tests did not
change, so their stamps stayed fresh. The \code{deposit} code is written by
the changed rule, whose footprint now also contains the examples, so its
stamp no longer matched. (c) Any admitted in-place delta, for example one
that only rewords the Developer's prompt: every in-place or hot delta clears
all escalations.

\what{Q12 (spoofed tenant header).} A user with \emph{any} valid key, for
example a free one of their own, could send
\code{x-agentm2m-tenant: <victim's key id>} and read, change or delete the
victim's projects. The attacker needs a valid key and the victim's key id.
Key ids are random, but they are identifiers, not secrets: they appear in
key records, logs and directory names. Security must not depend on keeping
them hidden, and the router makes sure it does not.

\what{Q13 (forwarding loop).} The server would forward every incoming call
to itself, which would forward it again, and so on, until a timeout or
resource limit ended the chain. No call would ever execute.

\what{Q14 (process sandbox).} Locally, the code belongs to your own team on
your own machine; it can do nothing you could not do yourself, and the
resource limits stop accidents. On a shared server the code comes from many
users; a same-user process could read the data directory with every tenant's
projects and the keys database. Only a sandbox that hides those files is
acceptable.

\what{Q15 (version skew).} The MCP server would run 0.4.0 while the hook
validates with 0.3.0: a rule or typed team that 0.4.0 accepts might be
reported as broken after every edit, or the reverse. The version-consistency
test in \file{test_plugin_package.py} (and \code{bump_version.py --check} in
\file{validate.sh}) fails before release.

% ===========================================================================
\section{Hints for the exercises}\label{app:hints}
% ===========================================================================

\begin{enumerate}[leftmargin=1.6em, itemsep=2pt]
\item Copy the shape of \Cref{lst:errors}; the unknown agent gives a tool error.
\item (a) The value stays where it was; status \code{stale}. (b) escalated. (c) \code{already\_accepted}.
\item A list of objects with a string field \code{op} and free other fields. A precise schema (one per kind) lets the client reject bad operations early, at the cost of a larger, harder-to-read schema.
\item Allowed tools: \mcptool{trace\_query}, \mcptool{model\_show}, \mcptool{impact}; read-only, so it can be user-invocable.
\item Put the style guide into the \emph{footprint}, for example as an attribute of a view that the rule reads; then a change of the guide is a footprint change.
\item Use the lost-update table of \Cref{sec:arch} and the \code{footprint\_version} story of \Cref{sec:host}.
\item The mock's canned text is not a runnable script, so \code{passesDryRun} rejects every \code{dryRun} value (the rule's comment says so).
\item Parse the oracle with Python's \code{ast} module and look for an \code{Assert} node; return \code{Rejected("...")} otherwise.
\item The validator must see the oracles: add a reference from the operation to its criteria's test cases, or a helper that collects them, and run them against the body. \W{4}.
\item \code{AGENTM2M\_SANDBOX=bwrap} if bubblewrap works on your machine; otherwise \code{AGENTM2M\_ALLOW\_UNSANDBOXED\_EXEC=1}, safe only if you are the only key holder.
\item \code{..} and \code{.}: \code{tenants/<id>/..} is the tenants directory. \code{..x} is an ordinary name inside the tenant's directory.
\item For example, add a Reviewer view with examples of edge cases that \code{Method2Impl} reads: a new view and hand-off are added, so the delta is \code{hot}.
\item Compare with \Cref{sec:arch} (one service layer) and \Cref{sec:security} (helpers run code).
\end{enumerate}

\bibliographystyle{ieeetr}
\bibliography{autom2m-references,autom2m-plugin-references}

\end{document}
